

\medterm{Acquired Methemoglobinemia} An acquired hematologic disorder in which hemoglobin iron is oxidized from the ferrous ($\text{Fe}^{2+}$) to the ferric ($\text{Fe}^{3+}$) state by exogenous oxidizing agents (such as dapsone, benzocaine, lidocaine, nitrites, or sulfonamides), producing clinical cyanosis, chocolate-brown arterial blood, and cellular hypoxia refractory to supplemental oxygen.

\synonyms
Toxic Methemoglobinemia; Drug-Induced Methemoglobinemia; Chemically Induced Methemoglobinemia

\medterm{Benign Melanocytic Lesion} An older or imprecise term for a benign melanocytic growth, such as a melanocytic nevus or mole; unlike malignant melanoma, it is not cancerous.

\medterm{Blood Biomarker} A measurable sign of a disease or condition that can be isolated or detected in a blood sample, such as a disease-associated antibody, hormone, protein, or other substance.

\medterm{Bradykinesia and Athetosis} Hypokinetic and dyskinetic movement patterns characterized by generalized slowness of movement initiation (bradykinesia) or continuous slow, writhing, twisting involuntary movements (athetosis), frequently observed in Parkinson disease, secondary parkinsonism, or basal ganglia encephalopathies.

\synonyms
Slowness of Movement; Hypokinesia and Athetosis; Extrapyramidal Slowness

\medterm{Cessation of Menstruation} Cessation of menstrual periods, called amenorrhea; it is a normal part of menopause but may also result from pregnancy, illness, endocrine disorders, medications, stress, excessive exercise, or inadequate nutrition.

\medterm{Chorea and Myoclonus} Sudden, rapid, irregular, jerky, and non-rhythmic involuntary muscle contractions or shock-like jerks. Chorea flows unpredictably from one muscle group to another, whereas myoclonus represents brief shock-like twitches, observed in metabolic encephalopathies, post-hypoxic syndromes, and neurodegenerative disorders.

\synonyms
Jerky Involuntary Movements; Choreiform Activity; Myoclonic Jerks

\medterm{Congenital Heart Murmur} Alternate name; see \textit{Heart Murmur}.

\medterm{Continuous Passive Motion Machine} A device that gently moves a limb through a set range of motion while the person relaxes. It may be used after selected operations, especially knee surgery, but does not replace active rehabilitation.

\medterm{Cryptococcal Meningitis} A subacute, life-threatening fungal meningitis caused by \textit{Cryptococcus neoformans} or \textit{Cryptococcus gattii}, occurring predominantly in immunocompromised individuals (such as advanced HIV infection or organ transplant recipients). Diagnosis involves lumbar puncture with opening pressure measurement, CSF cryptococcal antigen (CrAg) lateral flow assay, and India ink stain.

\synonyms
Cryptococcal Meningoencephalitis; Fungal Meningitis CrAg; Cryptococcosis of CNS

\medterm{Female Urethral Meatus} The external opening of the female urethra, through which urine leaves the body; it is located anterior and superior to the vaginal opening.

\medterm{German Measles} Alternate name; see \textit{Rubella}.

\medterm{Golfer's Elbow} An overuse tendinopathy affecting the common flexor-pronator tendon origin at the medial epicondyle of the humerus. Characterized by localized tenderness and pain exacerbated by resisted wrist flexion and forearm pronation.

\synonyms
Medial Epicondylitis; Medial Epicondylalgia; Pitcher's Elbow

\medterm{Gynecomastia, Male} Alternate name; see \textit{Gynecomastia}.

\medterm{Heart Murmur Entry} Alternate name; see \textit{Heart Murmur}.

\medterm{Involuntary Movements} A diverse group of neurologic motor disorders characterized by uncontrollable, purposeless, or abnormal spontaneous body movements, including tremors, chorea, athetosis, hemiballismus, dystonia, myoclonus, and tics, stemming from basal ganglia, cerebellar, or motor circuit pathology.

\synonyms
Hyperkinetic Movement Disorders; Involuntary Motor Activity; Dyskinetic Movements

\medterm{Lambert-Eaton Myasthenic Syndrome} An autoimmune disorder in which antibodies interfere with nerve signals reaching muscles. It causes weakness, especially in the hips and shoulders, and may cause dry mouth, reduced reflexes, or problems with automatic body functions.

\synonyms
LEMS; Lambert-Eaton Syndrome; Eaton-Lambert Syndrome

\medterm{MacConkey Agar} A laboratory growth medium used mainly to grow and distinguish certain intestinal bacteria. It helps show whether bacteria can break down lactose, a sugar.

\medterm{Macerate} To soften and break down tissue after prolonged contact with fluid. The term may describe skin, wounds, or changes in a fetus that has died before delivery.

\medterm{Maceration} Softening and breakdown of tissue caused by prolonged exposure to moisture or fluid; it may occur in skin, wounds, or a fetus after death.

\medterm{Macewen Sign} An unusually hollow sound heard when part of the skull is tapped. It may occur with enlarged fluid spaces in the brain or rarely a brain abscess, but it cannot diagnose either condition by itself.

\medterm{Machupo Virus} An arenavirus that causes Bolivian hemorrhagic fever, a potentially fatal illness characterized by fever, pain, and bleeding; it is endemic in parts of Bolivia.

\medterm{Macroamylasemia} Macroamylasemia is a benign condition in which amylase binds to immunoglobulin and remains in blood, causing high serum amylase with low urinary clearance.

\medterm{Macroangiopathy} A disease of the body's larger blood vessels, often caused by fatty buildup or a blood clot. It can reduce blood flow to the heart, brain, legs, or other organs.

\medterm{Macroautophagy} A cell-cleaning process in which a cell encloses damaged parts or unwanted proteins, breaks them down, and reuses some of their materials. It helps maintain cell health during stress and normal turnover.

\medterm{Macrocephalus} Having an abnormally large head, usually defined by head circumference above the expected range for age and sex.

\medterm{Macrocephaly} An unusually large head for a person's age and sex. It may run in a family or be linked to fluid buildup, increased brain growth, or another condition.

\medterm{Macrocyte} An abnormally large red blood cell, commonly seen in macrocytosis and megaloblastic or other anemias.

\medterm{Macrocytic} Describing a cell, especially a red blood cell, that is abnormally large; macrocytic anemia features enlarged red blood cells.

\medterm{Macrocytic Anaemia} Alternate British spelling; see \textit{Macrocytic Anemia}.

\medterm{Macrocytic Anemia} A form of anemia characterized by abnormally large circulating red blood cells, defined on a complete blood count by a mean corpuscular volume (MCV) greater than $100\,\text{fL}$. It is divided into two major categories: (1) megaloblastic anemia, caused by impaired DNA synthesis from vitamin B12 or folate deficiency (often featuring hypersegmented neutrophils, glossitis, and peripheral neuropathy in B12 deficiency), and (2) non-megaloblastic macrocytosis, commonly caused by chronic alcohol consumption, liver disease, hypothyroidism, myelodysplastic syndrome, or reticulocytosis. Symptoms include fatigue, pallor, shortness of breath, and tingling in the extremities.
\synonyms{Macrocytosis; Large-Cell Anemia; High-MCV Anemia}

\medterm{Macrodactyly} Fingers or toes that are unusually large because their bones and soft tissues overgrow. It is usually present from birth, may affect one digit or several, and can interfere with movement or footwear.

\medterm{Macrodontia} Teeth that are larger than expected. It may affect one tooth or many teeth and can cause crowding, bite problems, or difficulty cleaning; treatment depends on the cause and symptoms.

\medterm{Macrogenitosomia} Unusual enlargement of the external genitals caused by excessive hormone activity before birth or during childhood. Possible causes include disorders of the adrenal glands, testes, ovaries, or other hormone systems.

\medterm{Macroglobulin} A large protein in the blood, often made of several protein units joined together. Immunoglobulin M is a major example; abnormal amounts can occur in some immune or blood disorders.

\medterm{Macroglossia} Macroglossia is an abnormally enlarged tongue that can affect feeding, speech, airway, or dental development and may be congenital or acquired.

\medterm{Macrognathia} Abnormal enlargement of the jaw, which may be associated with pituitary gigantism, tumors, or other disorders; also called a prognathic mandible.

\medterm{Macrolide} A class of antibiotics that stops bacteria from making proteins. Examples include azithromycin, clarithromycin, and erythromycin; side effects and medicine interactions vary.

\medterm{Macrolides} Antibiotics that stop bacteria from making proteins and treat selected infections. Common examples include azithromycin, clarithromycin, and erythromycin.

\medterm{Macronutrients} Nutrients needed in large amounts, mainly carbohydrates, proteins, and fats. They provide energy, build and repair tissues, support hormones and cell membranes, and help maintain normal body functions.

\medterm{Macrophage} A large immune cell that surrounds and removes germs, dead cells, and debris. It also helps coordinate inflammation and tissue repair.

\medterm{Macrophage Activation} The process by which macrophages, a type of immune cell, respond to infection or tissue injury. They may destroy germs and damaged cells, alert other immune cells, and release substances that promote inflammation or repair.

\medterm{Macrophage Activation Syndrome} A dangerous complication of some rheumatic diseases in which immune cells become overactive and cause severe inflammation, fever, low blood counts, and organ problems.

\medterm{Macrophage Count} The number or percentage of macrophages found in a tissue, body fluid, or laboratory sample. The result may help assess infection, inflammation, tissue repair, or a blood disorder, but its meaning depends on the sample and other findings.

\medterm{Macropsia} A visual disturbance in which objects appear larger than they really are. It may occur with migraine, seizures, retinal problems, medicines, or brain disorders and can temporarily affect depth and distance judgment.

\medterm{Macroscopic} Large enough to be seen without a microscope. A macroscopic examination records a specimen’s visible size, shape, color, surface, and other gross features before microscopic testing.

\medterm{Macroscopic Hematuria} Visible blood in the urine, also called gross hematuria; it may arise from the kidneys, ureters, bladder, prostate, or urethra and requires evaluation for causes such as infection, stones, trauma, or malignancy.

\medterm{Macrosomia} Macrosomia means unusually large fetal or newborn size, often associated with diabetes, prolonged pregnancy, parental size, or other factors.

\medterm{Macrostomia} An abnormally wide mouth or oral opening, usually a congenital facial cleft caused by incomplete fusion of the embryonic facial processes.

\medterm{Macula} The central part of the retina responsible for sharp, detailed, and color vision. Damage from aging, diabetes, inflammation, or other disease can cause blurred, distorted, or missing central vision.

\medterm{Macula Lutea} The small yellowish area in the center of the retina responsible for sharp vision. Its center, called the fovea, provides the clearest vision.

\medterm{Macular} Relating to a macule, a flat, circumscribed change in skin color without a change in texture or elevation.
Macules may be caused by pigment, inflammation, vascular change, or other processes and are
distinguished from raised lesions such as papules.

\medterm{Macular Cyst} A cystic defect in the macula, the central retinal area responsible for sharp vision, often associated with vitreous traction and potentially causing distortion or loss of central vision.

\medterm{Macular Degeneration} Damage to the macula, the central part of the retina responsible for sharp, straight-ahead vision. The term most often refers to age-related macular degeneration. In the dry form, the macula becomes thinner and may develop deposits called drusen; advanced dry disease can involve loss of light-sensitive retinal cells and usually progresses gradually. In the wet form, abnormal blood vessels grow beneath the retina and leak blood or fluid, often causing faster loss or distortion of central vision. Straight lines may look wavy, and blank or blurry areas may develop. Peripheral vision is usually preserved, but reading, recognizing faces, and other detailed tasks may become difficult.

\medterm{Macular Degeneration, Dry} Alternate name; see \textit{Dry Macular Degeneration}.

\medterm{Macular Degeneration, Wet} Alternate name; see \textit{Wet Macular Degeneration}.

\medterm{Macular Drusen} Small yellow deposits beneath the retina in the macula, the area responsible for sharp central vision. A few small deposits may occur with aging, while many large deposits can signal increased risk of macular degeneration.

\medterm{Macular Edema} Swelling of the macula caused by fluid leaking from retinal blood vessels. It can blur or distort central vision and may result from diabetes, vein blockage, inflammation, surgery, or other eye disease.

\medterm{Macular Hole} A small opening in the center of the retina that can cause a missing or distorted area in central vision. Some cases need surgery.

\medterm{Macular Pucker} A thin layer of scar-like tissue on the retina that can wrinkle the macula and cause distorted or blurred central vision.

\medterm{Macular Skin Lesion} Relating to a macule, a small, flat, circumscribed change in skin color that is neither raised nor depressed.

\medterm{Macule} A flat, clearly defined change in skin color that cannot be felt as a raised or depressed area. Macules may be lighter, darker, red, or brown and can have many causes.

\medterm{Macules} Small, flat, circumscribed changes in skin color that are neither raised nor depressed; they are distinguished from raised papules and fluid-filled vesicles.

\medterm{Maculopapular Lesion} A skin lesion combining flat areas of altered color (macules) with small raised solid lesions (papules); it can occur in infections, drug eruptions, inflammatory diseases, and other conditions.

\medterm{Mad Cow Disease} Alternate name; see \textit{Bovine Spongiform Encephalopathy}.

\medterm{Madarosis} Loss or absence of eyelashes, eyebrows, or other hair, caused by local disease, inflammation, trauma, or systemic conditions.

\medterm{Madura Foot} A chronic localized infection of the foot or other tissue, also called mycetoma, producing swelling, draining sinuses, and grains.

\medterm{Madura Foot} A chronic localized infection of the skin and deeper tissues of the foot, caused by fungi or certain bacteria and characterized by swelling, draining sinuses, and granules.

\synonyms
Mycetoma Pedis; Maduromycosis

\medterm{Maduramycosis} Mycetoma, a chronic subcutaneous infection caused by true fungi or filamentous bacteria and characterized by swelling, sinuses, and granules.

\medterm{Madurella} Madurella is a genus of fungi that can cause mycetoma, a chronic infection of skin and deeper tissue after implantation through the skin.

\medterm{Madurella Mycetomatis} A fungus that commonly causes eumycetoma, a chronic infection producing nodules, draining sinuses, and grains, usually after traumatic implantation in the foot or limb.

\medterm{Maffucci Syndrome} A rare condition causing multiple benign cartilage growths and abnormal blood-vessel growths, usually beginning in childhood. It can deform bones, cause fractures, and increase the risk of certain cancers.

\medterm{MAG3 Renal Scan} A kidney scan that uses a small radioactive tracer to show how well each kidney works and how urine drains. It can help detect blockage or delayed drainage.

\medterm{Magenblase Syndrome} Stomach discomfort and bloating caused by swallowing too much air, often while eating or drinking. The swallowed air may cause burping, fullness, or visible swelling and can be worsened by chewing gum, fizzy drinks, eating quickly, or anxiety.

\medterm{Magnesium} A mineral needed for healthy nerves, muscles, bones, heart rhythm, and many body processes. Too little can cause weakness, tremor, muscle cramps, or abnormal heart rhythms.

\medterm{Magnesium Blood Test} A blood test measuring serum magnesium, an electrolyte involved in neuromuscular function, cardiac rhythm, and metabolism; it helps evaluate deficiency, excess, and related electrolyte disorders.

\medterm{Magnesium Carbonate} A magnesium salt used in some antacid or laxative preparations and as a magnesium supplement; excessive intake may cause toxicity, especially in renal failure.

\medterm{Magnesium Chloride} A magnesium salt used in some oral supplements and medical preparations to provide magnesium or, in concentrated solutions, to help correct magnesium deficiency; dosing depends on the elemental magnesium content.

\medterm{Magnesium Deficiency} An abnormally low body magnesium level, often caused by poor intake, gastrointestinal
loss, renal wasting, or medications. It may cause weakness, tremor, muscle cramps,
seizures, cardiac arrhythmias, hypokalemia, and hypocalcemia.

\medterm{Magnesium Hydroxide} An osmotic laxative and antacid that neutralizes gastric acid and draws water into the bowel.

\medterm{Magnesium In Diet} Dietary magnesium comes from foods such as nuts, seeds, legumes, whole grains, and leafy vegetables and supports nerve, muscle, bone, and heart function.

\medterm{Magnesium Oxide} An inorganic magnesium compound used in some supplements and antacids; it can also have a laxative effect, and excessive use may cause hypermagnesemia, especially with kidney impairment.

\medterm{Magnesium Sulfate} A magnesium salt used intravenously or intramuscularly to treat selected magnesium deficiencies and conditions such as eclampsia, and used orally as an osmotic laxative; it is also known as Epsom salt in its hydrated form.

\medterm{Magnesium Sulfate For Neuroprotection} Magnesium given through a vein before some very early preterm births to lower the baby's risk of brain injury. The timing and dose are guided by obstetric and newborn-care specialists.

\medterm{Magnesium Sulphate} A magnesium salt used medically for severe magnesium deficiency, prevention or treatment of eclamptic seizures, and selected arrhythmias or asthma indications.

\medterm{Magnetic Resonance Angiography} MRI techniques used to depict arteries or veins, with or without contrast. It may evaluate stenosis, aneurysm, dissection, thrombosis, and vascular malformation without ionizing radiation.

\medterm{Magnetic Resonance Cholangiopancreatography} A special MRI scan that shows the bile ducts, gallbladder, and pancreatic duct without inserting a camera into them. It can detect stones, narrowing, inflammation, leaks, and some blockages without using X-rays.

\synonyms
MRCP; MR Cholangiopancreatography

\medterm{Magnetic Resonance Imaging} A scan that uses a strong magnetic field and radio waves to make detailed pictures inside the body without X-rays. It is especially useful for the brain, spinal cord, joints, muscles, and many organs; some scans require injected contrast.

\medterm{Magnetic Resonance Safety} Precautions used during an MRI scan to prevent injury from the strong magnet, radio waves, heating, loud noise, and metal objects. Staff must check implants, metal fragments, medicines, pregnancy, and other risks before scanning.

\medterm{Magnetic Resonance Spectroscopy} A special MRI test that measures chemicals in tissue rather than only making a picture of its shape. It may help study brain disorders or provide extra information about a tumor, but results must be interpreted with ordinary imaging and other tests.

\medterm{Magnetic Resonance Venography} An MRI test that produces images of veins, with or without injected contrast. It can assess suspected clots or narrowing in veins of the brain, abdomen, pelvis, or limbs and does not use X-rays.

\synonyms
MRV; MR Venography

\medterm{Magnetic Seed} A tiny magnetic marker placed in a breast lesion or lymph node to help a surgeon find it. A special detector locates the marker during surgery, providing an alternative to some wire or radioactive-marker methods.

\medterm{Magnetic Therapy} Treatment that uses a magnetic field for a health purpose. Ordinary static magnets have not been shown to treat disease, while specialized magnetic stimulation devices have established uses for some neurologic or psychiatric conditions.

\medterm{Magnetocardiography} A test that records the very weak magnetic fields produced by the heart’s electrical activity using sensors outside the body. It may help study heart rhythm or electrical conduction, usually in specialist or research settings.

\medterm{Magnetoencephalography} A noninvasive technique that records the weak magnetic fields produced by neuronal electrical activity, using sensors outside the head to localize and study brain function and some epileptic discharges.

\medterm{Magnusiomyces Capitatus} An opportunistic yeast formerly known as \textit{Geotrichum capitatum} or \textit{Saprochaete capitata}; it can cause invasive infections, especially in people with severe immunosuppression or hematologic malignancy.

\medterm{Mahvash Disease} A rare inherited disorder in which the body’s cells do not respond normally to glucagon, a hormone that helps raise blood sugar. It can enlarge the pancreas, increase glucagon levels, and disturb blood-sugar control.

\medterm{Main-En-Griffe} A claw-like hand position usually caused by damage to the ulnar nerve, which controls several hand muscles. The ring and little fingers may bend at the smaller joints, weaken the grip, and become difficult to straighten.

\medterm{Maintenance} Ongoing care or treatment intended to preserve health, prevent relapse or deterioration, or sustain a therapeutic result; maintenance therapy is continued after initial control of a condition.

\medterm{Maintenance Dose} The regular amount of a medicine needed to keep its effect at the desired level. The dose depends on how quickly the body removes the medicine and how much the body absorbs.

\medterm{Maintenance Of Wakefulness Test} A sleep test that measures how long a person can remain awake in a quiet setting during the day.

\medterm{Maintenance Therapy} Treatment continued after a disease has improved or been controlled to maintain that benefit, prevent return or worsening, or delay progression. The treatment may be a medicine, procedure, or other care, and its duration depends on the condition, response, risks, and treatment goals.

\medterm{Maize} Corn, a cereal grain used as food and feed; contaminated maize can expose people to fungal toxins such as aflatoxin.

\medterm{Majeed Syndrome} A rare autoinflammatory disorder caused by LPIN2 variants, characterized by chronic recurrent multifocal osteomyelitis, congenital dyserythropoietic anemia, and inflammatory skin disease.

\medterm{Majewski Syndrome} A rare inherited skeletal disorder involving abnormal bone growth and distinctive craniofacial or limb findings; the exact features depend on the genetic subtype.

\medterm{Major} Larger, greater, or more significant than another related structure or condition; for example, the teres major muscle is larger than the teres minor muscle.

\medterm{Major Depression} Alternate name; see \textit{Major Depressive Disorder}.

\medterm{Major Depression With Psychotic Features} A major depressive episode accompanied by delusions or hallucinations, usually mood-congruent, requiring prompt psychiatric assessment and treatment.

\medterm{Major Depressive Disorder} A mood disorder involving one or more major depressive episodes. During an episode, five or more symptoms occur during the same two-week period, including either depressed mood or loss of interest or pleasure (anhedonia). Other symptoms may include changes in sleep, appetite, energy, concentration, movement, guilt, worthlessness, or thoughts of death. The symptoms cause clinically significant distress or impairment.

\synonyms
Major Depression; Clinical Depression

\medterm{Major Histocompatibility Complex} A group of genes and cell-surface proteins that display pieces of germs or abnormal cells to the immune system. In humans, these proteins are called HLA.

\medterm{Major Histocompatibility Complex Molecule} A cell-surface protein that displays small pieces of germs or abnormal cells to immune cells. This helps the immune system recognize what belongs to the body and what may need an immune response.

\synonyms
MHC Molecule

\medterm{Major Trauma} A serious injury that may threaten life, limb, or organ function. It can involve multiple body regions or major damage to the brain, spine, chest, abdomen, pelvis, or blood vessels.

\medterm{Malabsorption} A condition in which the body cannot properly digest or absorb nutrients from food. It may cause chronic diarrhea, greasy stools, weight loss, anemia, or vitamin and mineral deficiencies.

\medterm{Malabsorption Syndrome} Impaired absorption of nutrients from the small intestine, causing diarrhea, weight loss, anemia, vitamin deficiencies, or other nutritional complications.

\medterm{Malacia} Abnormal softening or weakening of tissue. It may affect cartilage, bone, or brain tissue after faulty development, injury, infection, reduced blood supply, or another disease; symptoms depend on the tissue involved.

\medterm{Malacoplakia} A rare long-term inflammation in which immune cells collect mineral deposits. It most often affects the bladder or digestive tract and can form a mass that looks like cancer.

\medterm{Maladaptation} An adjustment to a situation or environment that does not work well and harms health or daily functioning. The term may describe physical, psychological, or behavioral responses to illness or stress.

\medterm{Maladaptive Stress Response} A lasting or overly strong reaction to stress that interferes with health, daily life, or relationships. It may include constant worry, poor sleep, irritability, avoidance, or physical symptoms and can improve with support and appropriate treatment.

\medterm{Malaise} Malaise is a general feeling of illness, discomfort, or lack of well-being that commonly accompanies infection, inflammation, chronic disease, medication effects, or other systemic conditions.

\medterm{Malakoplakia} A rare long-term inflammation in which immune cells collect mineral deposits. It most often affects the bladder or digestive tract and can form a mass that resembles cancer.

\medterm{Malar} Relating to the cheekbone or cheek area. A malar rash is a red or purple rash across the cheeks and bridge of the nose, often seen in lupus. A malar fracture is a break in the cheekbone and may cause swelling, numbness, or a changed facial shape.

\medterm{Malar Rash} A red or violaceous rash across the cheeks and bridge of the nose, often sparing the folds beside the nose.
It is classically associated with systemic lupus erythematosus but can occur with other conditions;
appearance alone does not establish a diagnosis.

\medterm{Malaria} A life-threatening blood infection caused by single-celled \textit{Plasmodium} parasites and spread by the bite of infected female \textit{Anopheles} mosquitoes. Once in the body, parasites multiply in the liver and then invade red blood cells, causing recurring attacks that follow three stages: shaking chills (cold stage), high fever (hot stage), and heavy sweating (sweating stage) every 1 to 3 days. \textit{Plasmodium falciparum} is the most dangerous species because it clogs small blood vessels, leading to cerebral malaria (seizures, coma), severe anemia, dark urine (blackwater fever), lung fluid, or kidney failure. \textit{Plasmodium vivax} can stay dormant in the liver and cause relapses months later unless treated with primaquine. Diagnosis is made by examining a blood smear under a microscope or with a rapid test. Treatment uses artemisinin-based combination medications (ACTs) or intravenous artesunate for severe disease.

\synonyms
Plasmodium Infection; Paludism; Ague

\medterm{Malassezia} Malassezia is a yeast normally found on skin that can contribute to dandruff, seborrheic dermatitis, pityriasis versicolor, and some folliculitis.

\medterm{Malassezia Furfur} Malassezia furfur is a yeast that normally lives on skin but can multiply excessively. It is linked with dandruff, oily skin inflammation, and pityriasis versicolor, which causes lighter or darker patches.

\medterm{Malassimilation} Failure to properly digest food or absorb its nutrients. It may cause diarrhea, weight loss, bloating, anemia, or vitamin shortages and can result from pancreatic disease, bowel disorders, surgery, infection, or food intolerance.

\medterm{Malate Dehydrogenase} An enzyme in energy-producing pathways that changes malate into oxaloacetate. It helps cells process nutrients and is involved in transferring chemical energy within mitochondria and other cellular compartments.

\medterm{Malathion} An organophosphate insecticide that inhibits acetylcholinesterase; topical formulations have also been used to treat head lice, but exposure or ingestion can cause cholinergic toxicity.

\medterm{Malathion Poisoning} Poisoning caused by the organophosphate insecticide malathion. It can produce sweating, salivation, vomiting, diarrhea, wheezing, slow heartbeat, muscle twitching, weakness, seizures, or breathing failure and needs urgent assessment.

\medterm{MALDI-TOF Mass Spectrometry} A laboratory method that identifies bacteria or fungi by comparing their protein pattern with a reference database. It can identify an organism quickly, but separate testing is needed to determine antibiotic sensitivity.

\medterm{Male} Describing a sex associated with reproductive anatomy that produces or is organized around small gametes (sperm); chromosomes, hormones, anatomy, and gender identity do not always fit a simple binary pattern.

\medterm{Male Breast Cancer} Cancer developing in breast tissue in a man. It is uncommon, but a new lump, nipple discharge, skin change, or nipple retraction needs medical evaluation.

\synonyms
Male Mammary Carcinoma; Breast Carcinoma in Men

\medterm{Male Breast Pathology} Disorders affecting male breast tissue, including gynecomastia, benign epithelial lesions, inflammatory conditions, and carcinoma. Evaluation of a persistent mass or nipple change requires clinical examination and appropriate imaging, with biopsy when indicated.

\medterm{Male Condom} A thin sheath placed over the penis before sexual contact to reduce the risk of pregnancy and transmission of many sexually transmitted infections by acting as a physical barrier.

\medterm{Male Genitalia} The external and internal reproductive organs of a male, including the penis, scrotum, testes, epididymides, vas deferens, seminal vesicles, prostate, and associated ducts and glands.

\medterm{Male Hypoactive Sexual Desire Disorder} Persistently reduced sexual thoughts or desire in a man that causes personal distress or relationship difficulty. It differs from erection problems, although both can occur together; medicines, illness, mood, and relationship factors may contribute.

\medterm{Male Hypogonadism} Reduced production of testosterone, impaired sperm production, or both due to disease of the testes or
the hypothalamic-pituitary system. Possible features include low libido, erectile difficulty,
infertility, reduced muscle mass, fatigue, and loss of bone density.

\synonyms
Low Testosterone; Testosterone Deficiency

\medterm{Male Infertility} Reduced ability to achieve a pregnancy because of factors affecting sperm production, sperm
transport, ejaculation, sexual function, or other reproductive processes.

\synonyms
Male Subfertility; Male Factor Infertility

\medterm{Male Menopause} Alternate informal term; see \textit{Male Hypogonadism}.

\medterm{Male Pattern Baldness} Androgenetic alopecia, a patterned progressive loss of scalp hair caused by genetic susceptibility and androgen effects on hair follicles.

\medterm{Male Pelvis} The lower part of the abdomen between the hip bones in a male; compared with the female pelvis, it is generally narrower, taller, and more robust, with a narrower pubic arch and sacrum.

\medterm{Male Phenotype} The observable reproductive anatomy, secondary sex characteristics, and other traits associated with male sex development, produced by interacting genetic, hormonal, and developmental factors and showing natural variation.

\medterm{Male Reproductive System} The organs and tissues that produce sperm and male sex hormones and deliver sperm, including the testes, ducts, accessory glands, penis, and scrotum.

\medterm{Malformation} An abnormal structural development present from birth, caused by disturbed formation of an organ or body part.

\medterm{Malignancy} Cancerous disease in which abnormal cells grow uncontrollably and may invade nearby tissue or spread to distant organs. Its likely behavior and treatment depend on the cancer type, extent, grade, and important molecular findings.

\medterm{Malignancy-Associated Hypercoagulability} A cancer-related tendency toward excessive clot formation caused by tumor-associated coagulation activation, inflammation, platelet interaction, treatment, immobility, or vascular injury. See \textit{Cancer-Associated Thrombosis}.

\medterm{Malignant} Cancerous, with the ability to invade nearby tissues or spread to distant parts of the body. Malignant cells grow abnormally and may damage organs, but behavior varies among cancer types and stages.

\medterm{Malignant Acanthosis Nigricans} A paraneoplastic syndrome in which velvety, hyperpigmented, verrucous plaques develop
rapidly in flexural and mucosal areas as a marker of an underlying internal malignancy,
most commonly gastric adenocarcinoma. Unlike benign acanthosis nigricans, the malignant
form is extensive, abrupt in onset, and may involve the oral mucosa.

\medterm{Malignant Ameloblastoma} A rare cancer arising from ameloblastoma-forming tissue, usually in the jaw, with potential to invade locally or spread to distant organs.

\medterm{Malignant Ascites} Buildup of fluid in the abdomen caused by cancer involving the abdominal lining or nearby organs. It can cause swelling, discomfort, early fullness, nausea, or breathlessness and may require drainage and cancer-directed treatment.

\medterm{Malignant Bone Neoplasm} A cancer that begins in bone or nearby supporting tissue. It may cause persistent bone pain, swelling, or a fracture; scans and a biopsy identify the type and extent before treatment.

\medterm{Malignant Fibrous Histiocytoma} Alternate name; see \textit{Undifferentiated Pleomorphic Sarcoma}.

\medterm{Malignant Giant Cell Tumor Of Tendon Sheath} A rare locally aggressive malignant tumor resembling a giant cell tumor of tendon sheath, characterized by destructive growth and potential for metastasis; diagnosis requires expert soft-tissue pathology review.

\medterm{Malignant Glioma} An aggressive primary tumor arising from glial or glial-like cells of the central nervous system; high-grade examples include anaplastic astrocytoma and glioblastoma.

\medterm{Malignant Hypertension} A medical emergency in which very high blood pressure is causing new injury to organs such as the brain, heart, kidneys, or eyes. It can cause headache, confusion, chest pain, breathlessness, vision changes, seizures, or kidney failure.

\medterm{Malignant Hyperthermia} A rare, life-threatening reaction to certain anesthetic medicines that causes rapidly rising body temperature, severe muscle stiffness, fast heart rate, and dangerous chemical changes. It requires immediate emergency treatment and affects some genetically susceptible people.

\medterm{Malignant Mastocytosis} An older name for a severe disorder in which abnormal mast cells multiply and collect in organs. It can damage the bone marrow, liver, spleen, or other tissues and may cause flushing, faintness, bleeding, or organ failure.

\medterm{Malignant Melanoma} A cancer that begins in melanocytes, the cells that make skin pigment. It may appear as a changing dark or unusual mole and can spread to other tissues, especially when diagnosed at a later stage.

\medterm{Malignant Melanoma} Alternate name; see \textit{Melanoma}.

\medterm{Malignant Mesothelioma} A cancer of the thin lining around the lungs, abdomen, or heart. It is strongly linked to asbestos exposure and may cause chest or abdominal pain, fluid buildup, cough, or breathlessness.

\medterm{Malignant Myoepithelioma} A rare cancer arising from supportive contractile cells in a gland, most often a salivary gland. It can grow into nearby tissue and spread elsewhere, causing a slowly enlarging or painful mass.

\medterm{Malignant Neoplasm} A cancerous growth in which abnormal cells multiply without control, invade nearby tissues, and may spread to distant organs through the blood or lymphatic system.

\synonyms
Cancer; Malignancy

\medterm{Malignant Otitis Externa} An invasive infection of the external ear canal and skull base, usually caused by Pseudomonas aeruginosa in older adults with diabetes or immunocompromise; it can cause severe pain and cranial-nerve complications.

\medterm{Malignant Peripheral Nerve Sheath Tumor} An uncommon soft-tissue cancer arising from or near peripheral nerves. It may cause a growing painful mass, numbness, or weakness and is treated with specialist surgery, radiation, or medicines when appropriate.

\medterm{Malignant Pleural Effusion} Buildup of fluid around the lung because cancer involves the lung lining or has spread there. It may cause breathlessness, chest discomfort, or cough; drainage relieves symptoms, while cancer treatment may reduce recurrence.

\medterm{Malignant Pustule} An older name for the skin form of anthrax, caused by Bacillus anthracis and beginning as an itchy papule that develops into a painless black eschar.

\medterm{Malignant Syphilis} A severe form of secondary syphilis that causes fever and widespread crusted, ulcerated, or tissue-damaging skin lesions. It is more common with weakened immunity and requires prompt specialist treatment for syphilis.

\medterm{Malignant Teratoma Of The Mediastinum} A malignant germ-cell tumor containing tissues from more than one embryonic layer and arising in the mediastinum; evaluation uses imaging, tumor markers, and tissue diagnosis.

\medterm{Malignant Thymoma} An older term for an invasive thymic epithelial tumor; modern usage distinguishes invasive thymoma from thymic carcinoma, which is a separate, more aggressive thymic malignancy.

\medterm{Malignant Tumor} A cancerous growth capable of invading nearby tissue or spreading to distant sites. Its behavior and treatment depend on its type, grade, stage, and molecular features.

\medterm{Malingering} Deliberately pretending to have, or exaggerating, symptoms to obtain an outside benefit such as money, medicines, shelter, or avoiding an obligation. It is not itself a mental illness and should not be diagnosed without reliable evidence.

\medterm{Mallampati Classification} A grading system used before anesthesia to estimate how difficult it may be to place a breathing tube. The clinician looks inside the open mouth and records which parts of the throat can be seen; it is only one part of the airway assessment.

\medterm{Malleolus} A bony prominence at the ankle formed by the lower end of the tibia medially or fibula laterally.

\medterm{Mallet} A surgical hammer used with tools such as osteotomes to cut or shape bone. It delivers controlled force during operations, helping surgeons work on bone while limiting unnecessary injury to nearby tissues.

\medterm{Mallet Finger} An injury to the extensor tendon at the distal finger joint, causing inability to actively straighten the fingertip. It may result from a forced flexion injury and sometimes includes an avulsion fracture.

\medterm{Mallet Toe} A flexion deformity of the distal interphalangeal joint, causing the end of a lesser toe to bend downward. It may be flexible or rigid and can cause pain, rubbing, corns, or difficulty wearing shoes; it differs anatomically from hammertoe and claw toe.

\medterm{Mallet Toe and Hammertoe} Alternate name; see \textit{Hammertoe}.

\medterm{Malleus} The largest of the three small bones in the middle ear, also called the hammer. It is attached to the eardrum and passes sound vibrations to the incus and then the inner ear.
\commonname{Hammer}

\medterm{Mallory Body} A clump of damaged proteins inside a liver cell, seen most often in alcohol-related liver disease but also in other liver conditions. It is a laboratory finding, not a diagnosis by itself.

\medterm{Mallory-Weiss Tear} A superficial tear in the lining of the lower esophagus or upper stomach, usually caused by forceful vomiting or retching. It commonly causes vomiting of blood. Many tears stop bleeding on their own within a few days, but upper endoscopy may be needed for diagnosis and treatment.

\medterm{Malnutrition} Poor health caused by too little, too much, or an imbalance of nutrients. It can impair growth, immunity, wound healing, muscle function, and organ function and may involve undernutrition, nutrient deficiencies, or excess weight.

\medterm{Malocclusion} A mismatch in the position of the upper and lower teeth or jaws that affects how the teeth meet. It may cause crowding, an overbite, an underbite, difficulty chewing, speech problems, jaw discomfort, or trouble cleaning the teeth.

\medterm{Malocclusion Of Teeth} Misalignment of the upper and lower teeth or jaws that affects the bite, chewing, speech, appearance, or oral hygiene.

\medterm{Malondialdehyde} A chemical formed when fats are damaged by oxidation. Laboratory researchers may measure it as an indirect sign of oxidative stress, but the result alone does not diagnose a disease.

\medterm{Malpighian Bodies} Renal glomeruli within Bowman's capsules, where blood filtration begins; the term may also refer historically to lymphoid nodules in the spleen.

\medterm{Malposition} An abnormal position of an organ, body part, implanted device, or fetus; in obstetrics, fetal malposition refers to an abnormal orientation of the presenting part during labor.

\medterm{Malpresentation} A fetal presentation other than the usual vertex presentation, including breech, face, brow,
shoulder, or compound presentation. It may complicate labor and influence delivery planning.

\medterm{Malrotation} Abnormal rotation or fixation of an organ during development, most often referring to abnormal rotation of the
intestine and its risk of volvulus.

\medterm{MALS} Abbreviation; see \textit{Median Arcuate Ligament Syndrome}.

\medterm{MALT Lymphoma} Alternate name; see \textit{MALT Lymphoma}.

\medterm{Malta Fever} Brucellosis, a zoonotic infection caused by Brucella species and characterized by intermittent fever, sweats, fatigue, and joint or systemic symptoms.

\medterm{Maltase} An intestinal enzyme that breaks maltose into two glucose molecules.

\medterm{Maltodextrin} A mixture of short glucose polymers produced by partial hydrolysis of starch, used as a rapidly digestible carbohydrate and as a thickener, filler, or carrier in foods and medicines.

\medterm{Maltose} A reducing disaccharide composed of two glucose molecules joined by an alpha-1,4 glycosidic bond, produced during starch digestion and hydrolyzed by maltase.

\medterm{Malt-Worker's Lung} An occupational hypersensitivity pneumonitis caused by inhalation of organic dusts associated with malting or grain processing.

\medterm{Mammaglobin A} A secretory protein expressed predominantly in mammary tissue and used as an immunohistochemical marker that can help identify breast origin in some metastatic carcinomas.

\medterm{Mammaplasty} Surgery to change the size or shape of the breast or to rebuild it after breast removal. It may enlarge, reduce, lift, or reconstruct the breast; possible risks include bleeding, infection, altered sensation, scarring, and problems with healing.

\medterm{Mammary} Relating to the breast or milk-producing gland. Mammary glands develop during puberty and pregnancy and produce milk after childbirth under hormone control. The term may describe normal breast tissue or diseases such as cysts, infection, benign growths, or breast cancer.

\medterm{Mammary Duct Ectasia} A benign inflammatory breast condition in which major subareolar milk ducts become abnormally dilated, shortened, and filled with stagnant thick lipid-rich secretions, leading to periductal chronic inflammation and fibrosis. Most common in perimenopausal and postmenopausal women, it typically presents with thick, sticky, multicolored (green, brown, or creamy) nipple discharge, nipple tenderness, or a slit-like retraction of the nipple. It can produce a firm subareolar lump that closely mimics breast cancer on mammography and physical examination. Management is typically conservative with warm compresses and reassurance, though surgical excision of affected ducts (Hadfield's procedure) is performed for persistent symptoms.

\synonyms
Duct Ectasia; Plasma Cell Mastitis; Periductal Mastitis; Comedomastitis

\medterm{Mammary Epithelium} The epithelial lining of the mammary ducts and secretory alveoli; these cells help produce, secrete, and transport milk and can give rise to many breast carcinomas.

\medterm{Mammary Gland} The milk-producing gland in the breast. It contains small milk-making lobules connected to ducts that carry milk to the nipple, and its growth and activity are controlled by hormones.

\medterm{Mammary Glands} Paired milk-producing glands in the breasts, composed of lobules and ducts that develop and function under hormonal control.

\medterm{Mammary Glands} Paired glands of the anterior chest that produce milk after childbirth. Each breast contains lobules, ducts, connective tissue, and fat and is influenced by hormonal and reproductive changes.

\synonyms
Breasts; Lactiferous Glands

\medterm{Mammary Implant} An artificial device placed in the breast for reconstruction or enlargement. Implants may contain saline or silicone gel and may be placed above or below the chest muscle. Possible complications include infection, rupture, hard scar tissue, changes in sensation, and rarely implant-associated lymphoma.

\medterm{Mammary Paget Disease} An intraepidermal spread of malignant glandular cells into the nipple epidermis, usually associated with underlying breast carcinoma and presenting as persistent eczematous or crusted nipple change.

\medterm{Mammillary Body} One of a pair of rounded nuclei on the underside of the hypothalamus that participate in memory circuits connecting the hippocampus with the thalamus; injury can contribute to memory impairment.

\medterm{Mammogram} A mammogram is a low-dose x-ray examination of breast tissue used for screening and evaluation of masses, asymmetry, calcifications, or other abnormalities.

\medterm{Mammographic Calcifications} Mineral deposits within breast tissue visualized on screening or diagnostic mammography. Macrocalcifications are typically benign (associated with fibroadenomas, cysts, or fat necrosis), whereas grouped, pleomorphic, or linear branching microcalcifications warrant biopsy to exclude ductal carcinoma in situ (DCIS) or invasive carcinoma.

\synonyms
Breast Microcalcifications; Mammographic Microcalcifications; Breast Calcifications

\medterm{Mammography} X-ray imaging of the breast used to screen for or evaluate breast abnormalities. Screening intervals
depend on age, risk, symptoms, and the guideline being followed. Standard views include
the craniocaudal and mediolateral-oblique projections.

\medterm{Mammoplasty} Plastic or reconstructive breast surgery that changes breast size, shape, or contour. It may reduce or enlarge the breasts, lift them, rebuild them after cancer surgery, or correct differences present from birth.

\medterm{Managing Pain During Labor} Pain relief during labor may include breathing and movement, support, injected or intravenous medicines, nitrous oxide, or neuraxial anesthesia, chosen according to preference and clinical circumstances.

\medterm{Mandatory Reporting} A legal or regulatory duty to report specified findings or circumstances—such as suspected child or elder abuse, certain communicable diseases, or some injuries—to designated authorities; requirements vary by jurisdiction.

\medterm{Mandible} The large movable bone forming the lower jaw and holding the lower teeth. It supports chewing and speech and contains nerves and blood vessels that supply the teeth, gums, and lower lip.

\medterm{Mandible} The large movable bone forming the lower face and holding the lower teeth. It supports chewing and speech and contains nerves and blood vessels that supply the teeth, gums, and lip.

\synonyms
Inferior Maxilla; Jawbone

\medterm{Mandibular} Relating to the mandible, the lower jawbone. It holds the lower teeth and moves at the jaw joints. A mandibular fracture can cause pain, swelling, numbness, loose teeth, or a changed bite.

\medterm{Mandibular Advancement Device} An oral appliance that holds the lower jaw or tongue forward during sleep to enlarge or stabilize the upper airway, usually for selected patients with obstructive sleep apnea or snoring.

\medterm{Mandibular Condyle} The rounded superior end of each mandibular ramus that articulates with the temporal bone at the temporomandibular joint and permits jaw movement.

\medterm{Mandibular Fracture} A break in the lower jaw bone typically caused by trauma, presenting with pain,
swelling, malocclusion, and difficulty opening the mouth; may require surgical fixation.

\medterm{Mandibular Prosthesis} An artificial implant or device used to replace, reconstruct, stabilize, or restore part of the mandible after trauma, tumor removal, congenital deformity, or other loss of jaw structure.

\medterm{Mandibulofacial Dysostosis} A congenital craniofacial disorder, classically Treacher Collins syndrome, involving underdevelopment of the mandible and cheekbones with variable abnormalities of the ears, eyes, and palate.

\medterm{Manganese} An essential trace element required in small amounts. Excessive exposure, especially by inhalation, can injure the nervous system and cause manganism, a movement disorder resembling parkinsonism.

\medterm{Mania} A distinct period of abnormally elevated, expansive, or irritable mood and increased energy or activity, with associated symptoms such as reduced need for sleep, rapid speech, racing thoughts, inflated self-esteem, or risky behavior. It causes marked impairment, requires hospitalization, or includes psychotic features; the episode usually lasts at least one week when hospitalization is not required.

\medterm{Manic} Relating to mania, a pathological episode of abnormally elevated, expansive, or irritable mood with increased energy or activity and associated changes in thought, speech, sleep, or judgment.

\medterm{Manic Disorder} Alternate name; see \textit{Bipolar Disorder}.

\medterm{Manic Episode} A distinct period of abnormally elevated, expansive, or irritable mood and increased energy or activity, lasting at least seven days unless hospitalization is required sooner or psychotic features occur. It includes associated changes such as reduced need for sleep, rapid speech, racing thoughts, inflated self-esteem, or risky behavior and causes marked impairment.

\medterm{Manic Mood} An abnormally elevated, expansive, or persistently irritable mood accompanied by increased energy or activity, reduced need for sleep, rapid speech, racing thoughts, inflated self-esteem, or risky behavior; it is characteristic of a manic episode.

\medterm{Manic-Depression} An older term for bipolar disorder, a mood disorder characterized by episodes of mania or hypomania and episodes of depression.

\medterm{Manic-Depressive Illness} Alternate name; see \textit{Bipolar Disorder}.

\medterm{Manifest} Clearly apparent or expressed; in clinical usage, a disease may manifest through particular signs or symptoms.

\medterm{Manifestation} A sign, symptom, or other observable expression of a disease or health condition. Manifestations may be noticed by the person, found during examination, or detected through laboratory tests or imaging.

\medterm{Manipulation} Deliberate physical movement of a body part or joint, often performed therapeutically or diagnostically by a trained clinician.

\medterm{Manipulation Procedure} Deliberate movement of a body part by a trained clinician to restore motion, alignment, or function. It may be performed with or without anesthesia; soreness, bleeding, nerve injury, fracture, or worsening symptoms are possible depending on the procedure.

\medterm{Mannan} A polysaccharide composed primarily of mannose units, found in the cell walls or capsules of many fungi and in some plants; fungal mannans can be recognized by the immune system and used in laboratory testing.

\medterm{Mannheimia} A genus of gram-negative bacteria that mainly infects livestock, especially the respiratory tract of cattle and sheep.

\medterm{Mannitol} A medicine that pulls water into the urine. It is given by injection in selected emergencies to reduce swelling or pressure around the brain or inside the eye and in some kidney-related situations; it can worsen dehydration or heart failure.

\medterm{Mannose} A six-carbon monosaccharide related to glucose that forms part of glycoproteins and glycolipids and is important in protein glycosylation and cell recognition.

\medterm{Mannose-Binding Lectin Deficiency} Low levels or reduced activity of a blood protein that helps the immune system recognize and attack germs. Many people have no symptoms, but some are more prone to repeated infections, especially if other immune defenses are weak.

\medterm{Mannosidosis} An inherited lysosomal storage disorder caused by deficient alpha-mannosidase activity, leading to accumulation of mannose-rich oligosaccharides and variable intellectual, skeletal, immune, hearing, and neurologic abnormalities.

\medterm{Manometer} An instrument that measures pressure in a liquid or gas. In medicine, it may measure pressure in blood vessels, the lungs, the bladder, or another body space.

\medterm{Manometry} A test that measures pressure and muscle coordination inside an organ or passage, such as the esophagus, stomach, rectum, or urinary tract. It helps identify movement and sphincter disorders.

\medterm{Mansonella Perstans} A filarial worm that causes mansonellosis, usually acquired through bites of infected midges and often producing mild or no symptoms.

\medterm{Mansonella Streptocerca} A filarial worm transmitted by biting midges that lives mainly in the skin and can cause itchy skin disease or small nodules.

\medterm{Mansonelliasis} Infection with filarial nematodes of the genus \textit{Mansonella}, transmitted by biting midges or blackflies; many infections are asymptomatic, but some cause itching, skin lesions, fever, or eosinophilia.

\medterm{Mantle Cell Lymphoma} A usually fast-growing cancer of B lymphocytes that often enlarges lymph nodes and the spleen and may involve the bone marrow, blood, or digestive tract. A characteristic chromosome change helps confirm the diagnosis and guide treatment.

\medterm{Mantle Field} A historical radiation-treatment field encompassing lymph-node regions in the neck, mediastinum, and upper chest, formerly used for Hodgkin lymphoma; modern treatment generally uses more limited involved-site fields.

\medterm{Mantoux Test} A skin test for immune sensitization to tuberculosis. A small amount of purified protein derivative is injected into the forearm, and the raised area is measured 48 to 72 hours later; a positive result does not by itself prove active disease.

\medterm{Manual} Performed by hand rather than automatically, or a written guide explaining procedures and instructions. In healthcare, manual techniques may assess, move, treat, or support body structures safely.

\medterm{Manual Removal Of Placenta} A procedure in which a clinician removes a placenta that has not separated or cannot be delivered after birth. The timing and need depend on bleeding, maternal condition, and local clinical protocols.

\medterm{Manual Vacuum Aspiration} A procedure that uses a handheld syringe and a thin tube to remove tissue from the uterus. It may treat early pregnancy loss or end a pregnancy, depending on local law and clinical circumstances, and is performed by trained clinicians.

\synonyms
MVA; Mini-Suction Curettage

\medterm{Manus} The hand, including the wrist and fingers, in anatomical terminology.

\medterm{MAOIs} A class of medicines used mainly for depression and some other conditions. They slow the breakdown of certain brain chemical messengers, but have important interactions with foods and medicines and must be prescribed and monitored carefully.

\medterm{Map-Dot-Fingerprint Type Corneal Dystrophy} A usually bilateral corneal epithelial basement membrane dystrophy producing map-like, dot-like, or fingerprint-like patterns and sometimes recurrent corneal erosions or visual blurring.

\medterm{Maple Syrup Urine Disease} An inherited disorder in which the body cannot properly break down three amino acids from protein. Harmful substances build up, causing sweet-smelling urine, poor feeding, sleepiness, abnormal movements, seizures, or brain injury without treatment.

\medterm{Mapping} Charting the location of genes or other identifiable regions on chromosomes; the term may also describe charting functional or anatomical regions in other medical contexts.

\medterm{Marantic} Relating to a sterile, noninfective process associated with wasting; marantic endocarditis refers to nonbacterial thrombotic endocarditis.

\medterm{Marasmus} Severe undernutrition caused by prolonged lack of calories and protein. It leads to marked loss of body fat and muscle, poor growth in children, weakness, and increased vulnerability to infection; urgent nutritional and medical care is needed.

\medterm{Maraviroc} A CCR5 antagonist (entry inhibitor) used in combination antiretroviral therapy for HIV-1
infection in patients with CCR5-tropic virus. Maraviroc binds to the human CCR5
coreceptor on CD4+ T cells, preventing HIV-1 gp120 binding and viral entry.

\medterm{Marble Bone Disease} Osteopetrosis, an inherited disorder of defective bone resorption causing abnormally dense but fragile bones and possible marrow failure.

\medterm{Marburg Virus} A filovirus that causes Marburg virus disease, a severe viral hemorrhagic fever in humans and some nonhuman primates. The virus is associated with Egyptian rousette fruit bats and can first spread to people after contact with infected bats. It then spreads between people through direct contact with infected blood, other body fluids, or materials contaminated with them, not through ordinary airborne spread. The disease may begin with high fever, severe headache, tiredness, and muscle aches, followed by vomiting, diarrhea, rash, bleeding, shock, or organ failure.

\medterm{Marburgvirus} A group of filoviruses that includes Marburg virus and Ravn virus, the viruses that cause Marburg virus disease. The disease is a severe viral hemorrhagic fever that can spread from fruit bats to people and between people through direct contact with infected body fluids.

\medterm{Marden-Walker Syndrome} A rare inherited disorder characterized by a distinctive facial appearance, low muscle tone, limited joint movement, poor growth, and developmental delay. Care is supportive and involves several specialties.

\medterm{Marek's Disease} A contagious herpesvirus disease of chickens that can cause nerve paralysis, tumors, weight loss, or death.

\medterm{Marfan Syndrome} An inherited connective-tissue disorder usually caused by a disease-causing variant in the \textit{FBN1} gene, which encodes fibrillin-1. It may affect the aortic root and heart valves, eyes, skeleton, joints, lungs, and dura. Features can include aortic-root enlargement, lens dislocation, tall or long-limbed stature, long fingers, scoliosis, chest-wall differences, and joint laxity; the severity varies among individuals.

\medterm{Marfanoid Habitus} A body habitus characterized by tall stature, long limbs, long fingers (arachnodactyly),
and joint laxity. Seen in Marfan syndrome, but also in other connective tissue disorders
and homocystinuria. Differentiated from Marfan syndrome by absence of aortic root
dilation and lens subluxation.

\medterm{Marfan's Syndrome} Alternate spelling; see \textit{Marfan Syndrome}.

\medterm{Marfan's Syndrome} Alternate spelling; see \textit{Marfan Syndrome}.

\medterm{Marginal Ulcer} An ulcer occurring at or near a gastrojejunal anastomosis, often after gastric bypass or other gastric surgery, where acid, ischemia, smoking, nonsteroidal anti-inflammatory drugs, or other factors may contribute.

\medterm{Marginal Zone Lymphoma} A slow-growing non-Hodgkin lymphoma that begins in immune cells in the marginal zone of lymphoid tissue.

\synonyms
MZL

\medterm{Maribavir} Maribavir is an antiviral medicine used to treat certain cytomegalovirus infections that are resistant or refractory to other treatment after transplantation.

\medterm{Marie-Strumpell Disease} Ankylosing spondylitis, a chronic inflammatory disease primarily affecting the sacroiliac joints and spine that can cause back pain, stiffness, and progressive spinal fusion.

\medterm{Marijuana} Cannabis preparations containing cannabinoids such as THC; medical effects and risks depend on dose, formulation, route, and patient factors.

\synonyms
Cannabis; Weed; Pot

\medterm{Marijuana Intoxication} Temporary changes after using cannabis, including impaired attention, memory, coordination, judgment, or perception. Anxiety, panic, drowsiness, vomiting, or psychosis may occur; severe symptoms, accidental ingestion by a child, or unsafe behavior require urgent help.

\medterm{Marijuana/Cannabis} Cannabis is a plant containing active chemicals such as THC and cannabidiol. THC can affect memory, coordination, mood, and judgment; cannabis may also interact with medicines and can cause dependence.

\medterm{Marine Animal Stings Or Bites} Injuries caused by marine animals through biting, stinging, spines, or venom; effects range from local pain and swelling to allergic, neurologic, or systemic toxicity.

\medterm{Marine Toxin} A poisonous substance produced or accumulated by marine organisms, such as some fish, shellfish, algae, or bacteria, that can cause neurologic, gastrointestinal, or other poisoning.

\medterm{Marital Disharmony} Persistent conflict or distress in a marriage or partnership that may contribute to psychological symptoms or require relationship-focused support.

\medterm{Marital Status} A demographic or social-status variable describing whether a person is unmarried, married or partnered, separated, divorced, widowed, or in another legally or culturally recognized relationship; it may be relevant to social and health assessments.

\medterm{Marital Therapy} Counseling for couples that addresses communication, conflict, emotional distress, intimacy, and relationship problems. It may help partners improve interaction and make informed decisions about their relationship.

\medterm{Marked} Pronounced, conspicuous, or substantially greater than usual; in a clinical report, the exact meaning depends on the finding and the comparison standard.

\medterm{Marker} An identifiable biological feature used to track a gene, cell, disease, or process; a genetic marker is a detectable DNA region inherited with a nearby gene, while a blood or tumor marker is a measurable substance associated with a condition.

\medterm{Maroteaux-Lamy Syndrome} A rare inherited disorder in which certain complex sugars build up in the body. It can cause short stature, stiff joints, enlarged liver or spleen, cloudy corneas, heart-valve disease, breathing problems, and coarse facial features. Intelligence is often normal. Treatment may include enzyme replacement, supportive care, or stem-cell transplantation.

\synonyms
Mucopolysaccharidosis Type Vi; MPS VI; Arylsulfatase B Deficiency

\medterm{Marrow} The soft tissue inside bones, including red marrow that produces blood cells and yellow marrow rich in fat.

\medterm{Marrow Hyperplasia} Increased activity and numbers of blood-forming cells in the bone marrow. It may occur when the body is replacing blood lost or destroyed, or when tissues are not receiving enough oxygen.

\medterm{Marshall Syndrome} A rare periodic-fever syndrome, usually associated with PFAPA features, including recurrent fever, aphthous ulcers, pharyngitis, and cervical adenitis, often beginning in childhood.

\medterm{Marsupialization} A surgical technique that opens a cyst or abscess and sutures its edges to the surrounding skin or mucosa to permit continuing drainage.

\medterm{Martin-Bell Syndrome} An older name for fragile X syndrome, an inherited neurodevelopmental disorder caused by CGG-repeat expansion in the FMR1 gene.

\medterm{MASA Syndrome} An X-linked neurodevelopmental disorder characterized by intellectual disability, absent or limited speech, shuffling gait, spasticity, and adducted or clasped thumbs; it may also include short stature and lumbar lordosis. The name is an older acronym and terminology should be used sensitively.

\medterm{Masculinization} Development of male-typical physical characteristics, caused by normal puberty, androgen excess, androgen treatment, or certain disorders.

\medterm{MASH} Abbreviation; see \textit{Metabolic Dysfunction-Associated Steatohepatitis}.

\medterm{Mask} A device placed over the face to deliver oxygen, anesthesia, ventilation, or protection; the meaning depends on clinical context.

\medterm{Masked Data} Data altered, obscured, or replaced so that identifying values or sensitive details are hidden while selected analytical properties may be retained; masking is used to protect privacy and can be reversible or irreversible.

\medterm{Masklike Face} An expressionless face with little spontaneous animation, also called masklike facies; it may occur in Parkinson disease, myotonic dystrophy, and other disorders.

\medterm{MASLD} Abbreviation; see \textit{Metabolic Dysfunction-Associated Steatotic Liver Disease}.

\medterm{Masochism} Sexual or nonsexual arousal or gratification associated with receiving pain, humiliation, or degradation; distress or impairment determines clinical significance.

\medterm{Masoprocol} A synthetic derivative of nordihydroguaiaretic acid that inhibits lipoxygenase pathways and has been investigated for topical treatment of actinic keratoses and for antineoplastic activity.

\medterm{Mass} A discrete area of abnormal tissue that is larger or more substantial than a nodule. A mass may be benign or malignant and can occur in many organs.

\medterm{Mass Casualty Incident} An event causing more injured or ill people than nearby medical services can immediately manage. Health teams use triage to prioritize urgent treatment, transport, and limited supplies.

\medterm{Mass Screening} Testing a large group of people who do not yet have symptoms to find disease or risk early. A useful screening program has a suitable test, an effective follow-up pathway, and benefits that outweigh harms.

\medterm{Mass Spectrometer} A laboratory instrument that identifies chemicals by giving their particles an electrical charge and separating them according to size and charge. It can measure medicines, toxins, proteins, and other substances.

\medterm{Mass Spectrometry} An analytical method that identifies and measures molecules by their mass-to-charge ratios; coupled with chromatography, it is widely used in clinical chemistry, toxicology, pharmacology, and proteomics.

\medterm{Massage} Manual manipulation of soft tissues using techniques such as stroking, kneading, or pressure, used for relaxation or symptom relief; it is not a substitute for diagnosis or treatment of an underlying disease.

\medterm{Massage Therapy} Massage therapy uses hands-on manipulation of soft tissues to promote relaxation or help manage selected pain and movement problems.

\medterm{Masseter} A powerful muscle of chewing that elevates the mandible and lies over the lateral surface of the mandibular ramus.

\medterm{Masseter Muscle} A strong chewing muscle running from the cheekbone to the lower jaw. It raises the jaw to close the mouth and may become painful or enlarged with teeth grinding.

\medterm{Massive Hemoptysis} Life-threatening coughing of blood that threatens the airway, breathing, or circulation; a fixed
volume threshold is not required. It requires urgent airway assessment, stabilization, and
definitive treatment of the bleeding source.

\medterm{Massive Ovarian Edema} Marked stromal edema causing enlargement of one or both ovaries, usually due to partial torsion or hormonal stimulation. It can mimic an ovarian neoplasm and may cause pain or virilization.

\medterm{Massive Pulmonary Embolus} A high-risk, catastrophic pulmonary embolism characterized by extensive blood clot obstruction of the main pulmonary arterial tree (or major lobar arteries), leading to acute right ventricular failure and hemodynamic collapse. It is clinically defined by sustained systemic hypotension (systolic blood pressure below 90 mmHg or a drop of at least 40 mmHg), obstructive shock, or cardiac arrest. Patients present with sudden severe breathlessness, syncope (fainting), cyanosis, tachycardia, and elevated jugular venous pressure. Emergency management requires immediate supportive resuscitation, advanced mechanical clot removal (catheter-directed or surgical embolectomy), or systemic thrombolytic therapy (clot-dissolving medications).

\synonyms
Massive PE; High-Risk Pulmonary Embolism; Fulminant Pulmonary Embolism

\medterm{Massive Transfusion Protocol} A coordinated protocol for life-threatening hemorrhage that provides rapid, balanced
delivery of red cells, plasma, and platelets with monitoring and correction of
fibrinogen, calcium, temperature, and acid-base abnormalities.

\medterm{Mast Cell} An immune cell found in tissues that releases substances involved in allergy, inflammation, and defense against some germs. Its products can cause itching, swelling, mucus production, and changes in blood pressure.

\medterm{Mast Cell Activation Syndrome} A condition involving recurrent symptoms from inappropriate release of mast-cell mediators, such as flushing,
itching, hives, abdominal symptoms, wheezing, or low blood pressure. Diagnosis requires compatible episodes,
objective evidence of mediator release, and response to appropriate treatment; many similar conditions must be excluded.

\medterm{Mast Cell Sarcoma} An exceptionally rare, aggressive neoplasm composed of immature or atypical mast cells forming a destructive mass in tissue or bone.
Unlike systemic mastocytosis, it is characterized by a localized sarcomatous tumor and may progress to disseminated mast-cell disease or acute leukemia.
Clinical manifestations depend on the involved site and may include pain, swelling, organ dysfunction, or mediator-related symptoms such as flushing and hypotension.
Diagnosis requires histopathology, mast-cell immunophenotyping, and systemic staging; management generally involves specialist-directed intensive antineoplastic therapy.

\medterm{Mastalgia} Breast pain or tenderness, which may be cyclical with menstruation or noncyclical. Causes include hormonal changes, cysts, infection, trauma, and musculoskeletal pain.

\medterm{Mast-Cell Sarcoma} An extremely rare cancer made of abnormal mast cells, immune cells involved in allergic and inflammatory reactions.

\medterm{Mastectomy} Surgical removal of all or part of a breast, performed to treat or prevent selected breast cancers and sometimes for other breast conditions. Depending on the operation, the nipple, areola, and nearby lymph nodes may also be removed.

\medterm{Masticate} To chew food, using the teeth, tongue, and muscles of mastication to mechanically break it down and mix it with saliva before swallowing.

\medterm{Mastication} Chewing food with the teeth and moving it with the tongue before swallowing. It breaks food into smaller pieces and mixes it with saliva, making it easier to swallow and begin digesting.

\medterm{Masticatory Muscle} One of the muscles that moves the mandible during chewing, chiefly the masseter, temporalis, medial pterygoid, and lateral pterygoid muscles.

\medterm{Mastitis} Inflammation of breast tissue, often occurring during lactation but also possible outside lactation. It may
cause localized pain, warmth, swelling, and fever; infection and inflammatory causes
must be distinguished clinically.

\synonyms
Puerperal Mastitis; Breast Inflammation

\medterm{Mastocytoma} A localized collection or tumor of mast cells, usually presenting in the skin as a solitary cutaneous mastocytoma; rubbing the lesion can cause itching, redness, or swelling from mediator release.

\medterm{Mastocytosis} A disorder in which too many mast cells, immune cells involved in allergic reactions, collect in the skin or internal organs. It may cause itchy spots, flushing, abdominal symptoms, wheezing, or low blood pressure.

\medterm{Mastodynia} Pain or tenderness in one or both breasts, also called mastalgia. It may follow the menstrual cycle or result from injury, infection, medicines, cysts, or other breast conditions; a new persistent lump, skin change, or bloody discharge needs assessment.

\medterm{Mastoid} Relating to the mastoid part of the temporal bone behind the ear. This area contains air spaces connected with the middle ear and can become infected, causing pain, swelling, or drainage.

\medterm{Mastoid Antrum} An air-containing cavity in the mastoid portion of the temporal bone that communicates with the middle-ear attic.

\medterm{Mastoid Bone} The portion of the temporal bone located behind the ear.

\medterm{Mastoid Process} The conical bony projection of the temporal bone behind the ear, containing air cells and providing attachment for several neck muscles; infection of its air cells is mastoiditis.

\medterm{Mastoidectomy} Surgery to remove infected or diseased air cells in the bone behind the ear. It may treat chronic middle-ear infection, cholesteatoma, mastoiditis, or complications that have spread beyond the ear.

\medterm{Mastoiditis} Mastoiditis is infection and inflammation of the mastoid air cells behind the ear, usually following middle-ear infection and potentially causing serious spread.

\medterm{Mastopexy} Breast-lift surgery that removes extra skin and reshapes breast tissue to raise the breasts and nipples. It can improve shape but does not necessarily reduce breast size; scars, altered sensation, asymmetry, and changes in breastfeeding ability are possible risks.

\synonyms
Breast Lift

\medterm{Masturbation} Self-stimulation of the genitals for sexual pleasure. It is a common human behavior; it is generally
not harmful unless it causes injury, distress, or interferes with daily life.

\synonyms
Autoeroticism; Solitary Sex

\medterm{Materia Medica} An older term for the body of medicinal substances and information about their properties and uses.

\medterm{Maternal} Maternal means relating to a mother or the pregnant person, especially in connection with pregnancy, childbirth, or postpartum health.

\medterm{Maternal Age} Maternal age is the age of the pregnant person and can influence fertility, pregnancy risks, screening choices, and obstetric care.

\medterm{Maternal Expulsive Efforts} The pregnant person's efforts to push the baby down and out during the second stage of labor. Pushing may follow the body's urge or be guided by the birth team, depending on the situation and chosen pain relief.

\synonyms
Bearing Down; Second Stage Expulsive Efforts

\medterm{Maternal Health} The physical and mental health of a person during pregnancy, childbirth, and the postpartum period. Care includes prevention, antenatal monitoring, safe delivery, and treatment of complications.

\medterm{Maternal Mortality} Maternal mortality is death during pregnancy or within the defined postpartum period from a pregnancy-related cause, measured as a population health indicator.

\medterm{Maternal Mortality Rate} The number of maternal deaths related to pregnancy per 100,000 live births in a specified population and period.

\medterm{Maternal Obesity} Obesity present before or during pregnancy. It can increase the risk of high blood pressure, diabetes, cesarean birth, and health problems for the pregnant person or baby.

\medterm{Maternal Physiologic Changes In Pregnancy} Normal body changes during pregnancy that support the growing fetus and prepare for birth. They affect blood volume, heart function, breathing, digestion, kidneys, hormones, breasts, skin, and the musculoskeletal system.

\medterm{Maternal Sepsis} Organ dysfunction caused by a dysregulated maternal response to infection during
pregnancy, childbirth, or the postpartum period. It may arise from uterine, urinary,
respiratory, wound, or other infection and requires prompt cultures, antimicrobial
therapy, source control, and hemodynamic support.

\medterm{Maternal Serum Alpha-Fetoprotein} A blood test during pregnancy, usually at 15--20 weeks, that measures a protein made by the fetus. An unusually high or low result may be linked to an incorrect pregnancy date, more than one fetus, or certain birth defects or chromosome conditions. It is a screening test, not a diagnosis, and abnormal results need follow-up.

\medterm{Maternal Vaccine} A vaccine given during pregnancy to protect the pregnant person and, through transferred antibodies, provide early protection to the fetus or newborn; recommended products and timing depend on current guidance.

\medterm{Maternity} The state or service relating to pregnancy, childbirth, and the care of the mother and newborn.

\medterm{Maternity Services} Healthcare for pregnancy, labor and birth, the postpartum period, and the newborn. Services may include prenatal checks, midwifery or obstetric care, delivery support, breastfeeding help, and emergency care.

\medterm{Materteral} Relating to a maternal aunt or to the genetic relationship between an aunt and her niece or nephew; synonymous with materterine.

\medterm{Mathematics Disorder} A learning disorder involving persistent difficulty understanding numbers, calculation, quantity, or mathematical reasoning despite appropriate instruction.

\medterm{Mating} Sexual pairing or reproduction between organisms that permits exchange or transfer of gametes and may lead to fertilization.

\medterm{Matrix Band} A thin dental strip placed around a prepared tooth to temporarily replace a missing wall and contain restorative material while a filling or crown buildup is formed.

\medterm{Matrix Metalloproteinase} An enzyme that cuts proteins in the extracellular matrix surrounding cells. It helps with normal tissue repair and remodeling, but excessive or poorly controlled activity can promote inflammation, tissue destruction, or cancer spread.

\medterm{Maturation} The process through which cells, tissues, organs, or a person develop and become fully functional. Normal maturation may involve changes in size, structure, ability, hormone activity, and response to the environment.

\medterm{Mature Onset Diabetes} An older term generally referring to type 2 diabetes mellitus, which often develops in adulthood, though diabetes classification depends on cause rather than age alone.

\medterm{Mature Teratoma} A germ-cell tumor composed of well-differentiated tissues from one or more embryonic germ layers; many are benign, but behavior depends on the site, age, and tumor type, and some can undergo malignant transformation.

\medterm{Maxilla} The paired bones forming the upper jaw and supporting the upper teeth. They also form parts of the cheeks, nose, eye sockets, hard palate, and the air-filled sinuses behind the cheeks.

\medterm{Maxilla} The paired bones forming the upper jaw, the front part of the hard palate, and parts of the walls of the nose and eye sockets.

\synonyms
Superior Maxilla; Upper Jawbone

\medterm{Maxillary} Relating to the maxilla, the upper jawbone. It supports the upper teeth, forms part of the eye socket and roof of the mouth, and contains the maxillary sinuses. Problems may include tooth infection, sinusitis, or facial fractures.

\medterm{Maxillary Sinus} The maxillary sinus is a paired, air-filled cavity in the maxilla beneath each orbit. It drains into the nasal cavity and may become inflamed or infected in sinusitis.

\medterm{Maxillary Sinusitis} Inflammation or infection of the air-filled spaces behind the cheekbones. It may cause cheek or tooth pain, blocked nose, thick nasal discharge, reduced smell, cough, or fever and can follow a cold, allergy, or dental infection.

\medterm{Maxillofacial} Relating to the jaw, face, mouth, and associated bones and soft tissues. Maxillofacial care may include treatment of injuries, congenital differences, infections, and tumors.

\medterm{Maxillofacial Prosthesis} An artificial replacement for part of the face, jaw, mouth, nose, eye area, or ear after injury, cancer surgery, or a birth difference. Examples include devices that close a palate opening or replace an eye, nose, or ear.

\medterm{Maximal Breathing Capacity} The greatest volume of air that can be moved in and out of the lungs in a specified period during maximal effort; it is related to but distinct from standard spirometric measures.

\medterm{Maximum} The maximum is the greatest value in a set of measurements or the highest level reached by a variable, such as a recorded temperature, blood pressure, or drug concentration.

\medterm{Maximum Plasma Concentration} The highest measured plasma concentration after a dose. Cmax and the time to reach it
are used to describe absorption rate, peak-related toxicity, and comparative
pharmacokinetics.

\synonyms
Cmax

\medterm{Maximum Thickness} Maximum thickness is the greatest measured thickness of a structure or specimen. It may be reported in imaging, pathology, or assessment of a lesion, organ, or tissue layer.

\medterm{Maximum Tolerated Dose} The highest dose of a drug that produces acceptable, manageable toxicity without causing
unacceptable adverse effects, determined through dose-escalation studies, most commonly
in Phase I clinical trials. The MTD identifies the dose range for subsequent efficacy
studies.

\medterm{Mayer-Rokitansky-KüSter-Hauser Syndrome} A condition present from birth in which the uterus and much of the vagina do not develop, although the ovaries and usual female puberty changes develop normally. It often first becomes noticeable when menstruation does not begin. Kidney, hearing, or spine abnormalities may also occur. Treatment and fertility options depend on the individual’s anatomy and goals.

\synonyms
MRKH Syndrome; Müllerian Agenesis; Vaginal Agenesis

\medterm{May-Hegglin Anomaly} An inherited MYH9-related disorder with macrothrombocytopenia, giant platelets, and characteristic leukocyte inclusions; hearing, kidney, or cataract abnormalities may occur.

\medterm{Mayo Scissors} Strong surgical scissors designed to cut fascia, muscle, sutures, and other tough materials. Their broad blades provide force and control, and they may be straight or curved for different surgical tasks.

\medterm{Maytansine} Maytansine is a potent antimitotic compound that binds tubulin and inhibits microtubule assembly. Its derivatives are used as cytotoxic payloads in some antibody--drug conjugates.

\medterm{May-Thurner Syndrome} Compression of the left pelvic vein by the right pelvic artery, which can slow blood flow and increase the risk of a clot or swelling in the left leg.

\medterm{Maze Procedure} A heart operation for atrial fibrillation. The surgeon creates carefully placed lines of scar tissue in the upper heart chambers to block abnormal electrical signals and help restore a regular rhythm. It may be performed with other heart surgery; recurrence, bleeding, stroke, or the need for a pacemaker are possible.

\synonyms
Cox-Maze Procedure; Surgical Ablation for Atrial Fibrillation; Maze Operation

\medterm{Mazindol} Mazindol is a sympathomimetic appetite suppressant that affects central catecholamine signaling. It has been used for short-term obesity treatment in some countries and may cause stimulant-related adverse effects.

\medterm{MB ChB} Abbreviation for Bachelor of Medicine and Bachelor of Surgery, a medical degree awarded in many countries.

\medterm{MC Lymphoma} Abbreviation; see \textit{Mantle Cell Lymphoma}.

\medterm{MCA Peak Systolic Velocity} An ultrasound measurement of blood speed in a major artery of the fetus's brain. An unusually high value may suggest fetal anemia and can help guide monitoring or further testing during pregnancy.

\medterm{MCAD Deficiency} An inherited disorder in which the body cannot properly break down certain medium-chain fats for energy.

\medterm{McArdle Disease} A rare inherited muscle disorder in which muscles cannot release usable energy from stored sugar during exercise. Activity can cause early muscle pain, cramps, exhaustion, or dangerous muscle breakdown with dark urine. Some people improve after slowing down briefly and restarting gently, called the second-wind phenomenon. Avoiding sudden strenuous exercise helps prevent attacks.

\synonyms
GSD V; Myophosphorylase Deficiency; Glycogen Storage Disease Type V

\medterm{McBurney Sign} Tenderness when pressing on the lower right side of the belly about halfway between the belly button and the hip bone, which is a key sign of an infected appendix (appendicitis).

\medterm{McBurney's Point} A location in the lower-right abdomen where tenderness may occur when the appendix is inflamed. Pain there can support a diagnosis of appendicitis but cannot confirm it by itself.

\medterm{Mcc Gene} A gene that helps regulate cell growth and is considered a tumour-suppressor gene. Loss or alteration of its function has been found in some colorectal and other cancers, but the finding alone does not diagnose cancer.

\medterm{McCune-Albright Syndrome} A condition caused by a change affecting hormone signals in some body cells. It may cause irregular skin coloring, fragile or unevenly growing bones, and early puberty or other hormone problems.

\medterm{McDonald Cerclage} A common operation that places a strong stitch around the cervix through the vagina. It may help keep the cervix closed when there is a high risk of very early birth; timing and suitability depend on the pregnancy history and examination.

\medterm{MCH} The average amount of hemoglobin in one red blood cell, reported in a complete blood count. Hemoglobin carries oxygen, and the result helps classify some types of anemia.

\medterm{MCHC} The average concentration of hemoglobin inside red blood cells, reported in a complete blood count. It helps clinicians assess and classify anemia together with the other blood-count results.

\medterm{MCL Injury} Abbreviation; see \textit{Medial Collateral Ligament Injury}.

\medterm{McMurray Test} A knee examination in which the leg is flexed, rotated, and extended to place stress on the menisci. Pain or a palpable click may suggest a meniscal lesion but is not diagnostic alone.

\medterm{McRoberts Maneuver} A positioning maneuver for managing shoulder dystocia where the mother’s legs are
hyperflexed onto her abdomen, with thighs against the abdomen. This flattens the sacrum,
increases the anteroposterior diameter of the pelvis, and rotates the symphysis pubis
superiorly. Often combined with suprapubic pressure.

\medterm{MCRPC} Abbreviation; see \textit{Metastatic Castration-Resistant Prostate Cancer}.

\medterm{MCV} Mean corpuscular volume: the calculated average volume of a red blood cell, derived from hematocrit and red-cell count and reported as part of a complete blood count.

\medterm{MDMA} 3,4-Methylenedioxymethamphetamine, a psychoactive stimulant and hallucinogen commonly known as ecstasy or molly; it can produce serious cardiovascular, neurologic, and temperature-regulation effects.

\medterm{MDS} Abbreviation; see \textit{Myelodysplastic Syndromes}.

\medterm{Mead Acid} An omega-9 fatty acid made when the body lacks enough essential omega-6 and omega-3 fatty acids. Increased levels can support laboratory assessment of essential-fatty-acid deficiency, although results must be interpreted with nutrition and illness history.

\medterm{Meadow Saffron} Colchicum autumnale, a plant containing colchicine; ingestion can cause severe gastrointestinal illness, multiorgan toxicity, and death.

\medterm{Mean} The mean is the arithmetic average, calculated by adding all values and dividing by the number of values. It is commonly used to summarize clinical measurements but can be distorted by extreme values.

\medterm{Mean Arterial Pressure} An estimate of the average pressure pushing blood through the body's organs during each heartbeat. Clinicians use it with other findings to judge whether organs are receiving enough blood.

\medterm{Mean Cell Volume} The average volume of a red blood cell, calculated from the hematocrit and red-cell count and reported as MCV in a complete blood count.

\medterm{Mean Corpuscular Hemoglobin Concentration} The average concentration of oxygen-carrying hemoglobin inside red blood cells. It is part of a complete blood count and helps classify the cause of anemia.

\medterm{Mean Corpuscular Volume} The average size of red blood cells, reported in a complete blood count. Small cells or large cells can point toward different causes of anemia and must be interpreted with the other blood results.

\medterm{Mean Diameter} Mean diameter is the average diameter of a structure or lesion, usually calculated from one or more measured dimensions. The method used should be specified when measurements are clinically compared.

\medterm{Mean Thickness} Mean thickness is the average thickness calculated from multiple measurements of a tissue, structure, or specimen. It is used in imaging, pathology, and physiologic measurements.

\medterm{Measles} A highly contagious viral infection causing fever, cough, runny nose, red eyes, and a rash that spreads from the face downward. It can cause pneumonia, brain inflammation, or other serious complications.

\medterm{Measles Virus} A negative-sense RNA virus that causes measles, a highly contagious respiratory illness.
The virus is transmitted by respiratory droplets and is one of the most infectious
pathogens known. Symptoms include fever, cough, coryza, conjunctivitis, and rash.

\medterm{Measurable Residual Disease} Cancer cells that remain after treatment and can be detected with a sensitive, validated test. The amount may be too small to see on scans or by routine examination, but it can help estimate relapse risk and guide further treatment.

\medterm{Measurement} Measurement is the quantitative determination of a physical, biologic, or clinical characteristic using a defined method and unit. Interpretation depends on accuracy, precision, and the clinical context.

\medterm{Meatal} Relating to a meatus, a natural opening or passage in the body. It commonly refers to the opening of the urethra, the ear canal, or one of the passages inside the nose. Meatal narrowing can obstruct urine or air flow.

\medterm{Meatal Stenosis} Narrowing of the urethral meatus, which may cause a deflected or narrow urinary stream, spraying, dysuria, or difficulty emptying the bladder.

\medterm{Meatus} A natural opening or passage into or out of a body structure. Examples include the opening of the urethra, the ear canal, and the small passages inside the nose.

\medterm{Meatuses} Natural openings or passages in the body. In the nose, the meatuses are channels beneath the nasal folds that help air move through the nose and allow drainage from the sinuses and tear ducts.

\medterm{Mebendazole} A medicine used to treat several intestinal worm infections. It disrupts structures the worms need to absorb nutrients and survive; treatment choices and dosing depend on the parasite, age, pregnancy, and local guidance.

\medterm{Mebrofenin Scan} A nuclear-medicine scan that follows a tracer through the liver and bile ducts. It can assess gallbladder function, bile blockage, leaks, or selected causes of jaundice and abdominal pain.

\medterm{Mecamylamine} Mecamylamine is an orally active ganglionic nicotinic-receptor blocker that reduces sympathetic nerve transmission and lowers blood pressure. It has been used for severe hypertension and may cause postural hypotension, dizziness, constipation, and dry mouth.

\medterm{Mecarzole} Mecarzole is a synonym for carbendazim, a benzimidazole fungicide that interferes with fungal microtubule formation. It is an agricultural chemical rather than a routine human medicine and may be toxic if exposure occurs.

\medterm{ME/CFS} Abbreviation; see \textit{Myalgic Encephalomyelitis}.

\medterm{Mechanical Circulatory Support} Machines or pumps that help the heart move blood when it cannot pump enough on its own. They may be used temporarily during severe illness or longer term while awaiting, or instead of, a heart transplant.

\medterm{Mechanical Insufflation-Exsufflation} A cough-assist device that gives a breath in followed quickly by suction-like pressure out. It helps people with weak breathing muscles clear mucus, especially in neuromuscular disease or spinal-cord injury.

\medterm{Mechanical Ventilation} Breathing support from a machine that pushes air and oxygen into the lungs and helps remove carbon dioxide. It may be delivered through a mask or a breathing tube when illness or injury prevents adequate breathing.

\medterm{Mechanical Ventilator} A mechanical ventilator is a machine that supports or replaces spontaneous breathing by delivering controlled or assisted breaths, usually through a mask or an endotracheal tube. Settings are adjusted to oxygenation, ventilation, mechanics, and patient needs.

\medterm{Mechanism Of Death} The physiologic or biochemical derangement that causes death (the final failure) as
opposed to the underlying disease or injury (cause of death) and the circumstances
(manner). Examples: ventricular fibrillation, cerebral hemorrhage, sepsis, respiratory
failure, exsanguination, cardiac tamponade.

\medterm{Mechanism Of Toxicity} The way a substance causes injury or illness in the body. It may interfere with an enzyme or receptor, damage cell membranes or DNA, disrupt energy production, or trigger an abnormal immune response. Understanding the mechanism helps guide prevention and treatment.

\medterm{Mechanoreceptor} A sensory receptor that detects physical forces such as touch, pressure, stretch, vibration, sound, or movement. It converts mechanical stimulation into electrical signals carried by nerves to the brain or spinal cord.

\medterm{Mechlorethamine} Mechlorethamine is an alkylating antineoplastic agent that cross-links DNA and inhibits cell division. Topical gel is used for cutaneous T-cell lymphoma; systemic use is limited by toxicities such as myelosuppression and tissue injury.

\medterm{Meckel Diverticulectomy} Surgical removal of a Meckel diverticulum, usually performed when it causes bleeding, inflammation, obstruction, perforation, or is found incidentally in a high-risk setting.

\medterm{Meckel Diverticulum} The most common congenital anomaly of the gastrointestinal tract, occurring as a true diverticulum (containing all intestinal wall layers) on the antimesenteric border of the distal ileum due to incomplete closure of the embryonic vitelline (omphalomesenteric) duct. It classically follows the ``Rule of 2s'': found in approximately 2\% of the population, located roughly 2 feet proximal to the ileocecal valve, measuring about 2 inches in length, and frequently containing ectopic tissue (most commonly ectopic gastric mucosa or ectopic pancreatic tissue). Acid secretion from ectopic gastric mucosa causes ulceration and painless rectal bleeding (hematochezia), especially in young children. Other complications include diverticulitis (which closely mimics acute appendicitis), bowel obstruction, and intussusception. Diagnosis of bleeding diverticula is confirmed by a Meckel scan (technetium-99m pertechnetate scintigraphy), and symptomatic cases are treated with surgical diverticulectomy.

\synonyms
Meckel's Diverticulum; Vitelline Duct Remnant; Omphalomesenteric Diverticulum; Congenital Ileal Diverticulum

\medterm{Meckel-Gruber Syndrome} A severe inherited birth disorder that causes major abnormalities of the brain, kidneys, and limbs. Typical findings include a sac of brain tissue at the back of the skull, enlarged cystic kidneys, and extra fingers or toes. Poor lung development is common, and affected pregnancies usually do not survive to birth.

\synonyms
Meckel Syndrome; Dysencephalia Splanchnocystica

\medterm{Meckel's Diverticulum} Alternate name; see \textit{Meckel Diverticulum}.

\medterm{Meclizine} Meclizine is a first-generation antihistamine with anticholinergic and vestibular-suppressant effects. It is used to prevent or treat motion sickness and vertigo and may cause drowsiness, dry mouth, and blurred vision.

\medterm{Meclofenamate Overdose} Taking too much meclofenamate, a nonsteroidal anti-inflammatory drug, may cause nausea, vomiting, stomach pain or bleeding, drowsiness, kidney injury, seizures, or metabolic acidosis. Urgent medical or poison-control advice is needed after a suspected overdose.

\medterm{Meconium} The thick, dark first stool of a newborn, composed of intestinal secretions, shed cells, mucus, and swallowed amniotic material.

\medterm{Meconium Aspiration} Breathing meconium, a newborn's first stool, into the airways before or around birth. It can block airflow and inflame the lungs, causing breathing difficulty and low oxygen.

\medterm{Meconium Aspiration Syndrome} Breathing problems in a newborn after meconium, the first stool, enters the airways before or during birth. It can block airways and inflame the lungs, causing rapid breathing, low oxygen, and need for respiratory support.

\medterm{Meconium Ileus} A blockage of a newborn’s small intestine by unusually thick first stool. It is strongly associated with cystic fibrosis and may cause a swollen abdomen, green vomiting, and failure to pass stool soon after birth.

\medterm{Meconium Peritonitis} Inflammation of the tissue lining a fetus’s abdomen after a hole in the intestine allows first stool to escape before birth. It can cause abdominal fluid or calcium deposits and may be linked to bowel blockage or cystic fibrosis.

\medterm{Meconium-Stained Amniotic Fluid} Green or yellow amniotic fluid caused by passage of fetal stool before birth. It may occur with a mature or distressed fetus and increases the risk of breathing problems after delivery; newborns are assessed and supported as needed.

\medterm{MECP2 Duplication Syndrome} A rare X-linked genetic disorder caused by extra copies or increased dosage of the MECP2 gene. It often causes developmental delay, low muscle tone, seizures, recurrent infections, and breathing or feeding problems, especially in boys.

\medterm{Medial} Medial means toward or situated nearer the body's midline; it contrasts with lateral and is used to describe anatomical position, movement, or imaging findings.

\medterm{Medial Collateral Ligament Injury} A stretch or tear of the strong ligament on the inner side of the knee, often caused by a blow to the outer knee or a twisting injury. It causes inner-knee pain and swelling; severe tears may make the knee unstable.

\medterm{Medial Cuneiform} The largest of three small bones in the inner middle part of the foot. It joins the navicular bone and the first toe bone, helps support the arch, and forms part of the joint that can be injured in a Lisfranc injury.

\medterm{Medial Epicondylitis} Alternate medical name; see \textit{Golfer's Elbow}.

\medterm{Medial Meniscus} The medial meniscus is a crescent-shaped fibrocartilage pad between the medial femoral condyle and medial tibial plateau. It distributes load and contributes to knee stability; tears may cause joint-line pain, swelling, or locking.

\medterm{Medial Rotation} Turning a limb toward the body's midline. For example, turning the thigh so the knee and foot point inward is medial rotation; the movement is assessed at joints such as the shoulder and hip.

\medterm{Medial Tibial Stress Syndrome} Alternate name; see \textit{Shin Splints}.

\medterm{Median} Located in or relating to the middle plane or midpoint of a structure. In medicine, it describes anatomy, measurements, statistical values, nerve positions, or locations between two sides.

\medterm{Median Arcuate Ligament Syndrome} Compression of the celiac artery and nearby nerves by the median arcuate ligament, potentially causing post-meal abdominal pain and weight loss.

\synonyms
Dunbar Syndrome; MALS

\medterm{Median Nerve} A major nerve running from the neck through the arm and forearm to the hand. It helps bend the wrist and fingers, moves the thumb, and provides feeling to the thumb, index, middle, and part of the ring finger; compression can cause carpal-tunnel symptoms.

\medterm{Median Nerve Compression} Alternate name; see \textit{Carpal Tunnel Syndrome}.

\medterm{Median Plane} A sagittal plane that passes directly through the midline of the body, dividing it into
equal right and left halves.

\medterm{Median Rhomboid Glossitis} A smooth, red, diamond-shaped depapillated area on the dorsum of the tongue caused by
chronic Candida infection, usually asymptomatic and managed with antifungal therapy.

\medterm{Median Survival} The time by which half of the people in a study have experienced the specified event, such as death. It is a summary measure of survival in a defined population.

\medterm{Mediastinal Cyst} A mediastinal cyst is a fluid- or tissue-filled cystic lesion in the mediastinum. Common types include bronchogenic, thymic, and pericardial cysts; symptoms depend on size and compression, and imaging guides evaluation.

\medterm{Mediastinal Disease} A disorder affecting structures in the mediastinum, including the heart, great vessels, thymus, lymph nodes, nerves, and tumors or cysts within the central chest.

\medterm{Mediastinal Emphysema} Mediastinal emphysema, or pneumomediastinum, is the presence of free air in the mediastinum. It can follow alveolar rupture, trauma, surgery, or esophageal injury and may cause chest pain, neck crepitus, or respiratory compromise.

\medterm{Mediastinal Hemorrhage} Bleeding into the mediastinum, commonly caused by trauma, vascular rupture, surgery, or invasive procedures, and potentially causing airway or cardiovascular compression.

\medterm{Mediastinal Mass} An abnormal growth or collection of tissue in the mediastinum, the central compartment of the chest. Causes include tumors, enlarged lymph nodes, cysts, and vascular abnormalities.

\medterm{Mediastinal Neoplasm} A mediastinal neoplasm is an abnormal growth in the central space of the chest between the lungs; it may arise from lymph nodes, thymus, nerves, or other tissues.

\medterm{Mediastinal Pleura} The mediastinal pleura is the portion of parietal pleura lining the lateral surfaces of the mediastinum. It separates the mediastinal contents from the pleural cavity and reflects around the lung root.

\medterm{Mediastinal Shift} Mediastinal shift is displacement of the mediastinal structures away from their usual position. It may result from increased pressure on one side, such as a tension pneumothorax, or loss of volume on the other, such as lung collapse.

\medterm{Mediastinal Tumor} A mediastinal tumor is a mass in the central chest compartment between the lungs and may be benign or malignant.

\medterm{Mediastinitis} Mediastinitis is inflammation or infection of the central chest compartment, often after esophageal rupture, surgery, trauma, or spread of infection.

\medterm{Mediastinoscopes} Mediastinoscopes are slender endoscopic instruments used during mediastinoscopy to inspect the mediastinum and obtain tissue samples, especially from mediastinal lymph nodes. The procedure is performed through a small incision near the sternum.

\medterm{Mediastinoscopy} A surgical procedure that accesses the mediastinum through a small incision near the sternum to
inspect or biopsy mediastinal lymph nodes or masses. It is one option for staging lung cancer
when less-invasive sampling is nondiagnostic or unsuitable.

\medterm{Mediastinoscopy With Biopsy} A procedure using an instrument inserted through a small incision near the sternum to inspect the mediastinum and obtain lymph-node or mass samples for diagnosis and staging.

\medterm{Mediastinum} The central part of the chest between the lungs. It contains the heart, major blood vessels, windpipe, food pipe, thymus, nerves, and lymph nodes.

\medterm{Mediators Of Inflammation} Molecules that start, strengthen, limit, or coordinate inflammation. They include histamine, prostaglandins, cytokines, and other signals that change blood vessels, attract immune cells, cause symptoms, and support healing.

\medterm{Medical Cannabis And Cannabidiol} Cannabis and cannabidiol are substances from the cannabis plant used in some medical settings. Their effects and risks vary by product and dose; they may cause sleepiness, impaired judgment, interactions, or dependence and do not treat every condition.

\medterm{Medical Certificate} A document completed by an authorized healthcare professional that records a medical finding, diagnosis, fitness status, or period of illness.

\medterm{Medical Condition} A state affecting health or body function, whether temporary, long-lasting, symptomatic, or silent. A medical condition may result from disease, injury, inherited traits, infection, or normal biological changes.

\medterm{Medical Device} An instrument, implant, machine, software program, or other product used to diagnose, prevent, monitor, or treat disease or injury. Examples include pacemakers, syringes, imaging equipment, and blood-glucose meters.

\medterm{Medical Entomology} The study of insects and related arthropods that cause disease or transmit disease-causing organisms to humans, including their biology, ecology, physiology, and genetics.

\medterm{Medical Ethics} Moral principles that guide healthcare decisions, including respect for a person's choices, helping rather than harming, fairness, privacy, and honest communication. Decisions consider the patient's wishes, likely benefits and harms, professional duties, and relevant law.

\medterm{Medical Examiner} A physician who investigates sudden, unexpected, violent, or suspicious deaths. The examiner may order an autopsy, determine the medical cause of death, and provide information for the official death certificate; exact duties vary by location.

\medterm{Medical History} A medical history is a structured account of a person’s current concerns, past illnesses and operations, medications, allergies, family history, social history, and relevant exposures. It helps guide examination, diagnosis, and treatment.

\medterm{Medical Illness} The subjective experience of being unwell, which may or may not correspond to an
identifiable disease.

\medterm{Medical Jargon} Specialized words and phrases used by healthcare professionals. If they are not explained, patients may misunderstand their condition, test results, or treatment, so clinicians should use plain language whenever possible.

\medterm{Medical Jurisprudence} The branch of law (and its study) that deals with the rights and duties of physicians
and the legal aspects of medical practice, distinct from but overlapping with forensic
medicine.

\medterm{Medical Malpractice} Legal tort occurring when a healthcare provider fails to exercise the standard of
reasonable care, skill, or diligence expected of a prudent practitioner of similar
standing, resulting in injury, harm, or death to a patient.

\synonyms
Clinical Malpractice; Medical Malpractice

\medterm{Medical Marijuana} Cannabis used under medical supervision to relieve symptoms or support treatment of selected conditions. Products contain cannabinoids such as THC or cannabidiol; effects, interactions, impairment, and dependence risk vary by product and dose.

\medterm{Medical Marijuana} Alternate name; see \textit{Medical Cannabis}.

\medterm{Medical Nutrition Therapy} Personalized food and nutrition care provided by a qualified dietitian to manage a health problem. It may support people with diabetes, kidney disease, digestive disorders, cancer, obesity, or malnutrition.

\medterm{Medical Officer} A physician appointed to provide clinical care or medical oversight for an institution, service, or organization.

\medterm{Medical Physicist} A specialist who applies physics to medical imaging, radiation treatment, and radiation safety. Medical physicists check equipment, calculate radiation doses, and help ensure that scans and cancer treatments are accurate and as safe as possible.

\medterm{Medical Physics} The application of physics to medicine, including diagnostic imaging, radiation therapy, radiation safety, and physiological
measurement.

\medterm{Medical Pin} A medical pin is a slender metal fixation device inserted into bone or across a joint to stabilize a fracture or support orthopedic treatment. Pins may be temporary or permanent and require monitoring for infection and displacement.

\medterm{Medical Practitioner} A licensed physician qualified to diagnose, prevent, and treat illness.

\medterm{Medical Privacy} Protection of information that can identify a patient and control over its sharing. It includes keeping health records secure, limiting access to appropriate people, and obtaining permission when disclosure is not otherwise allowed.

\medterm{Medical Procedure} A planned clinical action or intervention performed by healthcare professionals to diagnose, measure, monitor, treat, or prevent a health condition, or to improve bodily function and structure. Medical procedures range in complexity from non-invasive evaluations (such as an electrocardiogram, ultrasound, or blood pressure check) to minimally invasive tests (such as drawing blood, inserting a catheter, performing an endoscopy, or taking a needle biopsy) and major invasive surgeries (such as tumor resections, joint replacements, or organ transplants). Procedures are classified by their primary clinical purpose: diagnostic (identifying the underlying cause of disease), therapeutic or curative (treating an illness, correcting a defect, or removing damaged tissue), palliative (relieving symptoms, obstruction, or pain to improve quality of life), or preventive/prophylactic (such as screening colonoscopies to remove precancerous polyps). Depending on the level of invasiveness and tissue involvement, procedures may be performed awake, with local numbing, under procedural sedation, or under general anesthesia, requiring strict sterile technique to reduce risks of bleeding, infection, and tissue injury.

\synonyms
Clinical Procedure; Healthcare Procedure; Medical Intervention; Surgical Procedure

\medterm{Medical Records} Written or electronic information about a person's health, including symptoms, examinations, test results, diagnoses, medicines, procedures, and follow-up. Accurate records help healthcare professionals coordinate safe care and track changes over time.

\medterm{Medical Repository} The first medical magazine and scientific journal published in the United States; founded in New York in 1797, it was published quarterly until 1824.

\medterm{Medical Symbol} A commonly recognized sign representing medicine, especially the Rod of Asclepius: a staff with one serpent. The winged, two-serpent caduceus is often used incorrectly as a medical symbol.

\medterm{Medical Tattooing} Medical tattooing (permanent makeup, micropigmentation) covers depigmented skin
(vitiligo, vitiligo camouflage), restores areolas after mastectomy, and masks scars.
Trichopigmentation simulates hair on bald scalp. Medical tattooing uses specialized
pigments and techniques.

\medterm{Medical Tests} Diagnostic procedures used to screen for, confirm, or monitor disease through analysis
of biologic specimens, imaging, physiologic measurements, or genetic material.

\medterm{Medicinal} Used for the prevention, diagnosis, relief, or treatment of disease.

\medterm{Medicine} A substance or treatment used to prevent, diagnose, relieve, or treat disease. Medicines may be prescription, over-the-counter, biological, or traditional, and all can have benefits, risks, and interactions.

\medterm{Medicines For Back Pain} Medicines for back pain may include selected anti-inflammatory drugs, acetaminophen, muscle relaxants, or other treatments depending on the cause and patient risks.

\medterm{Medicines For Osteoporosis} Osteoporosis medicines reduce bone loss or increase bone formation and may include bisphosphonates, denosumab, or anabolic treatments chosen according to fracture risk.

\medterm{Medicines For Sleep} Medicines for sleep may help short-term insomnia but can cause dependence, falls, confusion, or next-day impairment; treatment should also address the cause of poor sleep.

\medterm{Medicolegal} Pertaining to the application of medical knowledge to legal questions and the
intersection of medicine and the law. Covers forensic pathology, toxicology, psychiatry,
and clinical issues: informed consent, malpractice, medical records, disability,
personal injury, occupational disease, and public health law.

\medterm{Meditation} Meditation is a practice of intentionally directing attention, often toward the breath, a sensation, or an object, while observing thoughts without judgment. It may promote relaxation and stress management but does not replace indicated medical treatment.

\medterm{Mediterranean Anemia} Alternate name; see \textit{Thalassemia}.

\medterm{Mediterranean Diet} An eating pattern rich in vegetables, fruits, beans, whole grains, nuts, olive oil, and fish, with less red and processed meat and added sugar. It can support cardiovascular health as part of an overall balanced diet.

\medterm{Mediterranean Fever} An older name for brucellosis, a bacterial infection that can cause prolonged fever, sweats, fatigue, and joint or muscle pain.

\medterm{Medium} A substance or condition through which something occurs; in microbiology, a culture medium supports the growth of microorganisms.

\medterm{MEDLARS} Medical Literature Analysis and Retrieval System, an electronic bibliographic system developed by the U.S. National Library of Medicine to index and retrieve biomedical literature; it preceded MEDLINE.

\medterm{Medroxyprogesterone} Medroxyprogesterone is a progestogen used for contraception, abnormal uterine bleeding, endometriosis, and hormone therapy; depot injections may cause irregular bleeding, weight change, delayed fertility return, and reduced bone density.

\medterm{Medroxyprogesterone Acetate} Medroxyprogesterone acetate is a synthetic progestin used in hormonal contraception, treatment of endometriosis or abnormal uterine bleeding, and hormone therapy. It suppresses ovulation and produces progestational effects; formulation-specific risks and contraindications apply.

\medterm{Medulla} The inner or central part of an organ, such as the renal medulla or adrenal medulla.

\medterm{Medulla Oblongata} The lower part of the brainstem, joined to the spinal cord. It helps control breathing, heart rate, blood pressure, swallowing, coughing, vomiting, and other automatic functions.

\medterm{Medullary} Relating to the medulla or inner portion of an organ.

\medterm{Medullary Carcinoma} A cancer whose cells grow in broad groups and are surrounded by many immune cells. The term is used most often for cancers of the breast or thyroid, and its behavior depends on the organ and tumor findings.

\medterm{Medullary Sponge Kidney} A condition present from birth in which tiny urine-collecting tubes inside the kidneys are widened. Urine minerals may form stones or deposits, causing blood in the urine, repeated infections, pain, or no symptoms.

\medterm{Medullary Thyroid Carcinoma} A cancer of hormone-producing thyroid cells that make calcitonin. It may occur alone or as part of an inherited syndrome; blood tests, imaging, and a biopsy help confirm it.

\medterm{Medulloblastoma} A fast-growing cancer that usually begins in the cerebellum, the part of the brain involved in balance and coordination. It mainly affects children and may cause headache, vomiting, unsteady walking, or other neurologic problems.

\medterm{Medulloepithelioma} A medulloepithelioma is a rare embryonal tumor arising most often from the nonpigmented ciliary epithelium of the eye in children. It may cause visual symptoms, neovascular glaucoma, or a visible mass and is treated according to its extent and pathology.

\medterm{Mefenamic Acid} A nonsteroidal anti-inflammatory medicine used for short-term pain, including menstrual pain. It can irritate or bleed from the stomach, affect kidney function, and increase cardiovascular risk, especially with prolonged use or in people with risk factors.

\medterm{Mefloquine} Mefloquine is an antimalarial drug used for prevention and treatment of susceptible malaria. It has a long half-life but can cause neuropsychiatric, vestibular, and cardiac adverse effects, so patient-specific contraindications and warnings are important.

\medterm{Megabecquerel} A megabecquerel (MBq) is a unit of radioactivity equal to one million nuclear disintegrations per second. It is used to state the activity of radiopharmaceuticals and other radioactive materials.

\medterm{Megacalycosis} Megacalycosis is congenital, nonobstructive dilatation of the renal calyces, usually with otherwise preserved renal architecture and function. It is often incidental but must be distinguished from obstructive hydronephrosis.

\medterm{Megacolon} Abnormal enlargement and dilation of the colon, which may be congenital, acquired, toxic, or inflammatory.

\medterm{Megakaryocyte} A very large cell in the bone marrow that makes platelets, the small blood parts that help stop bleeding. It grows long extensions into nearby blood vessels, where pieces break off and enter the bloodstream as platelets.

\medterm{Megakaryocytopoiesis} Megakaryocytopoiesis is the formation and maturation of megakaryocytes in the bone marrow, followed by platelet production. It is regulated in large part by thrombopoietin and other hematopoietic signals.

\medterm{Megaloblast} An unusually large, immature cell in bone marrow that should develop into a red blood cell. It usually reflects faulty DNA production, commonly caused by vitamin B12 or folate deficiency.

\medterm{Megaloblastic Anaemia} Macrocytic anemia caused by defective DNA synthesis, most often from vitamin B12 or folate deficiency.

\medterm{Megaloblastic Anemia} Anemia caused by impaired DNA synthesis, usually related to vitamin B12 or folate deficiency, with abnormally large red blood cells.

\medterm{Megalocephaly} Abnormally large head size, which may be familial, developmental, metabolic, or associated with brain disease.

\medterm{Megamonas Funiformis} Megamonas funiformis is an anaerobic bacterium of the genus Megamonas, commonly associated with the intestinal microbiota. Its detection in a clinical specimen requires interpretation alongside symptoms, specimen quality, and other microbiologic findings.

\medterm{Megamonas Hypermegale} Megamonas hypermegale is an anaerobic bacterium in the genus Megamonas, primarily associated with the gastrointestinal microbiota. Its presence in a specimen is interpreted in relation to the site, clinical symptoms, and other microbiologic results.

\medterm{Megasphaera} Megasphaera is a genus of anaerobic bacteria found in the oral cavity and gastrointestinal and genital microbiota. Species identification may be relevant in microbiome studies or when recovered from a clinical specimen.

\medterm{Megaureter} An abnormally widened ureter, the tube that carries urine from the kidney to the bladder. It may result from a blockage, urine flowing backward, or abnormal ureter development and can cause infection, kidney swelling, or reduced kidney function.

\medterm{Megestrol} Megestrol is a synthetic progestin used in some settings to treat hormone-responsive cancers and, in certain formulations, anorexia or severe weight loss. It can cause glucocorticoid-like effects, thromboembolic risk, and adrenal suppression.

\medterm{Megestrol Acetate} Megestrol acetate is the acetate salt of the progestin megestrol. Depending on formulation and indication, it is used for advanced endometrial or breast cancer and for appetite stimulation in selected patients with severe weight loss.

\medterm{Meglitinides} Short-acting diabetes medicines that encourage the pancreas to release insulin around mealtimes. They lower blood sugar but can cause low blood sugar and weight gain, so doses are usually linked to meals and skipped when a meal is skipped.

\medterm{Megrim} An older term for migraine, a recurrent headache disorder often associated with nausea and sensitivity to light or sound.

\medterm{Meibomian Glands} Sebaceous glands in the eyelids that produce the oily layer of the tear film and help prevent tear evaporation.

\medterm{Meibomianitis} Meibomianitis is inflammation or blockage of oil-producing eyelid glands, causing irritated eyelid margins, dry eye, crusting, or recurrent styes.

\medterm{Meige Syndrome} A movement disorder causing repeated involuntary eyelid squeezing and muscle spasms around the mouth, jaw, or tongue. Symptoms may worsen with stress, speaking, chewing, or touch.

\medterm{Meigs Syndrome} A rare combination of a usually benign ovarian fibroma, fluid in the abdomen, and fluid around the lungs. The fluid often improves after the ovarian tumor is removed, but cancer and other causes must be excluded.

\medterm{Meigs' Syndrome} A rare combination of a usually benign ovarian fibroma, fluid in the abdomen, and fluid around the lungs. The fluid often improves after the ovarian tumor is removed, but cancer and other causes must be excluded.

\medterm{Meiosis} A specialized cell division that produces haploid eggs or sperm from a diploid precursor cell and creates
genetic variation through recombination and chromosome assortment.

\medterm{Meiotic Recombination} Meiotic recombination is exchange of DNA between homologous chromosomes during meiosis, usually through crossing over. It creates new combinations of alleles and contributes to genetic variation.

\medterm{Meissner's Corpuscles} Small touch sensors in hairless skin, especially the fingertips, palms, and soles. They help the nervous system detect gentle contact, texture, and movement across the skin.

\medterm{MEK Inhibitor} A targeted medicine that blocks MEK proteins in the MAPK signaling pathway, which can slow growth of some cancers with pathway changes. Side effects vary and may include rash, diarrhea, heart effects, and eye problems.

\medterm{Melancholia} An older term for severe depression, often describing pervasive sadness, loss of pleasure, psychomotor change, and biological symptoms.

\medterm{Melancholic Features} A pattern of severe depression marked by almost complete loss of pleasure, very low mood, major changes in movement or sleep, poor appetite or weight loss, excessive guilt, and symptoms that are often worse in the morning.

\medterm{Melanin} Melanin is a pigment produced by melanocytes that gives skin, hair, and eyes their colour and absorbs ultraviolet radiation; its amount and distribution influence pigmentation disorders and melanoma biology.

\medterm{Melanoblast} A precursor cell that develops into a melanocyte and produces pigment after migrating into the skin and other tissues.

\medterm{Melanocyte} A melanocyte is a pigment-producing cell derived from the neural crest. It makes melanin in specialized organelles called melanosomes and transfers the pigment to neighboring keratinocytes, helping protect skin from ultraviolet radiation.

\medterm{Melanocytes} Pigment-producing cells in the skin and hair follicles. They make melanin, the pigment that gives skin and hair much of their color and helps protect skin cells from ultraviolet radiation.

\medterm{Melanocytic Nevi} Common moles formed by collections of pigment-producing melanocytes. Most are harmless, but a mole that changes in size, shape, color, or symptoms should be assessed for possible melanoma.

\medterm{Melanocytic Nevus} A melanocytic nevus is a usually benign, localized proliferation of melanocytes that appears as a macule, papule, or nodule. Changes in size, color, border, or symptoms may require assessment for atypia or melanoma.

\medterm{Melanocytic Tumor Of The Central Nervous System} A tumor showing melanocytic differentiation within the CNS or meninges, ranging from benign melanocytoma to malignant melanoma. Evaluation includes imaging, histology, and molecular assessment to determine behavior.

\medterm{Melanoderma} Increased pigmentation or darkening of the skin caused by increased melanin production or an increased number of melanocytes.

\medterm{Melanoma} Melanoma is a malignant tumor of melanocytes, most often arising in skin but also in eye or mucosal sites; early detection matters because it can invade and metastasize.

\synonyms
Cutaneous Melanoma; Malignant Melanoma

\medterm{Melanoma Of The Eye} A cancer arising from pigment-producing cells in the eye, most often the uveal tract. It may cause blurred vision, flashes, a growing dark spot, or no symptoms at first and requires specialist eye assessment.

\medterm{Melanoma Surveillance} Ongoing examination of the skin and, when appropriate, the eyes or other previously affected sites after melanoma or in people at increased risk. It may include self-checks, clinician examination, dermoscopy, photographs, and evaluation of new or changing lesions.

\medterm{Melanoma Treatment} Treatment depends on the tumor’s depth, spread, molecular findings, overall health, and previous treatment. Localized melanoma is usually removed surgically; selected cases need lymph-node assessment or additional medicine. Disease that has spread may require immune-based or targeted treatment, surgery, radiation, or combinations.

\medterm{Melanophore} A melanophore is a pigment-containing cell, especially one in lower vertebrates, whose melanin granules can disperse or aggregate to alter visible coloration. In mammals, melanocytes are the principal melanin-producing cells.

\medterm{Melanosis} Abnormal accumulation or increased production of melanin in skin, mucosa, an organ, or other tissue.

\medterm{Melanosis Coli} Melanosis coli is dark pigmentation of the colon lining caused by lipofuscin accumulation, classically associated with prolonged stimulant-laxative use. It is usually benign and may fade after stopping the trigger, but the underlying constipation and medication use should be reviewed.

\medterm{Melanosome} A melanosome is a specialized organelle in a melanocyte in which melanin is synthesized, stored, and transported. Melanosomes are transferred to keratinocytes and influence skin, hair, and eye pigmentation.

\medterm{MELAS Syndrome} A mitochondrial disorder involving encephalopathy, lactic acidosis, and stroke-like episodes, often beginning before adulthood and caused by pathogenic mitochondrial DNA variants.

\medterm{Melasma} Melasma is acquired, usually symmetric brown or gray facial hyperpigmentation associated with ultraviolet exposure, hormones, pregnancy, and genetic susceptibility; it is benign but often recurrent.

\synonyms
Chloasma; Mask of Pregnancy

\medterm{Melasma} An acquired, chronic pigmentary disorder characterized by symmetric, hyperpigmented,
light-to-dark brown macules and patches on sun-exposed areas of the face (malar,
forehead, upper lip, chin). It primarily affects women of childbearing age and
individuals with darker skin types (Fitzpatrick types III-V).

\synonyms
Facial Melanosis; Mask of Pregnancy

\medterm{Melatonin} A hormone produced mainly by the pineal gland that helps coordinate sleep and wake timing. Darkness usually increases its release, while light suppresses it; supplements can cause drowsiness and medicine interactions.

\medterm{Melatonin Receptor} A cell-surface receptor activated by the hormone melatonin. The main types, MT1 and MT2, help regulate the body's sleep-wake clock and seasonal rhythms and are targets of some sleep medicines.

\medterm{Melatonin Receptor Agonists} Medicines that act like melatonin at receptors involved in the body's sleep-wake clock. They may help some people fall asleep, but choice of treatment depends on the cause of insomnia, other medicines, and health conditions.

\medterm{MELD} Model for End-Stage Liver Disease, a laboratory-based score used to estimate severity of advanced liver disease and help prioritize adults for liver transplantation. Traditional versions use bilirubin, creatinine, and INR; current implementations may include additional variables.

\medterm{MELD Score} A medical score estimating the severity and short-term risk of death from advanced liver disease. It uses laboratory results and helps guide urgency for liver transplantation; the calculation may vary as official versions are updated.

\medterm{Melena} Black, sticky, often foul-smelling stool caused by digested blood, usually from bleeding in the stomach or upper intestine. It can signal serious bleeding and requires prompt medical evaluation.

\medterm{Melibiose} A disaccharide composed of galactose and glucose, produced by hydrolysis of raffinose and metabolized by some microorganisms.

\medterm{Melio} An older shortened form related to melioidosis, the infection caused by Burkholderia pseudomallei.

\medterm{Melioidosis} An infectious disease caused by Burkholderia pseudomallei, usually acquired from
contaminated soil or water through inoculation or inhalation. It may cause pneumonia,
skin infection, abscesses, sepsis, or chronic disease and is more severe in people with
diabetes or other risk factors.

\medterm{Melitensis} A species designation, especially Brucella melitensis, a bacterium that causes brucellosis and is associated with goats and sheep.

\medterm{Melitten} Melitten, more commonly spelled melittin, is the principal peptide toxin in honeybee venom. It disrupts cell membranes and contributes to pain, inflammation, hemolysis, and other effects of envenomation.

\medterm{Melorheostosis} A rare sclerosing bone dysplasia producing irregular, flowing cortical hyperostosis on
imaging. It may cause pain, stiffness, limb deformity, restricted movement,
contractures, or soft-tissue fibrosis and can affect one or several contiguous bones.

\medterm{Meloxicam} Meloxicam is a nonsteroidal anti-inflammatory drug that inhibits cyclooxygenase enzymes and reduces prostaglandin synthesis. It is used for pain and inflammation, but can cause gastrointestinal, renal, cardiovascular, and hypersensitivity adverse effects.

\medterm{Melphalan} A chemotherapy medicine that damages cancer cells by interfering with their DNA. It is used for selected blood cancers and other cancers, but can also suppress bone marrow and increase infection or bleeding risk.

\medterm{Memantine} An NMDA receptor antagonist (uncompetitive, low-to-moderate affinity) used for moderate
to severe Alzheimer disease, often added to an acetylcholinesterase inhibitor. Blocks
excessive glutamate-mediated excitotoxicity while preserving normal synaptic
transmission. Once or twice daily oral.

\medterm{Memantine Hydrochloride} Memantine hydrochloride is an NMDA-receptor antagonist used to help manage moderate-to-severe Alzheimer disease. It reduces pathologic glutamate-mediated excitotoxicity; common adverse effects include dizziness, headache, constipation, and confusion.

\medterm{Membrane} A thin sheet of tissue or material that surrounds a cell, organ, or body cavity. Membranes can protect structures, separate spaces, allow selected substances to pass, and produce or absorb fluids.

\medterm{Membrane Channel} A membrane channel is a protein-forming pore across a cell membrane that permits selected ions or molecules to pass down an electrochemical or concentration gradient. Many channels are gated by voltage, ligands, or mechanical forces.

\medterm{Membrane Fusion} Membrane fusion is merging of two lipid bilayers into one continuous membrane. It is required for processes such as synaptic vesicle release, enveloped-virus entry, fertilization, and intracellular vesicle trafficking.

\medterm{Membrane Lipids} Fat-like molecules that form the main structure of cell membranes. They create a flexible barrier, help cells communicate, and provide attachment sites for proteins; abnormal lipid composition can affect signaling and cell function.

\medterm{Membrane Microdomain} A membrane microdomain is a small, specialized region of a cell membrane with a distinct composition of lipids and proteins. It can organize receptors, signaling molecules, transporters, and cytoskeletal interactions.

\medterm{Membrane Part} A membrane part is a structural or molecular component of a biological membrane, such as a lipid, protein, carbohydrate, channel, or receptor. These components determine membrane integrity, transport, signaling, and cell recognition.

\medterm{Membrane Potential} The electrical potential difference across a cell membrane, approximately -70 mV in most
neurons (resting membrane potential). Established by the unequal distribution of ions
(primarily K+, Na+, Cl-) across the membrane and selective permeability.

\medterm{Membrane Protein} A protein attached to or embedded in a cell membrane. It may transport substances, receive signals, act as an enzyme, join cells, anchor the cytoskeleton, or help the immune system identify cells.

\medterm{Membrane Raft} A membrane raft is a cholesterol- and sphingolipid-enriched membrane microdomain that helps cluster selected proteins and signaling molecules. Raft organization is dynamic and may change during receptor signaling or membrane trafficking.

\medterm{Membranoproliferative Glomerulonephritis} A pattern of kidney-filter damage that can cause blood and protein in the urine, swelling, and reduced kidney function. Causes include immune disorders, infections, and some blood cancers.

\medterm{Membranous} Relating to, consisting of, or resembling a membrane. In medicine, the term may describe a thin tissue layer or a disease pattern involving such a layer.

\medterm{Membranous Nephropathy} A kidney disease in which the filtering units are injured by immune proteins deposited along their delicate filter. It often causes large amounts of protein in the urine and swelling and may be linked to autoimmune disease, infection, medicines, or cancer.

\medterm{Memory} The ability to learn, store, and recall information. Memory depends on several connected areas of the brain and can be affected by sleep problems, stress, medicines, injury, illness, aging, or brain disorders.

\medterm{Memory Cell} A long-lasting B cell or T cell formed after an infection or vaccination. If the same germ is encountered again, memory cells respond faster and more strongly, helping prevent illness or reduce its severity.

\medterm{Memory Disorder} Impairment of encoding, storage, retrieval, or recognition of information that exceeds normal variation and may result from neurodegeneration, injury, vascular disease, medication, or other causes.

\medterm{Memory Impairment} Memory impairment is reduced ability to learn, store, retrieve, or recognize information. It may be transient or persistent and can result from normal aging, medications, sleep problems, neurologic disease, psychiatric illness, or other medical conditions.

\medterm{Memory Loss} Reduced ability to recall or learn information. It may result from normal aging, medical illness, medicines,
brain injury, or a neurocognitive disorder.

\medterm{Memory Loss} Reduced ability to recall past information, form new memories, or both. Causes include normal aging, medications, head injury, seizures, stroke, infection, substance use, neurocognitive disorders, and psychological conditions; sudden memory loss requires urgent assessment.

\medterm{Memory Span} The number of items, such as words or digits, that a person can briefly retain and recall, usually in the original order; it is a measure of short-term or working memory.

\medterm{Memory T Cells And B Cells} Long-lived immune cells that remember earlier exposures and respond faster when they meet the same target again. They support lasting protection after infection or vaccination, though protection may fade.

\synonyms
Memory Lymphocytes

\medterm{MEN 2} Abbreviation; see \textit{Multiple Endocrine Neoplasia Type 2}.

\medterm{Menarche} The first menstrual period, marking the onset of menstruation. Its timing varies with genetics, nutrition, health, and pubertal development; markedly early or delayed menarche may require assessment.

\medterm{Mendelian Inheritance} Inheritance of traits determined by single genes and transmitted according to dominant, recessive, or sex-linked patterns.

\medterm{Mendelism} The principles of inheritance based on Gregor Mendel's work, including segregation of alleles and independent assortment of genes; Mendelian traits may follow autosomal or X-linked dominant or recessive patterns.

\medterm{Mendelson Syndrome} A chemical pneumonitis caused by inhalation of acidic stomach contents, classically during anesthesia
or childbirth. It can cause cough, breathing difficulty, and impaired oxygenation.

\medterm{Meniere Disease} An inner-ear disorder causing repeated attacks of spinning dizziness, usually with fluctuating hearing loss, ringing in the ear, and a feeling of pressure or fullness. Attacks often last from 20 minutes to several hours.

\synonyms
Endolymphatic Hydrops; Meniere's Syndrome

\medterm{Meniere's Disease} Alternate spelling; see \textit{Meniere Disease}.

\medterm{Meniere's Disease} An inner-ear disorder causing episodes of vertigo, fluctuating hearing loss, tinnitus, and a feeling of fullness in one ear. Treatment aims to reduce attacks and protect hearing.

\medterm{Meningeal Carcinomatosis} Spread of cancer cells to the thin coverings and fluid around the brain and spinal cord. It may cause headache, nausea, weakness, seizures, vision changes, or problems with movement and needs specialist evaluation.

\medterm{Meningeal Leukemia} Meningeal leukemia is leukemia involving the membranes or fluid around the brain and spinal cord, potentially causing headache, weakness, or cranial-nerve problems.

\medterm{Meningeal Lymphoma} Meningeal lymphoma is lymphoma affecting the membranes or fluid spaces around the brain and spinal cord.

\medterm{Meningeal Neoplasm} A meningeal neoplasm is an abnormal growth arising in or near the membranes surrounding the brain and spinal cord.

\medterm{Meninges} Three protective layers surrounding the brain and spinal cord. Fluid between two of the layers cushions the nervous system; bleeding or infection in these layers can be dangerous.

\medterm{Meningioma} A usually slow-growing tumor arising from cells associated with the meninges. Most are benign, but some
are atypical or malignant and can cause symptoms by compressing the brain, spinal cord, or
cranial nerves.

\medterm{Meningism} Meningism is a syndrome of neck stiffness, headache, and sometimes photophobia or vomiting that resembles meningeal irritation without necessarily indicating meningitis. Serious causes, including infection or subarachnoid hemorrhage, require prompt evaluation.

\medterm{Meningismus} Neck stiffness and other signs of irritation around the brain and spinal cord without proven inflammation of those coverings. It can occur with fever, severe headache, dehydration, bleeding, or other illnesses and needs evaluation.

\medterm{Meningitis} Inflammation of the membranes around the brain and spinal cord, usually caused by infection but sometimes by medicines or immune disease. Fever, severe headache, stiff neck, light sensitivity, or confusion requires urgent assessment.

\medterm{Meningocele} A birth defect in which the protective coverings around the brain or spinal cord bulge through an opening in the skull or spine, forming a fluid-filled sac. It may cause nerve problems and needs specialist assessment.

\medterm{Meningocele Repair} Surgery to close an opening in the spine and return the protective coverings of the spinal cord to their proper position. It aims to prevent injury or infection; nerve, bladder, bowel, movement, or fluid problems may remain or develop.

\medterm{Meningococcal Disease} A serious infection caused by Neisseria meningitidis that can produce meningitis or rapidly spreading bloodstream infection. Fever, severe headache, stiff neck, confusion, vomiting, or a purple rash requires emergency care.

\medterm{Meningococcal Infection} An infection caused by the bacterium Neisseria meningitidis. It may remain mild or cause life-threatening meningitis or bloodstream infection, with fever, headache, stiff neck, vomiting, confusion, or a purple rash developing quickly.

\medterm{Meningococcal Meningitis} Acute bacterial meningitis caused by Neisseria meningitidis, often accompanied by bloodstream infection; it is a medical emergency requiring immediate antibiotics and public-health measures.

\medterm{Meningococcemia} Meningococcemia is bloodstream infection with Neisseria meningitidis, causing fever, rash, shock, clotting failure, and potentially fatal sepsis.

\medterm{Meningococcus} Neisseria meningitidis, a bacterium that can cause meningitis or rapidly spreading bloodstream infection. It spreads through respiratory secretions and may cause fever, headache, rash, shock, or death without urgent treatment.

\medterm{Meningoencephalitis} Inflammation of both the meninges surrounding the brain and the brain tissue itself, usually caused by
infection or autoimmune disease. Fever, headache, neck stiffness, confusion, seizures, or reduced
consciousness requires urgent medical assessment.

\medterm{Meningomyelocele} A birth defect in which the protective coverings and spinal-cord tissue bulge through an opening in the spine. It can cause weakness or paralysis of the legs and bladder or bowel problems.

\medterm{Meniscal Allograft Transplantation} Surgery placing a donated meniscus into a knee after substantial loss of the person's own meniscus. It may reduce pain in carefully selected younger patients, but healing is variable and infection, stiffness, graft failure, or later arthritis can occur.

\medterm{Meniscal Injuries} Tears or other damage to the knee menisci, which are crescent-shaped pieces of cartilage
that help distribute load and stabilize the knee. Symptoms may include joint-line pain,
swelling, clicking, or restricted motion. Diagnosis is based on examination and, when needed, imaging.

\medterm{Meniscal Tear} A split in one of the two rubbery cartilage pads that cushion the knee. It may follow a twisting injury or gradual wear and can cause pain along the joint line, swelling, clicking, catching, locking, or a feeling that the knee gives way.

\medterm{Meniscectomy} An operation that removes part or all of a damaged meniscus, the rubbery pad that cushions the knee. Surgeons usually remove only the torn portion when possible, because preserving meniscus helps protect the joint from later wear.

\medterm{Meniscus} A crescent-shaped fibrocartilage structure in the knee, or a curved fluid surface in a container.

\medterm{Meniscus Tear} Alternate name; see \textit{Torn Meniscus}.

\medterm{Menkes Disease} A rare inherited disorder in which the body cannot transport copper properly. Low copper in the brain and other tissues can cause poor growth, developmental decline, seizures, weak muscles, and sparse, brittle, or unusually twisted hair.

\medterm{Menkes Kinky Hair Syndrome} An X-linked disorder of copper transport caused by ATP7A dysfunction. It produces
progressive neurodegeneration, seizures, developmental delay, hypopigmented or brittle
kinky hair, connective-tissue abnormalities, and failure to thrive.

\medterm{Menogaril} Menogaril is a semisynthetic anthracycline derivative of nogalamycin. It intercalates into DNA and inhibits topoisomerase II, thereby impairing DNA replication and repair; it was investigated as an antineoplastic agent.

\medterm{Menometrorrhagia} Excessive uterine bleeding occurring both at expected menstrual times and at irregular intervals; the older term corresponds broadly to heavy and irregular menstrual bleeding and warrants evaluation for underlying causes.

\medterm{Menopausal} Relating to menopause or the years around the permanent end of menstrual periods. Hormone changes may cause hot flashes, sleep problems, mood changes, vaginal dryness, and increased bone-loss risk.

\medterm{Menopausal Hormone Therapy} Treatment with estrogen, with progestogen when needed to protect the uterus, for troublesome hot flashes, night sweats, or vaginal symptoms during menopause. It may affect clot, stroke, breast-cancer, and other risks, so treatment is individualized.

\medterm{Menopausal Hot Flashes} Alternate name; see \textit{Hot Flashes}.

\medterm{Menopause} The permanent end of menstrual periods because of loss of ovarian activity, usually diagnosed after 12 consecutive months without menstruation. Hormonal changes may cause hot flashes, night sweats, vaginal symptoms, sleep problems, or other symptoms.

\medterm{Menopause And Perimenopause} Perimenopause is the transition before menopause, when menstrual periods and hormone levels become irregular. Menopause is reached after 12 consecutive months without a period when no other cause explains it.

\medterm{Menopause Status} Menopause status describes whether a person is premenopausal, perimenopausal, or postmenopausal. Natural menopause is clinically defined after 12 consecutive months without menstruation when no other cause explains the amenorrhea.

\medterm{Menopause Syndrome} A nontechnical term for the group of symptoms that may occur during the menopausal transition, such as hot flashes, night sweats, sleep problems, mood changes, vaginal dryness, and menstrual irregularity. Menopause itself is a normal life stage, not a disease.

\medterm{Menorrhagia} Heavy or prolonged menstrual bleeding that interferes with daily life. Causes include fibroids, adenomyosis, polyps, bleeding disorders, thyroid disease, or irregular ovulation; ongoing blood loss may cause iron-deficiency anemia.

\medterm{Menotropin} Menotropin is a urinary-derived preparation containing follicle-stimulating and luteinizing hormone activity, used to stimulate follicular development or spermatogenesis; ovarian hyperstimulation and multiple pregnancy are important risks.

\medterm{Menses Duration} Menses duration is the number of days of menstrual bleeding in a cycle, counted from the first day of bleeding until bleeding stops. Persistently prolonged, heavy, or irregular bleeding warrants clinical assessment.

\medterm{Menstrual} Relating to menstruation or the menstrual cycle, including the monthly shedding of the uterine lining and the bleeding that follows. Cycle length and bleeding vary, but very heavy, painful, irregular, or unexpected bleeding may need evaluation.

\medterm{Menstrual Cramps} Cramping pain in the lower abdomen during or around menstruation, medically called dysmenorrhea. Pain may radiate to
the back or thighs and may occur with nausea, diarrhea, headache, or fatigue. Severe, worsening, or newly
developed pain requires evaluation for causes such as endometriosis or fibroids.

\synonyms
Dysmenorrhea; Period Pain; Painful Menstruation

\medterm{Menstrual Cycle} The recurring hormonal and uterine changes that prepare the body for possible pregnancy. Cycle length and
bleeding pattern vary among individuals and across life stages.

\medterm{Menstrual Disorders} Conditions involving abnormal menstrual timing, regularity, duration, amount of bleeding, or pain. They include absent or infrequent periods, irregular cycles, unusually heavy or prolonged bleeding, painful periods, and bleeding between periods or after sex. Causes include pregnancy, hormonal or ovulatory disorders, fibroids, endometriosis, thyroid disease, medicines, structural abnormalities, and perimenopause; persistent, heavy, or unexpected bleeding requires assessment.

\medterm{Menstruation} The periodic shedding of the endometrial lining as part of the menstrual cycle when pregnancy does not occur.
The timing, duration, and amount of bleeding vary with age, hormonal state, medicines,
and health conditions.

\medterm{Menstruation Disturbance} Alternate name; see \textit{Menstrual Disorders}.

\medterm{Mental} Relating to thinking, memory, emotion, behavior, or other functions of the mind. It does not by itself mean that a condition is psychiatric.

\medterm{Mental Age} An estimate of cognitive functioning expressed as the age level at which a standardized test score is typical; it is not the same as chronological age.

\medterm{Mental Disorder} Alternate singular form; see \textit{Mental Disorders}.

\medterm{Mental Disorders} Health conditions involving a clinically significant disturbance in cognition, emotion regulation, or behavior that causes distress, impairment, or both. Examples include mood, anxiety, psychotic, neurodevelopmental, trauma-related, personality, and substance use disorders. A mental-health concern or symptom does not necessarily constitute a mental disorder.

\medterm{Mental Disorders In Children} Mental disorders in children involve persistent, developmentally atypical disturbances in mood, thinking, behavior, or functioning that cause clinically significant distress or impairment. Symptoms must be interpreted in relation to the child’s developmental stage and distinguished from normal variation, temporary reactions, and age-appropriate behavior.

\medterm{Mental Fatigue} Mental fatigue is a subjective feeling of reduced mental energy, concentration, or cognitive endurance after sustained effort or illness. Persistent or severe symptoms may accompany sleep disorders, depression, neurologic disease, medication effects, or other medical conditions.

\medterm{Mental Health} Mental health encompasses emotional, psychological, and social well-being, affecting how individuals think, feel, behave, and handle stress, relationships, and decision-making.

\medterm{Mental Health Crisis} A period of severe emotional, behavioral, or psychological distress in which a person’s usual ability to cope or maintain safety is overwhelmed. It may involve suicidal thoughts, severe agitation, psychosis, substance intoxication, or inability to care for basic needs; it may or may not meet criteria for a psychiatric emergency.

\medterm{Mental Illness} A broad, commonly used term for mental disorders and related conditions affecting mood, thinking, behavior, perception, or daily functioning. See \textit{Mental Disorders}.

\medterm{Mental Illness In Children} Broad, commonly used term; see \textit{Mental Disorders In Children}.

\medterm{Mental Process} A mental process is a cognitive or affective activity of the mind, such as perception, attention, memory, reasoning, language, emotion, or decision-making.

\medterm{Mental Recall} Mental recall is retrieval of previously learned or stored information from memory. It may be assessed by asking a person to remember facts, events, words, or experiences after a delay.

\medterm{Mental Status Examination} A structured part of a clinical interview that assesses a person's appearance, behavior, speech, mood, thinking, perceptions, attention, memory, awareness, judgment, and understanding of their condition.

\medterm{Mental Status Testing} A structured assessment of appearance, behavior, alertness, orientation, attention, memory, language, thought, perception, mood, insight, and judgment.

\medterm{Mentation} The process of thinking, perceiving, reasoning, and other mental activity.

\medterm{Menthol} A mint-derived or synthetic compound producing a cooling sensation and used in some topical and respiratory preparations.

\medterm{Menthol Poisoning} Poisoning after swallowing, inhaling, or having substantial skin exposure to menthol. It may cause nausea, abdominal pain, dizziness, drowsiness, agitation, seizures, coma, or breathing problems; suspected significant exposure needs urgent poison-control advice.

\medterm{Mepacrine} An older antimalarial and antiprotozoal medicine, also called quinacrine, now used only in selected indications.

\medterm{Meperidine} An opioid analgesic, also called pethidine, whose use is limited by adverse effects and an active neurotoxic metabolite.

\medterm{Meperidine Hydrochloride} Meperidine hydrochloride is the hydrochloride salt of the opioid analgesic meperidine. It produces analgesia through opioid-receptor activity and carries risks of sedation, respiratory depression, seizures, dependence, and serotonin toxicity with interacting drugs.

\medterm{Meperidine Hydrochloride Overdose} Taking too much meperidine hydrochloride, an opioid, can cause extreme sleepiness, pinpoint pupils, slow or stopped breathing, low blood pressure, seizures, or coma. Delayed toxicity and serotonin toxicity are possible, so emergency assessment is required.

\medterm{Mephenesin} An older centrally acting muscle relaxant with sedative effects. It is rarely used now and can impair alertness, coordination, and breathing when combined with other sedatives.

\medterm{Mephentermine} A stimulant medicine formerly used in selected situations to raise dangerously low blood pressure. It increases activity of certain adrenaline-like receptors and can cause fast heartbeat, anxiety, or abnormal heart rhythms.

\medterm{Mephenytoin} Mephenytoin is a hydantoin anticonvulsant formerly used to treat seizures. It is metabolized partly to ethotoin and has limited modern use because of potentially serious adverse effects, including hypersensitivity and blood dyscrasias.

\medterm{Mepivacaine} A local anesthetic medicine that temporarily blocks pain signals from nerves. It is injected for numbing skin, dental areas, individual nerves, or larger body regions during procedures.

\medterm{Meprobamate} An older medicine that reduces anxiety and causes sleepiness. It can cause dependence, dangerous breathing and nervous-system depression, and overdose, especially with alcohol or other sedatives, so it is rarely used today.

\medterm{Meprobamate Overdose} Taking too much meprobamate can cause severe drowsiness, confusion, poor coordination, low blood pressure, slowed breathing, coma, or death. Alcohol and other sedatives increase the danger; suspected overdose requires emergency evaluation.

\medterm{Mepyramine} An antihistamine that blocks H1 receptors and may cause sedation; it is also called pyrilamine.

\medterm{Meralgia Paresthetica} Burning, tingling, or numbness over the outer thigh caused by compression of the lateral femoral cutaneous nerve.

\synonyms
Burning Thigh Pain; Bernhardt-Roth Syndrome

\medterm{Merbromin Poisoning} Poisoning from merbromin, an older mercury-containing antiseptic. Swallowing or repeated exposure may injure the mouth, stomach, kidneys, or nervous system; the amount and route determine risk and require poison-control or medical assessment.

\medterm{Mercaptopurine} A thiopurine antimetabolite used in selected leukemias and inflammatory disorders, requiring monitoring for marrow and liver toxicity.

\medterm{Mercurialism} Long-term poisoning caused by mercury exposure. It may damage the nervous system, kidneys, digestive tract, mouth, or mental health, with effects depending on the mercury form, dose, and exposure duration.

\medterm{Mercuric Chloride Poisoning} Severe poisoning from mercuric chloride, a corrosive inorganic mercury salt. It can burn the mouth and digestive tract and cause vomiting, bloody diarrhea, shock, kidney failure, or neurologic injury; emergency treatment is required.

\medterm{Mercuric Oxide Poisoning} Poisoning from mercuric oxide, an inorganic mercury compound that can irritate or burn tissues and damage the kidneys and nervous system. Swallowing or inhalation requires urgent poison-control or emergency assessment.

\medterm{Mercury} A metal used in some industrial and laboratory products. Exposure can harm the nervous system, kidneys, lungs, or digestive tract, so it must be handled and disposed of safely.

\medterm{Mercury Poisoning} Illness caused by exposure to elemental, inorganic, or organic mercury. Effects depend on the form and dose and may include nerve, kidney, lung, or digestive injury; urgent advice is needed after significant exposure.

\medterm{Meridian} A reference line or plane; in traditional acupuncture it denotes a proposed pathway for qi, while in anatomy it may describe a longitudinal line.

\medterm{Merkel Cell Carcinoma} An aggressive neuroendocrine skin cancer arising from Merkel cells. It presents as a
rapidly growing, painless, red-violet nodule on sun-exposed skin of elderly patients.
Approximately 80 percent are associated with Merkel cell polyomavirus. Risk factors
include immunosuppression, UV exposure, and fair skin.

\synonyms
MCC; Neuroendocrine Carcinoma of the Skin

\medterm{Merkel Cells} Specialized touch-sensing cells in the deepest layer of the skin that help detect light touch and sustained pressure.

\medterm{Mermaid Syndrome} An extremely rare congenital anomaly in which the legs are fused to varying degrees, also called sirenomelia.

\medterm{Meropenem} Meropenem is a broad-spectrum carbapenem antibiotic that inhibits bacterial cell-wall synthesis and treats severe susceptible infections, including meningitis; allergy, gastrointestinal effects, seizures, and resistance are possible.

\medterm{Meropenem-Vaborbactam} An intravenous antibiotic combination used for some serious infections caused by bacteria resistant to many other antibiotics. The choice depends on laboratory susceptibility testing, infection site, kidney function, and local guidance.

\medterm{MERRF Syndrome} A rare inherited disorder affecting how cells produce energy. It commonly causes muscle jerks, seizures, balance problems, and muscle weakness, and may also affect hearing, vision, or other organs.

\medterm{MERS-CoV} Middle East respiratory syndrome coronavirus, the virus that causes MERS. Infection can cause fever and respiratory illness ranging from mild disease to severe pneumonia, especially in people with other health conditions.

\medterm{Merthiolate Poisoning} Poisoning from merthiolate, commonly referring to the mercury-containing antiseptic thimerosal. Small skin exposures usually cause local irritation, but swallowing or repeated exposure may cause corrosive, kidney, or neurologic injury and needs advice.

\medterm{Merycism} Rumination, the repeated regurgitation and rechewing of food after eating.

\medterm{Mesalamine} Mesalamine, or 5-aminosalicylic acid, is an anti-inflammatory drug used to induce and maintain remission in ulcerative colitis and some other inflammatory bowel conditions. It acts mainly topically in the intestinal mucosa and can rarely affect renal function.

\medterm{Mesangiolysis} Mesangiolysis is dissolution or disruption of the mesangial tissue within a renal glomerulus. It may occur in severe endothelial injury and can be associated with microvascular lesions and glomerular capillary aneurysmal changes.

\medterm{Mesangium} The mesangium is the supportive tissue between glomerular capillary loops, composed of mesangial cells and extracellular matrix. Mesangial cells provide structural support, regulate capillary flow, and clear certain substances; abnormal expansion occurs in several glomerular diseases.

\medterm{Mescaline} A naturally occurring psychedelic substance found in certain cacti that alters perception and may cause hallucinations and autonomic effects.

\medterm{Mesencephalon} The mesencephalon, or midbrain, is the part of the brainstem between the pons and diencephalon. It contains pathways and nuclei involved in eye movements, motor control, arousal, and auditory and visual reflexes.

\medterm{Mesenchymal Cell} A connective-tissue precursor that can develop into bone, cartilage, fat, muscle, or tendon cells.

\medterm{Mesenchymal Chondrosarcoma} Mesenchymal chondrosarcoma is a rare aggressive cancer containing cartilage-forming tissue, often arising in bone or soft tissue.

\medterm{Mesenchymal Tumor Of The Central Nervous System} A rare tumor arising from connective or supporting tissue in or around the brain or spinal cord rather than from nerve cells. Its behavior is determined by examination of tissue, imaging, and sometimes genetic testing.

\medterm{Mesenchyme} Embryonic connective tissue that gives rise to much of the skeleton, connective tissue, blood vessels, and smooth muscle.

\medterm{Mesenchymoma} A mesenchymoma is a rare tumor containing more than one type of mesenchymal tissue, sometimes including cartilage, bone, fat, or vascular elements. Behavior depends on the site, age, histology, and presence of malignant components.

\medterm{Mesenteric} Relating to the mesentery, a fold of tissue supporting the intestines. It contains blood vessels, nerves, lymph vessels, and fat that help supply, drain, anchor, and protect the intestines.

\medterm{Mesenteric Angiography} Contrast imaging of the arteries supplying the intestines, used to evaluate mesenteric ischemia, bleeding, aneurysm, stenosis, occlusion, or vascular anatomy.

\medterm{Mesenteric Artery} A mesenteric artery supplies blood to portions of the intestine and associated abdominal organs. The superior and inferior mesenteric arteries are the major vessels; obstruction or embolism can cause acute mesenteric ischemia.

\medterm{Mesenteric Artery Ischemia} Inadequate blood flow through the mesenteric arteries, usually caused by embolism,
thrombosis, or severe low-flow states. Acute disease presents with severe abdominal pain
out of proportion to examination and can progress to bowel infarction.

\medterm{Mesenteric Cyst} A rare, benign, fluid-filled cystic mass that develops within the folds of the abdominal mesentery, most frequently originating from sequestered embryonic lymphatic tissue in the small bowel mesentery (especially the ileum). On physical examination, it demonstrates the classic Tillaux's sign: mobility in a transverse direction perpendicular to the mesenteric root, with limited or no movement vertically. While often asymptomatic and found incidentally, large cysts can cause abdominal distension, chronic pain, or acute surgical emergencies such as cyst torsion, internal hemorrhage, infection, rupture, or mechanical bowel obstruction. Diagnosis is made with abdominal ultrasound, CT, or MRI, and definitive management is complete surgical enucleation or laparoscopic excision to prevent recurrence.

\synonyms
Mesenteric Lymphangioma; Chylolymphatic Cyst; Enteric Mesenteric Cyst

\medterm{Mesenteric Ischemia} Reduced blood flow to the intestines, caused suddenly by a clot or gradually by narrowed arteries. It can cause severe abdominal pain, bowel injury, infection, or death and requires urgent assessment.

\medterm{Mesenteric Lymphadenitis} Inflammation of lymph nodes in the tissue supporting the intestines, often after a viral or bacterial infection. It can cause right-sided abdominal pain, fever, nausea, and tenderness and may resemble appendicitis.
\synonyms
Mesenteric Adenitis

\medterm{Mesenteric Veins} Mesenteric veins drain blood from the small and large intestines and carry it mainly through the superior and inferior mesenteric veins to the portal venous system. Thrombosis or impaired flow can cause intestinal congestion and ischemia.

\medterm{Mesenteric Venous Thrombosis} A blood clot in a mesenteric vein that impairs intestinal venous drainage and can cause abdominal pain, bowel edema, ischemia, infarction, or gastrointestinal bleeding.

\medterm{Mesentery} A fold of the abdominal lining that attaches the intestines to the back wall of the abdomen. It carries blood vessels, nerves, and lymph vessels to the bowel and helps hold the intestines in place.

\medterm{Mesh} A sheet-like implant made from synthetic or biological material that reinforces weakened tissue. It may be used in selected hernia or pelvic-floor operations; possible problems include pain, infection, scarring, movement of the implant, or injury to nearby organs.

\medterm{Mesial} Located toward the body's midline; in dentistry, located toward the front or midline of a tooth. The term helps describe tooth surfaces, positions, cavities, gum disease, and dental restorations.

\medterm{Mesial Surface} The mesial surface is the surface of a tooth or other structure oriented toward the midline. In dentistry, mesial surfaces contact adjacent teeth and can be involved in caries, periodontal disease, or restoration margins.

\medterm{Mesna} A protective medicine given with certain chemotherapy drugs to reduce irritation and bleeding of the bladder. It binds harmful breakdown products in the urine and does not treat the cancer itself.

\medterm{Mesoderm} The middle germ layer in an embryo, which develops into muscles, bones, cartilage, blood
vessels, and the kidneys.

\medterm{Mesonephric} Mesonephric means relating to the mesonephros, a temporary embryonic kidney, or to structures derived from its ducts and tubules. Mesonephric remnants can give rise to uncommon benign or malignant lesions of the genital tract.

\medterm{Mesonephric Adenocarcinoma} A rare HPV-independent cervical adenocarcinoma arising from mesonephric remnants, usually presenting as a cervical mass or abnormal bleeding and requiring histopathologic diagnosis.

\medterm{Mesonephric Tumor} A mesonephric tumor is a rare tumor arising from remnants of embryonic mesonephric ducts, usually in the female reproductive tract.

\medterm{Mesothelial Neoplasm} A mesothelial neoplasm is a growth arising from the lining of body cavities, such as the pleura, peritoneum, or pericardium.

\medterm{Mesothelin} Mesothelin is a cell-surface glycoprotein expressed mainly by mesothelial cells and overexpressed in several cancers, including mesothelioma and some ovarian and pancreatic cancers. It is being studied as a biomarker and therapeutic target.

\medterm{Mesothelioma} A cancer arising from the thin lining of body cavities, most often the pleura around the lungs. It is strongly associated with asbestos exposure and may cause chest pain, breathlessness, or fluid around the lung.

\medterm{Mesothelium} The mesothelium is a single layer of specialized cells lining serous body cavities and covering their organs, including the pleura, peritoneum, and pericardium. It produces lubricating fluid and participates in inflammation and repair.

\medterm{Messenger RNA} A temporary copy of genetic instructions made from DNA. It carries those instructions to cell structures that build proteins and is then broken down; some vaccines use messenger RNA to teach the immune system to recognize a target.

\medterm{Met} Abbreviation for methionine, an essential sulfur-containing amino acid; it is also represented by the one-letter code M and the codons AUG.

\medterm{Metabolic} Relating to metabolism, the chemical processes by which cells produce and use energy and materials.

\medterm{Metabolic Acidosis} A buildup of acid or loss of bicarbonate that makes the blood too acidic. Causes include severe diabetes, poor circulation, kidney failure, starvation, diarrhea, or poisoning and require evaluation of the underlying cause.

\medterm{Metabolic Alkalosis} A condition in which the blood becomes too alkaline, often because of repeated vomiting, fluid loss, or certain medicines. It can cause weakness, cramps, confusion, or abnormal heart rhythms.

\medterm{Metabolic Bone Disease} A group of disorders in which abnormal handling of calcium, phosphate, vitamin D, parathyroid hormone, or other factors weakens or changes bone. Examples include osteomalacia, rickets, and some disorders caused by kidney or endocrine disease.

\medterm{Metabolic Disease} A disorder that changes how the body processes or uses substances such as sugars, fats, proteins, minerals, or hormones. It may be inherited or acquired and can affect energy, growth, blood chemistry, or organ function.

\medterm{Metabolic Disorders} Alternate plural form; see \textit{Metabolic Disease}.

\medterm{Metabolic Dysfunction-Associated Steatohepatitis} Alternate name; see \textit{Fatty Liver Disease}.

\medterm{Metabolic Dysfunction-Associated Steatotic Liver Disease} Alternate name; see \textit{Fatty Liver Disease}.

\medterm{Metabolic Equivalents} Units that compare the energy used during an activity with the energy used while resting. They provide a rough estimate of exercise intensity, but actual energy use varies between people.

\synonyms
METs

\medterm{Metabolic Myopathies} Muscle disorders caused by impaired energy metabolism, often producing exercise intolerance, cramps, weakness, muscle pain, or episodes of rhabdomyolysis.

\medterm{Metabolic Neuropathies} Nerve damage caused by a metabolic or systemic disorder, such as diabetes, kidney failure, vitamin deficiency, or an inherited metabolic disease. Symptoms commonly include numbness, burning pain, tingling, weakness, or problems with balance.

\medterm{Metabolic Rate} The speed at which the body uses energy to keep organs working and support activity. It changes with body size, age, muscle amount, hormones, illness, temperature, food intake, and exercise.

\medterm{Metabolic Response} A metabolic response is a change in energy use, substrate metabolism, or biochemical pathways in response to conditions such as fasting, exercise, infection, trauma, or medication. The response may be adaptive or contribute to illness when excessive or prolonged.

\medterm{Metabolic Syndrome} A combination of increased waist size, high blood pressure, high blood sugar, high triglycerides, or low HDL cholesterol. Together these findings increase the risk of heart disease, stroke, and type 2 diabetes.

\synonyms
Insulin Resistance Syndrome; Syndrome X

\medterm{Metabolic Tumor Volume} The estimated amount of tumor tissue that shows increased activity on a PET scan. It can help measure disease burden and treatment response, especially in lymphoma, but inflammation and scan technique can affect the result.

\medterm{Metabolism} The chemical processes cells use to turn food and medicines into energy and building materials, and to break down substances the body no longer needs.

\medterm{Metabolite} A substance made, changed, or used when the body processes a medicine, nutrient, or naturally occurring molecule. Measuring metabolites in blood or urine can show body function, disease, or drug effects.

\medterm{Metacarpal} Relating to one of the five bones forming the palm between the wrist and fingers. Metacarpal fractures can cause pain, swelling, deformity, reduced grip, and difficulty moving the hand.

\medterm{Metacarpal Bone} A metacarpal bone is one of the five long bones of the hand between the wrist bones and the proximal phalanges. Each has a base, shaft, and head; fractures can affect alignment, grip, and joint function.

\medterm{Metacarpals} The five bones forming the palm, numbered from the thumb to the little finger. Each connects the wrist to a finger and helps with grip and hand movement; breaks can cause pain, swelling, deformity, or reduced function.

\medterm{Metacarpophalangeal Joint} A knuckle joint connecting a palm bone to a finger or thumb bone. It allows the digit to bend, straighten, and move from side to side and can be injured by impact or twisting.

\medterm{Metacarpus} The part of the hand between the wrist and fingers, containing the five metacarpal bones.

\medterm{Metachromatic Leukodystrophy} A rare inherited disorder in which harmful fatty substances build up in nerve cells and damage the protective covering of nerves. It causes worsening movement, thinking, speech, or vision problems, often beginning in childhood.

\medterm{Metagonimus Yokogawai} Metagonimus yokogawai is a small intestinal fluke acquired by eating raw or inadequately cooked freshwater fish. Infection can cause abdominal symptoms and, rarely, ectopic migration of eggs or adult worms.

\medterm{Metal Artifact} Distortion or loss of detail on an X-ray, CT, or MRI image caused by metal such as an implant, clip, or bullet. It can hide nearby findings, so special imaging methods or another test may be needed.

\medterm{Metal Cleaner Poisoning} Poisoning from a metal-cleaning product containing acids, alkalis, solvents, or other chemicals. It may burn the skin, eyes, mouth, or digestive tract, or cause choking and lung injury if aspirated; do not induce vomiting.

\medterm{Metal Fume Fever} A short-lived influenza-like illness caused by inhalation of fumes from heated metals, especially zinc; symptoms include fever, chills, headache, muscle aches, and cough.

\medterm{Metal Polish Poisoning} Poisoning from metal polish, which may contain petroleum solvents, acids, or other chemicals. Swallowing or inhalation can cause vomiting, drowsiness, chemical burns, or aspiration-related lung injury; urgent product-specific advice is needed.

\medterm{Metalloendopeptidase} A metalloendopeptidase is an endopeptidase whose catalytic activity depends on a bound metal ion, commonly zinc. It cleaves peptide bonds within proteins and includes enzymes involved in extracellular matrix remodeling and signaling.

\medterm{Metalloporphyrins} Metalloporphyrins are porphyrin molecules containing a centrally coordinated metal ion. Heme, which contains iron, is a biologically important example and forms the functional group of hemoglobin, myoglobin, and many enzymes.

\medterm{Metalloprotease} A metalloprotease is a proteolytic enzyme that requires a metal ion, commonly zinc, for catalytic activity. Metalloproteases regulate extracellular matrix turnover, cell signaling, inflammation, and tissue remodeling; dysregulation contributes to disease.

\medterm{Metalloprotein} A protein that contains a metal ion or binds one tightly as part of its structure or function. Iron, zinc, copper, and magnesium in metalloproteins support oxygen transport, enzyme activity, electron transfer, or regulation.

\medterm{Metallothionein} Metallothionein is a small, cysteine-rich protein that binds metals such as zinc, copper, and cadmium. It helps regulate essential metals and may protect cells from metal toxicity and oxidative stress.

\medterm{Metanephrine} Metanephrine is a metabolite of epinephrine formed mainly by catechol-O-methyltransferase. Plasma or urinary metanephrine measurements are used in evaluation for catecholamine-secreting tumors such as pheochromocytoma and paraganglioma.

\medterm{Metaphase} The stage of cell division in which chromosomes align at the cell's equatorial plane.

\medterm{Metaphase Plate} The metaphase plate is the plane at the center of a dividing cell where chromosomes align during metaphase. It is a geometric arrangement rather than a discrete anatomical structure.

\medterm{Metaphyseal Corner Fractures} Small corner or bucket-handle breaks near the growth plates of an infant’s long bones. They are highly concerning for forceful injury, especially when multiple or at different healing stages, but must be interpreted with the full medical and imaging assessment.

\medterm{Metaphyseal Widening} Metaphyseal widening is enlargement of the growth-plate or metaphyseal region of a developing long bone. It may be seen in disorders such as rickets, skeletal dysplasia, infection, or healing injury and must be interpreted with imaging and clinical findings.

\medterm{Metaphysis} The wider part of a growing long bone between its shaft and end. In children it lies next to the growth plate, where new bone forms and the bone lengthens.

\medterm{Metaplasia} A usually reversible change in which one mature cell type is replaced by another that better tolerates ongoing irritation or stress. Some persistent forms may increase the risk of further abnormal changes.

\medterm{Metaplastic Carcinoma} A cancer whose cells have changed to resemble a different type of tissue. The term is used especially for an uncommon, often aggressive breast cancer; diagnosis requires examination of a tissue sample.

\medterm{Metapneumovirus} A respiratory virus that can cause cough, fever, bronchiolitis, or pneumonia, especially in infants, older adults, and people with chronic disease.

\medterm{Metastasectomy} Surgery to remove one or more cancer deposits that have spread from the original tumor. It may help selected people when the deposits are limited and removable, but benefit depends on the cancer type, location, and other treatment options.

\medterm{Metastases To The Breast} Secondary malignant deposits in breast tissue from a nonmammary primary tumor, most often melanoma, lymphoma, or carcinoma from another organ. Imaging and biopsy are needed because the treatment differs from primary breast carcinoma.

\medterm{Metastasis} The spread of cancer cells from the original tumor to a separate part of the body, where they form new tumors. Spread can occur through the blood, lymphatic system, or body cavities.

\medterm{Metastasis Of Colorectal Cancer} Spread of cancer from the colon or rectum to another part of the body, most often the liver, lungs, or abdominal lining. Symptoms and treatment depend on the sites involved, number of tumors, and tumor test results.

\medterm{Metastasis-Free Survival} The length of time after treatment before cancer spreads to a distant part of the body or the person dies. Researchers use it to compare treatments and estimate how well they delay spread.

\medterm{Metastatic Adenocarcinoma} A gland-forming cancer that has spread from where it began to distant organs such as the liver, lungs, bones, or brain. Identifying the original site is important because treatment depends on its origin and spread.

\medterm{Metastatic Brain Tumor} A metastatic brain tumor is cancer that has spread to the brain from a primary tumor elsewhere in the body.

\medterm{Metastatic Breast Cancer} Breast cancer that has spread from the breast to distant organs or tissues, such as bone, liver, lung, or brain. It is treatable, but usually considered advanced and not curable.

\medterm{Metastatic Calcification} Buildup of calcium salts in otherwise healthy tissues because the blood calcium level is too high. It can affect the kidneys, lungs, stomach, or blood vessels and may damage their function.

\medterm{Metastatic Cancer} Cancer that has spread from its original site to distant organs or tissues. The spread may occur through the blood or lymphatic system and can form new tumors. Metastatic cancer is often called stage 4 cancer, although staging depends on the cancer type. The new tumors retain the characteristics of the original cancer; for example, breast cancer that spreads to the lung remains metastatic breast cancer, not lung cancer.

\synonyms
Stage 4 Cancer; Advanced Metastatic Cancer; Secondary Cancer

\medterm{Metastatic Castration-Resistant Prostate Cancer} Prostate cancer that has spread and continues growing despite treatment that lowers male hormones. Care may include hormone-targeting medicines, chemotherapy, radiation, immune treatment, or medicines for bone protection.

\medterm{Metastatic Colon Cancer} Colon cancer that has spread beyond the bowel, most often to the liver, lungs, or abdominal lining. Symptoms and treatment depend on where it has spread, the tumor’s genetic findings, and the person’s overall health.

\medterm{Metastatic Colorectal Cancer} Alternate name; see \textit{Colorectal Cancer}.

\medterm{Metastatic Inguinal Lymphadenopathy} Enlargement of lymph nodes in the groin because cancer cells have spread to them, often from the leg, genital area, lower abdomen, or anus. Examination and imaging help identify the original cancer and its extent.

\medterm{Metastatic Lung Cancer} Lung cancer that has spread beyond the lungs to organs such as the brain, bones, liver, or adrenal glands. Treatment depends on the cancer type, molecular findings, areas affected, and overall health.

\medterm{Metastatic Pleural Tumor} A cancer deposit in the thin lining around the lung that has spread from another tumor, often from the lung, breast, or digestive tract. It may cause chest pain, breathlessness, or fluid around the lung.

\medterm{Metastin} Metastin is an older name for kisspeptin, a peptide derived from the KISS1 gene. It stimulates hypothalamic gonadotropin-releasing hormone secretion and helps regulate puberty, reproductive function, and tumor-cell signaling.

\medterm{Metatarsal} Relating to one of the five long bones in the middle and front of the foot. Metatarsal bones connect the ankle bones to the toe bones and may be injured by falls, twisting, or repeated stress.

\medterm{Metatarsal Bone} A metatarsal bone is one of the five long bones of the forefoot between the tarsal bones and proximal phalanges. Metatarsal fractures and alignment abnormalities can affect weight bearing and gait.

\medterm{Metatarsal Valgus} Metatarsal valgus is outward deviation of one or more metatarsal bones relative to the foot’s normal axis. It may contribute to forefoot deformity and altered load distribution and is assessed clinically and radiographically.

\medterm{Metatarsalgia} Pain in the front or ball of the foot, often arising from pressure or irritation around the metatarsal heads.

\medterm{Metatarsals} The five bones forming the front half of the foot between the ankle and toes. They support body weight and help with walking; overuse or injury can cause pain or fractures.

\medterm{Metatarsus} The part of the foot between the tarsal bones and toes, containing the five metatarsal bones.

\medterm{Metatarsus Adductus} A foot deformity in which the forefoot is angled inward relative to the hindfoot, often recognized in infancy and ranging from flexible to rigid.

\medterm{Metatarsus Varus} Metatarsus varus is inward deviation of the forefoot caused by medial angulation of the metatarsal bones. It is often recognized in infancy and may be flexible or rigid; management depends on severity and persistence.

\medterm{Metencephalon} The metencephalon is the embryonic brain region that develops into the pons and cerebellum. These structures participate in motor coordination, balance, cranial-nerve functions, and relay of signals between the brain and spinal cord.

\medterm{Meteorism} Abdominal distension caused by excessive gas in the gastrointestinal tract.

\medterm{Meter} A metric unit of length equal to 100 centimeters; medical usage may also refer to a measuring device.

\medterm{Metered-Dose Inhaler} A pressurized inhaler that releases a measured amount of medicine as a spray. Correct timing and breathing are needed for the medicine to reach the lungs; a spacer can make use easier and improve delivery.

\medterm{Metestrus} The stage of the reproductive cycle after ovulation in many mammals, when the corpus luteum develops and progesterone rises. The exact timing and biological features differ among species.

\medterm{Metformin} A commonly used medicine for type 2 diabetes that lowers the amount of sugar made by the liver and helps the body use insulin better. Nausea or diarrhea are common; serious illness or severe kidney problems increase rare toxicity risk.

\medterm{Metformin Hydrochloride} Metformin hydrochloride is the hydrochloride salt of metformin, an oral biguanide used chiefly for type 2 diabetes. It lowers blood glucose without directly stimulating pancreatic insulin release and is available in immediate- and extended-release formulations.

\medterm{Metformin In Prediabetes} Metformin may reduce progression from prediabetes to type 2 diabetes, especially in people at
high risk. Its use depends on age, body size, prior gestational diabetes, kidney function,
other conditions, and patient preference; lifestyle intervention remains important.

\medterm{Methacholine Challenge Test} A breathing test used when asthma is suspected but routine breathing tests are normal. The person inhales increasing amounts of methacholine, which can narrow sensitive airways; breathing is measured after each dose under medical supervision.

\medterm{Methacholine Chloride} Methacholine chloride is a muscarinic cholinergic agonist used in bronchoprovocation testing. In susceptible people it causes airway narrowing, allowing assessment for bronchial hyperresponsiveness under controlled medical supervision.

\medterm{Methadone} A long-acting opioid medicine used in selected pain cases and treatment of opioid-use disorder. It reduces withdrawal and cravings but can cause dangerous breathing suppression, sedation, dependence, or heart-rhythm changes.

\medterm{Methadone Hydrochloride} Methadone hydrochloride is the hydrochloride salt of methadone. It provides long-lasting opioid-receptor activity for selected analgesic and opioid-use-disorder indications, with risks of respiratory depression, dependence, overdose, and cardiac arrhythmia.

\medterm{Methadone Overdose} Taking too much methadone can cause profound sleepiness, pinpoint pupils, slow or stopped breathing, coma, and death. Effects may be delayed or prolonged, and abnormal heart rhythms can occur; call emergency services immediately.

\medterm{Methaemoglobinaemia} A condition in which hemoglobin is changed so it cannot carry or release oxygen normally. It may cause blue or gray skin, headache, breathlessness, confusion, or severe oxygen shortage after exposure to certain medicines or chemicals.

\medterm{Methamidophos} Methamidophos is a highly toxic organophosphorus insecticide that inhibits acetylcholinesterase. Exposure can cause cholinergic poisoning with sweating, salivation, vomiting, bronchospasm, bradycardia, muscle weakness, and potentially respiratory failure.

\medterm{Methamphetamine} Methamphetamine is a potent central nervous system stimulant that increases synaptic catecholamines. It has limited medical use in some settings but is commonly misused; toxicity and chronic use can cause cardiovascular complications, psychosis, dependence, and neurocognitive harm.

\medterm{Methamphetamine Overdose} Taking too much methamphetamine can cause agitation, dangerously high temperature, fast heartbeat, high blood pressure, chest pain, seizures, stroke, abnormal rhythms, or severe confusion. This is an emergency requiring immediate medical care.

\medterm{Methamphetamines} Powerful stimulant drugs that increase activity in the brain and body. Repeated or high-dose use can cause dependence, severe heart problems, overheating, sleeplessness, paranoia, hallucinations, seizures, stroke, or death.

\medterm{Methane} Methane is a colorless, odorless, highly flammable gas produced by anaerobic decomposition and normal intestinal microbial activity. It can displace oxygen in enclosed spaces and cause asphyxiation; it is not itself a direct systemic poison at ordinary exposure levels.

\medterm{Methane Sulfonate} Methane sulfonate, or mesylate, is the anion or ester derived from methanesulfonic acid. Mesylate salts are commonly used to improve the formulation of medicines; some methanesulfonate esters are potent alkylating agents and toxic impurities.

\medterm{Methanobrevibacter Smithii} A methane-producing archaeon commonly found in the human intestine. It consumes hydrogen made by other gut microorganisms and may influence intestinal fermentation and gas production.

\medterm{Methanol} A poisonous alcohol whose breakdown products can cause severe acid buildup, blindness, brain injury, or death. Suspected ingestion is a medical emergency requiring immediate specialized treatment, even before symptoms appear.

\medterm{Methanol Poisoning} Poisoning by methanol, a type of alcohol found in some industrial products. Its breakdown can cause severe acid buildup, blindness, brain injury, or death; urgent hospital treatment is needed even before symptoms appear.

\medterm{Methanol Test} A blood test that measures methanol, a poisonous alcohol. It helps confirm toxic exposure, but treatment may begin before results because methanol can cause severe metabolic acidosis and delayed blindness or brain injury.

\medterm{Methazolamide} Methazolamide is a carbonic anhydrase inhibitor used mainly to lower intraocular pressure in glaucoma. It reduces aqueous humor production and can cause metabolic acidosis, electrolyte abnormalities, kidney stones, or sulfonamide-related reactions.

\medterm{Methemoglobin} Hemoglobin in which iron has been oxidized so it cannot carry oxygen normally. High levels cause methemoglobinemia, with cyanosis, headache, breathlessness, confusion, or severe tissue oxygen deficiency.

\medterm{Methemoglobinemia} Methemoglobinemia occurs when hemoglobin iron is oxidized and cannot carry oxygen normally, causing cyanosis, headache, breathlessness, or severe tissue hypoxia.

\medterm{Methenamine} Methenamine is a urinary antiseptic that decomposes in acidic urine to release formaldehyde, which inhibits bacterial growth. It is used for prevention of recurrent urinary tract infection in selected patients, not for treatment of active upper-tract infection.

\medterm{Methicillin-Resistant Staphylococcus Aureus} A strain of Staphylococcus aureus that is resistant to several common antibiotics. It can cause skin, wound, lung, blood, bone, or other infections, and treatment is chosen using the infection site and laboratory testing.

\medterm{Methicillin-Resistant Staphylococcus Aureus Infection} Alternate name; see \textit{MRSA Infection}.

\medterm{Methidathion} Methidathion is an organophosphorus insecticide and acaricide that inhibits acetylcholinesterase. Exposure can produce a cholinergic toxidrome with secretions, bronchospasm, bradycardia, gastrointestinal symptoms, weakness, and respiratory failure.

\medterm{Methimazole} A medicine that reduces production of thyroid hormones and is used mainly for hyperthyroidism and Graves disease. It can rarely cause severe low white-cell counts or liver injury; fever or sore throat during treatment needs prompt medical advice.

\medterm{Methiocarb} Methiocarb is a carbamate pesticide that inhibits acetylcholinesterase reversibly. Poisoning causes excessive muscarinic and nicotinic cholinergic activity and may require urgent supportive care and antidotal treatment.

\medterm{Methionine} An essential amino acid containing sulfur that the body uses to build proteins and make other substances. People must obtain it from food, and it can be converted into homocysteine and related compounds.

\medterm{Methionine Enkephalin} Methionine enkephalin is an endogenous opioid pentapeptide that binds opioid receptors and participates in modulation of pain, stress, reward, and neuroendocrine function. It is one of the products of proenkephalin processing.

\medterm{Methocarbamol} Methocarbamol is a centrally acting skeletal-muscle relaxant used with rest and physical measures for acute painful musculoskeletal conditions. It can cause drowsiness, dizziness, impaired coordination, and additive central nervous system depression.

\medterm{Methohexital} Methohexital is an ultra-short-acting barbiturate anesthetic used for induction of general anesthesia and, in selected settings, procedural sedation. It depresses central nervous system activity and may cause respiratory or cardiovascular depression.

\medterm{Methomyl} Methomyl is a carbamate insecticide that reversibly inhibits acetylcholinesterase. Poisoning produces cholinergic symptoms such as salivation, sweating, vomiting, bronchospasm, bradycardia, weakness, and respiratory compromise.

\medterm{Methotrexate} An antimetabolite that inhibits folate-dependent enzymes and is used in cancer and selected autoimmune or inflammatory disorders.

\medterm{Methotrexate In Dermatology} An immunomodulatory antimetabolite used in selected inflammatory and autoimmune skin
diseases, including psoriasis, eczema, and some connective-tissue disorders. Treatment
requires attention to liver disease, blood counts, infection risk, drug interactions,
and strict pregnancy precautions.

\medterm{Methotrimeprazine} Methotrimeprazine, also called levomepromazine, is a phenothiazine with antipsychotic, sedative, antiemetic, and analgesic properties. It is used in some countries for psychiatric illness and palliative-care symptoms and can cause sedation, hypotension, and anticholinergic effects.

\medterm{Methoxamine} Methoxamine is a predominantly alpha-adrenergic agonist that increases vascular tone and blood pressure. It has been used as a vasopressor for hypotension but is not a routine first-line agent in many current practice settings.

\medterm{Methoxsalen} Methoxsalen is a psoralen that intercalates into DNA and, after exposure to UVA, forms cross-links that suppress abnormal lymphocyte or skin-cell activity. It is used with UVA phototherapy for selected skin diseases; photosensitivity and skin-cancer risk require precautions.

\medterm{Methoxyflurane} A vapor medicine formerly used to cause general anesthesia and, in some settings, provide short-term pain relief. High or repeated exposure can damage the kidneys, so use requires careful medical supervision.

\medterm{Methyclothiazide} Methyclothiazide is a thiazide diuretic that inhibits sodium and chloride reabsorption in the distal convoluted tubule. It is used for hypertension or edema and may cause hypokalemia, hyponatremia, hyperuricemia, and glucose intolerance.

\medterm{Methyl Alcohol} Another name for methanol, a toxic alcohol that can cause severe acidosis and blindness.

\medterm{Methyl Bromide} Methyl bromide, or bromomethane, is a volatile fumigant and alkylating chemical. Inhalation can cause severe neurologic, respiratory, renal, and cardiovascular toxicity, including delayed seizures and pulmonary edema.

\medterm{Methyl Chloride} Methyl chloride, or chloromethane, is a volatile industrial gas formerly used in refrigerants and other applications. High-level inhalation exposure can depress the central nervous system and cause headache, ataxia, seizures, or cardiac effects.

\medterm{Methyl Isocyanate} Methyl isocyanate is a highly volatile and reactive industrial chemical. Inhalation or eye and skin exposure can cause severe irritation, pulmonary edema, and delayed respiratory injury; it was the toxic agent in the Bhopal disaster.

\medterm{Methyl Parathion} Methyl parathion is an organophosphorus insecticide that inhibits acetylcholinesterase. Exposure can produce severe cholinergic poisoning with secretions, bronchospasm, bradycardia, weakness, seizures, and respiratory failure.

\medterm{Methyl Salicylate Overdose} Taking too much methyl salicylate, found in some oils and pain-relief products, can cause salicylate poisoning with vomiting, ringing in the ears, rapid breathing, confusion, seizures, acidosis, or coma. Emergency assessment is needed.

\medterm{Methylation} Attachment of a small chemical group to DNA, proteins, or other molecules. On DNA, this can change how strongly a gene is used and helps regulate normal development, cell function, and some diseases.

\medterm{Methylcellulose} A plant-fiber product used as a bulk-forming laxative and as a thickener in medicines and foods. It absorbs water and adds bulk to stool; it should be taken with enough fluid and avoided when bowel blockage is suspected.

\medterm{Methyldopa} A blood-pressure medicine that acts in the brain to reduce nerve signals that tighten blood vessels. It is used in selected cases, including pregnancy, and may cause sleepiness, dizziness, or liver problems.

\medterm{Methylene Blue} A medicine used mainly for serious methemoglobinemia, a condition in which blood cannot carry oxygen normally. It may be unsafe in some people with an inherited red-cell enzyme deficiency and can interact with certain antidepressants.

\medterm{Methylene Blue Test} A test or procedure using methylene blue dye to trace fluid flow, identify a structure, or detect a leak or abnormal connection. The purpose depends on the body site; allergy and blood-related complications are uncommon but possible.

\medterm{Methylene Chloride} Methylene chloride, or dichloromethane, is a volatile industrial solvent. Inhalation can cause central nervous system depression and metabolism to carbon monoxide can produce tissue hypoxia; high exposure may be fatal.

\medterm{Methylenedioxymethamphetamine} 3,4-Methylenedioxymethamphetamine, abbreviated MDMA and commonly called ecstasy or molly, is a psychoactive drug with stimulant and hallucinogenic effects that increases serotonergic signaling and can cause serious toxicity.

\medterm{Methylergometrine} An ergot alkaloid that increases uterine contraction and may be used to control postpartum hemorrhage when appropriate.

\medterm{Methylmalonic Acid Blood Test} A blood test measuring methylmalonic acid to help detect vitamin B12 deficiency, especially when the B12 result is borderline. Kidney impairment can raise the result and must be considered during interpretation.

\medterm{Methylmalonic Acid Level} The amount of methylmalonic acid measured in blood or urine. A high level can support vitamin B12 deficiency, especially when the B12 result is unclear, but kidney disease can also raise it.

\medterm{Methylmalonic Acidemia} An inherited metabolic disorder in which the body cannot properly process certain fats and proteins, causing methylmalonic acid to build up. Episodes may cause vomiting, dehydration, acidosis, low blood sugar, seizures, or developmental and kidney problems.

\medterm{Methylmalonic Aciduria} A disorder in which methylmalonic acid builds up in blood and urine, usually because of an inherited enzyme or vitamin B12-processing problem. It may cause vomiting, poor feeding, acidosis, developmental problems, kidney disease, or neurologic injury.

\medterm{Methylmercury Compounds} Methylmercury compounds are organic mercury compounds that can accumulate in aquatic food chains and cross the blood--brain and placental barriers. Exposure primarily injures the nervous system and can impair fetal neurodevelopment.

\medterm{Methylmercury Poisoning} Poisoning caused by methylmercury, usually through contaminated food. It primarily injures the nervous system and may cause numbness, poor coordination, vision or hearing changes, and impaired fetal brain development; exposure requires medical evaluation.

\medterm{Methylmethacrylate} Methyl methacrylate is a volatile monomer used to make acrylic plastics and dental or orthopedic bone cement. Vapors may irritate the eyes and airways, and liquid exposure can cause dermatitis or sensitization.

\medterm{Methylmorphine} Another name for codeine, an opioid medicine used for selected pain or cough indications. It can cause drowsiness, constipation, dependence, and dangerous breathing suppression, especially with alcohol or other sedatives.

\medterm{Methylnaltrexone} Methylnaltrexone is a peripherally acting mu-opioid-receptor antagonist used to treat opioid-induced constipation when laxatives are insufficient. It generally preserves central analgesia but can cause abdominal pain, diarrhea, and rarely gastrointestinal perforation.

\medterm{Methylnaltrexone Bromide} Methylnaltrexone bromide is the bromide salt of methylnaltrexone, a peripherally restricted opioid antagonist. It reverses opioid effects in the gastrointestinal tract and is used for opioid-induced constipation under indication-specific precautions.

\medterm{Methylphenidate} A stimulant medicine used mainly to treat attention-deficit/hyperactivity disorder and narcolepsy, a condition causing excessive sleepiness. It affects brain chemicals involved in attention and alertness; reduced appetite, sleep problems, increased blood pressure, and misuse are possible risks.

\medterm{Methylprednisolone} A synthetic glucocorticoid used to suppress inflammation and immune activity.

\medterm{Methyltestosterone} A synthetic androgen related to testosterone. It can affect sexual development and other androgen-sensitive tissues but may cause liver injury, adverse cholesterol changes, infertility, or virilization and is used only in selected medical situations.

\medterm{Metoclopramide} A medicine that helps the stomach empty and relieves nausea and vomiting. It is used for selected digestive problems and treatment-related nausea; prolonged or high-dose use can cause involuntary movements and other nervous-system effects.

\medterm{Metolazone} A water-removing medicine used for high blood pressure and swelling. It helps the kidneys remove salt and water and is often combined with another diuretic, but can disturb sodium, potassium, kidney function, or hydration.

\medterm{Metolcarb} Metolcarb is a carbamate insecticide that can inhibit acetylcholinesterase. Exposure may cause cholinergic poisoning with excessive secretions, gastrointestinal symptoms, bronchospasm, bradycardia, weakness, or respiratory compromise.

\medterm{Metopic Ridge} A metopic ridge is a palpable or visible line along the forehead from fusion of the metopic suture; it may be normal or associated with trigonocephaly.

\medterm{Metoprolol} Metoprolol is a relatively beta-1-selective blocker used for hypertension, angina, heart failure, rate control, and some post-infarction settings; bradycardia, hypotension, fatigue, and bronchospasm can occur.

\medterm{Metric System} The metric system is a decimal system of measurement based on units such as the meter, kilogram, and liter. Modern medicine generally uses the International System of Units, with conventional units retained for some measurements.

\medterm{Metritis} Inflammation or infection of the uterus, often after childbirth, miscarriage, abortion, or a uterine procedure. It may cause fever, lower abdominal pain, unusual discharge, bleeding, or tenderness and usually needs treatment.

\medterm{Metronidazole} An antibiotic and antiparasitic medicine used for certain infections caused by bacteria that grow without oxygen and by some parasites. It can cause nausea or a metallic taste; nerve problems may occur with prolonged use.

\medterm{Metronidazole-Associated Encephalopathy} A rare complication of prolonged or high-dose metronidazole therapy, causing cerebellar
dysfunction (ataxia, dysarthria, nystagmus) and sometimes confusion or seizures. MRI may
show T2-hyperintense lesions in the cerebellar dentate nuclei, corpus callosum splenium,
and brainstem. Symptoms typically resolve with drug discontinuation.

\medterm{Metronomic Chemotherapy} Cancer treatment given in frequent, usually smaller doses over a long period instead of occasional high doses. The goal is to slow cancer growth and sometimes its blood supply; use depends on the cancer and treatment plan.

\medterm{Metrorrhagia} Uterine bleeding that occurs at unexpected times, including between menstrual periods. Causes include pregnancy, hormonal changes, fibroids, polyps, infection, medicines, and cancer; new, heavy, prolonged, or postmenopausal bleeding needs medical assessment.

\medterm{METs} Metabolic equivalents used to estimate activity intensity or energy use. One MET is approximately the energy used while resting; higher values indicate more demanding activity and help describe exercise capacity.
\synonyms
Metabolic Equivalents

\medterm{Metyrapone} A medicine that blocks one step in the production of cortisol, a hormone made by the adrenal glands. It is used for selected hormone tests or treatment under specialist supervision because levels can change quickly.

\medterm{Metzenbaum Scissors} Fine, lightweight surgical scissors designed to separate and cut delicate soft tissue. Their relatively narrow blades are useful for careful dissection, but they are not intended for cutting hard material or heavy sutures.

\medterm{Mevinphos} Mevinphos is a highly toxic organophosphorus insecticide that inhibits acetylcholinesterase. Exposure can rapidly cause cholinergic crisis, including bronchorrhea, bronchospasm, muscle weakness, seizures, and respiratory failure.

\medterm{Mexacarbate} Mexacarbate is a carbamate insecticide that interferes with acetylcholinesterase. Toxic exposure may produce muscarinic and nicotinic cholinergic symptoms and requires prompt decontamination and medical evaluation.

\medterm{Mexiletine} Mexiletine is an oral class IB antiarrhythmic and sodium-channel blocker used for selected ventricular arrhythmias and sometimes neuropathic pain. It can cause gastrointestinal and neurologic effects and may worsen or provoke arrhythmias in some patients.

\medterm{Meyerozyma Guilliermondii} Meyerozyma guilliermondii is a yeast formerly classified as Candida guilliermondii. It is usually an environmental or colonizing organism but can cause opportunistic bloodstream or other infections, particularly in immunocompromised patients.

\medterm{MF} Abbreviation; see \textit{Mycosis Fungoides}.

\medterm{MFS} Abbreviation; see \textit{Myxofibrosarcoma}.

\medterm{MG} Abbreviation; see \textit{Myasthenia Gravis}.

\medterm{MGUS} See \textit{Monoclonal Gammopathy of Undetermined Significance}.

\medterm{MHC Molecule} A protein on the surface of a cell that displays small pieces of proteins to T cells. This helps the immune system recognize infected or abnormal cells and distinguish the body's own cells from foreign material.

\synonyms
Major Histocompatibility Complex Molecule

\medterm{Mianserin} An antidepressant medicine used in some countries to treat depression. It changes signaling by several brain chemicals and can cause sleepiness, increased appetite, weight gain, dizziness, or other medicine-related effects.

\medterm{MIBG Scan} A nuclear medicine scan that uses a small amount of radioactive MIBG to find certain tumors of nerve-related tissue. It is used especially for pheochromocytoma, paraganglioma, and neuroblastoma; medicines and thyroid protection may be needed before the scan.

\medterm{MIBG Scintiscan} A nuclear medicine scan using radiolabeled metaiodobenzylguanidine to image tissues derived from the sympathetic nervous system, including selected pheochromocytomas, paragangliomas, and neuroblastomas.

\medterm{Mica} A group of sheet silicate minerals; inhalation of mica dust can cause occupational lung disease.

\medterm{Micafungin} An echinocandin antifungal agent used in the treatment and prophylaxis of invasive
candidiasis and esophageal candidiasis. Micafungin inhibits beta-(1,3)-D-glucan
synthase, disrupting fungal cell wall synthesis and leading to osmotic instability and
cell death.

\medterm{Micelle} A tiny cluster of molecules that helps carry fats and fat-soluble substances through intestinal fluid for absorption.

\medterm{Miconazole} An antifungal medicine used in creams, gels, or other preparations for infections of the skin, mouth, or vagina. It stops fungi from making an essential part of their outer covering; irritation and medicine interactions are possible, especially with some oral products.

\medterm{Microabscess} A very small pocket of pus containing immune cells, damaged tissue, and sometimes germs. It may occur in the skin or an organ because of infection or inflammation and is usually a finding rather than a diagnosis.

\medterm{Microalbuminuria} A small but abnormal amount of albumin, a blood protein, in the urine. It can be an early sign of kidney damage, especially in diabetes or high blood pressure, but temporary illness or exercise can also raise it and repeat testing may be needed.

\medterm{Microalbuminuria Test} A urine test measuring small amounts of albumin, usually by an albumin-to-creatinine ratio, to detect early kidney damage, especially in diabetes or hypertension.

\medterm{Microaneurysm} A very small balloon-like bulge in a blood vessel wall. In the retina, microaneurysms are an early sign of diabetic eye disease and may leak fluid, causing swelling and blurred vision. Their number and associated retinal changes help show how diabetes is affecting the eyes.

\medterm{Microangiopathic Hemolysis} Destruction of red blood cells as they pass through damaged or narrowed small blood vessels. Blood tests may show fragmented cells, anemia, and other signs of red-cell breakdown; causes include severe high blood pressure, pregnancy complications, infection, and clotting disorders.

\medterm{Microangiopathic Hemolytic Anemia} An anemia caused by mechanical fragmentation of red blood cells as they pass through damaged or narrowed small blood vessels, producing schistocytes and often elevated hemolysis markers.

\medterm{Microangiopathy} Disease or damage affecting the body's smallest blood vessels. Thickened or injured vessel walls can reduce blood flow or allow leakage and may damage organs, especially in diabetes, high blood pressure, or clotting disorders.

\medterm{Microarray} A laboratory test using a small chip covered with DNA probes. It can measure activity of many genes or find selected missing or extra DNA segments and is used mainly for specialized genetic or research testing.

\medterm{Microarray Analysis} Microarray analysis measures the abundance or hybridization of many nucleic-acid sequences simultaneously using an array of immobilized probes. It can assess gene expression, copy-number changes, or selected genetic variants.

\medterm{Microbacterium} A genus of gram-positive bacteria found in soil and water that can rarely cause bloodstream or other opportunistic infections.

\medterm{Microbe} A microorganism, such as a bacterium, fungus, protozoan, or virus.

\synonyms
Microorganism

\medterm{Microbial Biofilms} Microbial biofilms are organized communities of microorganisms attached to a surface and enclosed in a self-produced extracellular matrix. Biofilms can form on tissues and medical devices and often show increased tolerance to antimicrobials and host defenses.

\medterm{Microbicide} A microbicide is a chemical or biological agent that kills microorganisms or inactivates them. The term may refer to disinfectants, topical agents, or investigational products intended to reduce sexual transmission of infection.

\medterm{Microbiology} The study of microorganisms, including bacteria, viruses, fungi, and parasites, and how they interact with people, animals, plants, and the environment. In medicine, it supports the diagnosis, prevention, and treatment of infectious diseases.

\medterm{Microbiome} The community of microorganisms and their genetic material living in a body site, especially the intestine. These organisms can influence digestion, immunity, and health, although the effects differ by person and condition.

\medterm{Microcalcification} A very small calcium deposit in tissue. On a mammogram, its shape and arrangement may reflect normal change, injury, infection, or early breast cancer, so suspicious findings may require further imaging or biopsy.

\medterm{Microcephaly} Microcephaly is head size substantially smaller than expected for age and sex, reflecting genetic, developmental, infectious, toxic, or acquired brain causes.

\medterm{Microcirculation} The movement of blood through the body's smallest blood vessels. In these vessels, oxygen and nutrients pass to tissues and waste products pass into the blood; poor microcirculation can slow healing and damage organs.

\medterm{Microcolon} Microcolon is an abnormally small-caliber colon, often identified on contrast imaging. It may result from reduced fetal intestinal use or obstruction and is associated with conditions such as small-left-colon syndrome, meconium ileus, or distal intestinal obstruction.

\medterm{Microcornea} Microcornea is an abnormally small cornea, usually defined in adults by a horizontal diameter below about 10 mm. It may be isolated or associated with microphthalmia, glaucoma, cataract, or genetic syndromes.

\medterm{Microcurie} A microcurie ($\mu$Ci) is a unit of radioactivity equal to one-millionth of a curie, or approximately 37 kilobecquerels. It may be used to express the activity of radioactive tracers or sources.

\medterm{Microcystic Adnexal Carcinoma} A rare, slow-growing but locally aggressive malignant adnexal tumor, usually on the head or neck, with follicular and ductal differentiation; it may extend deeply and requires complete excision.

\medterm{Microcytic Anemia} A form of anemia characterized by abnormally small and pale circulating red blood cells, defined on a complete blood count by a mean corpuscular volume (MCV) below $80\,\text{fL}$. It occurs when the body is unable to synthesize adequate hemoglobin, starving red cells of their internal content. The five primary causes are commonly remembered by the mnemonic TAILS: Thalassemia (inherited globin chain defect), Anemia of chronic disease/inflammation (hepcidin-mediated iron trapping), Iron deficiency anemia (the most common cause, from chronic blood loss, poor diet, or pregnancy), Lead poisoning (toxic inhibition of heme enzymes), and Sideroblastic anemia (impaired iron incorporation in the bone marrow). Symptoms include fatigue, pallor, weakness, and shortness of breath; diagnostic evaluation includes serum ferritin, iron studies, and hemoglobin electrophoresis.
\synonyms{Microcytic Hypochromic Anemia; Small-Cell Anemia; Low-MCV Anemia}

\medterm{Microdermabrasion} A cosmetic treatment that gently removes the outermost layer of skin with fine crystals or a rough-tipped device. It may improve the look of mild uneven texture, dullness, or some marks, but can cause redness, irritation, infection, or temporary darkening.

\medterm{Microdialysis} A testing method that uses a very small tube to collect chemicals from the fluid around cells. It can measure local levels of medicines or natural substances, including in the brain or beneath the skin.

\medterm{Microdiscectomy} A minimally invasive spinal operation used to relieve pain, numbness, or weakness caused by a herniated intervertebral disc compressing a spinal nerve root in the lower back (lumbar radiculopathy or sciatica). Through a small incision in the midline of the lower back, a spine surgeon or neurosurgeon uses a high-magnification surgical operating microscope or surgical loupes to gently retract the overlying back muscle, remove a tiny sliver of bone (laminotomy), and visualize the inflamed nerve. The herniated or extruded disc fragments pressing against the nerve are carefully removed while leaving the healthy structural portion of the disc intact. Because tissue trauma is minimal compared to traditional open spine surgery, patients experience dramatic pain relief, often go home the same day or following morning, and resume light activities within a few weeks.

\synonyms
Lumbar Microdiscectomy; Microdiskectomy; Minimally Invasive Discectomy

\medterm{Microdissection} Microdissection is separation or isolation of very small tissue regions or structures under microscopic or highly magnified visualization. It is used in surgery, pathology, and molecular research to obtain selected cells or tissue.

\medterm{Microdissection Testicular Sperm Extraction} A surgical procedure that uses a microscope to search the testicle for small areas making sperm. It may help some men with very low or absent sperm production obtain sperm for assisted reproduction. The procedure can cause pain, bleeding, infection, or reduced testicular function, and sperm may not be found.

\synonyms
Micro-TESE; Microsurgical TESE

\medterm{Microdontia} Abnormally small teeth, which may affect one or several teeth and can be inherited or developmental.

\medterm{Microfibril} A very small thread-like structure that supports tissues. In connective tissue, microfibrils help organize elastic fibers and provide strength; in plants, they are made mainly of cellulose.

\medterm{Microfilament} A microfilament is a thin cytoskeletal filament composed mainly of actin, about 7 nm in diameter. Microfilaments support cell shape, motility, contraction, cytokinesis, and intracellular transport.

\medterm{Microfilament Proteins} Proteins that make, organize, move, or regulate actin microfilaments. They help cells maintain shape, migrate, divide, contract, and transport materials; examples include actin, myosin, and actin-binding proteins.

\medterm{Microfilaria} The tiny, immature form of a filarial worm that may circulate in blood or live in the skin. Mosquitoes, flies, or other biting insects can pick them up and spread the infection to another person.

\medterm{Microglandular Adenosis} Microglandular adenosis is an uncommon benign proliferation of small glands in the breast, often associated with hormonal influences such as pregnancy or progestin use. It can mimic invasive carcinoma and is diagnosed histologically.

\medterm{Microglia} Microglia are resident immune cells of the central nervous system. They survey neural tissue, remove debris, and respond to injury or infection; persistent activation can contribute to neuroinflammation.

\medterm{Microgliosis} Microgliosis is proliferation or increased activation of microglia in the central nervous system, usually in response to injury, infection, degeneration, or other inflammation. It is a histopathologic finding rather than a specific disease.

\medterm{Micrognathia} Micrognathia is an unusually small lower jaw that may affect feeding, teeth, facial development, or airway breathing.

\medterm{Micrognathism} Micrognathism is abnormal smallness of the jaw, usually referring to the mandible. It may be congenital or acquired and can affect feeding, dentition, airway patency, speech, or facial development.

\medterm{Micrograms} A unit of mass equal to one millionth of a gram, written as mcg. It is commonly used for small medicine doses and nutrient amounts; confusing micrograms with milligrams can cause a serious dosing error.

\medterm{Micrographia} Abnormally small handwriting, often becoming progressively smaller within a sentence. It commonly occurs in Parkinson disease and other movement disorders because of reduced movement size and control.

\medterm{Microinjection} Microinjection is delivery of a very small volume of fluid into a cell, tissue, or other microscopic target using a fine needle or micropipette. It is used in assisted reproduction, experimental biology, and some therapeutic procedures.

\medterm{Micrometastasis} A small group of cancer cells that has spread from the original tumor but is too small to detect on routine scans or examination. It may be found in tissue samples and can affect staging and treatment decisions.

\medterm{Micrometer} An instrument or unit used to measure very small distances; one micrometer is one-millionth of a meter.

\medterm{Micromolar} A concentration unit equal to one millionth of a mole per liter, written as \(\mu\mathrm{mol}/L\).

\medterm{Micromole} A micromole ($\mu$mol) is an amount of substance equal to one-millionth of a mole, or approximately $6.022\times10^{17}$ entities. It is commonly used for small quantities and concentrations of biochemical substances.

\medterm{Micron} A unit measuring one-millionth of a meter, commonly used to describe the size of cells, bacteria, and particles.

\medterm{Microneme} A microneme is a small secretory organelle in many apicomplexan parasites. It releases proteins that support host-cell attachment, motility, and invasion during the parasite life cycle.

\medterm{Micronization} Micronization is a process that reduces particles to micrometer-scale dimensions, often to increase surface area and improve dissolution or absorption of a drug. The resulting particle size and formulation still determine clinical performance.

\medterm{Micronized} Processed into very small particles to improve dispersion or absorption.

\medterm{Micronutrient} A micronutrient is a vitamin or mineral required in small amounts for normal growth, metabolism, and physiologic function. Deficiency or excess can cause disease, and requirements vary with age, pregnancy, illness, and other factors.

\medterm{Micronutrients} Vitamins and minerals required in small amounts for enzymatic, hormonal, structural,
hematologic, neurologic, and immune functions. They include fat-soluble vitamins (A, D,
E, and K), water-soluble vitamins, and essential minerals.

\medterm{Microorchidism} Abnormally small testes, which may result from impaired development, genetic syndromes, or testicular atrophy after injury or infection.

\medterm{Microorganism} A microscopic living organism, such as a bacterium, fungus, protozoan, or alga; viruses are acellular infectious agents.

\medterm{Microorganism Genetics} The study of genetic material, inheritance, variation, and gene activity in microorganisms. It helps explain how microbes grow, cause disease, develop drug resistance, and exchange useful or harmful traits.

\medterm{Micropenis} An unusually small penis caused by inadequate androgen effect or other developmental conditions, defined clinically by a stretched penile length more than 2.5 standard deviations below the mean for age and population. Evaluation may identify endocrine or genetic causes.

\medterm{Microphallus} An abnormally small penis, also called micropenis, usually defined relative to age and developmental standards; it may reflect hormonal or developmental disorders.

\medterm{Microphthalmia} A birth defect in which one or both eyes are smaller than usual.
It may occur with other eye or genetic conditions and can reduce vision. Treatment
depends on the cause and may include glasses, treatment for poor vision, or surgery.

\medterm{Microphthalmos} A birth condition in which one or both eyeballs are unusually small. Vision may be reduced, and other eye or genetic problems may occur, so affected children need an eye examination and ongoing specialist care.

\medterm{Microplastics} Very small plastic particles, generally less than five millimeters across, found in the environment and sometimes in food, water, and air. Human health effects are still being studied.

\medterm{Micropsia} A change in vision in which objects appear smaller or farther away than they really are. It can occur with migraine, retinal disease, brain disorders, medicines, or other problems and should be assessed if it is new, persistent, or associated with neurological symptoms.

\medterm{MicroRNA} A very small, non-coding RNA that helps control which genes are active, mainly by reducing the production of specific proteins. MicroRNAs are being studied as possible disease markers and treatment targets.

\medterm{Microsatellite Instability} A tumor pattern in which short repeated DNA sequences gain or lose units because the cell’s DNA repair system is faulty. A high level can occur in inherited or sporadic cancers and may help predict benefit from some immune treatments.

\medterm{Microscope} An instrument that magnifies structures too small to see clearly with unaided eyes. Light, electron, and fluorescence microscopes help examine cells, tissues, blood, microbes, and other samples.

\medterm{Microscopic} Visible only with a microscope or relating to microscopic examination.

\medterm{Microscopic Anatomy} The study of cells, tissues, and fine structural details that require magnification for examination; it includes cytology and histology.

\medterm{Microscopic Colitis} Long-lasting inflammation of the colon that usually causes frequent, watery, non-bloody diarrhea while the colon looks normal during routine colonoscopy. Tiny tissue samples are needed for diagnosis; the main forms are collagenous and lymphocytic colitis.

\synonyms
Collagenous Colitis; Lymphocytic Colitis

\medterm{Microscopic Examination} Examination of cells or tissue with a microscope to look for normal structure, infection, inflammation, or cancer. The sample may be stained and assessed together with its visible appearance, imaging, and other laboratory results.

\medterm{Microscopic Hematuria} Alternate name; see \textit{Hematuria}.

\medterm{Microscopic Polyangiitis} An autoimmune disease that inflames and damages small blood vessels. It can affect the kidneys, lungs, skin, nerves, or other organs and may cause blood in the urine, breathing problems, rash, weakness, or bleeding.

\medterm{Microscopy} Microscopy is examination of cells, tissues, organisms, or materials with a microscope. Light, fluorescence, electron, and other microscopy methods provide different information about structure and composition.

\medterm{Microsecond} A unit of time equal to one-millionth of a second.

\medterm{Microspectrophotometry} Microspectrophotometry measures the absorption or emission spectrum of a microscopic region or individual cell. It can quantify pigments, nucleic acids, proteins, or other chromophores in small biological samples.

\medterm{Microsphere} A microsphere is a small spherical particle, often made of polymer, glass, or another material, used in drug delivery, imaging, embolization, or laboratory assays. Its size, composition, and surface properties determine its biologic behavior.

\medterm{Microspherophakia} Microspherophakia is an abnormally small, spherical crystalline lens. It can occur in inherited disorders such as Weill--Marchesani syndrome and may cause lens displacement, pupillary block, angle-closure glaucoma, or reduced vision.

\medterm{Microsporidia} A group of microscopic spore-forming organisms related to fungi. They can cause prolonged diarrhea, eye disease, muscle disease, or widespread infection, especially in people with weakened immune systems.

\medterm{Microsporidiosis} An infection caused by microsporidia, microscopic organisms related to fungi that produce resistant spores. It most often affects the intestine and may cause long-lasting diarrhea and weight loss. It can also affect the eyes, muscles, kidneys, brain, or several organs, especially in people with weakened immune systems.

\medterm{Microsporum} Microsporum is a genus of dermatophyte fungi that can cause ringworm of the skin, scalp, or hair.

\medterm{Microsporum Audouinii} Microsporum audouinii is a dermatophyte that mainly infects the scalp and can cause tinea capitis, especially in children.

\medterm{Microsporum Distortum} Microsporum distortum is a rare dermatophyte related to M. canis that can cause skin or scalp infection.

\medterm{Microsporum Ferrugineum} A fungus that can cause ringworm of the scalp or skin, particularly in children. Laboratory examination can help confirm the infection and distinguish it from other causes of rash or hair loss.

\medterm{Microsporum Gypseum} Microsporum gypseum is a soil-associated dermatophyte that can cause ringworm of the skin or scalp.

\medterm{Microstomia} Microstomia is abnormal narrowing or smallness of the mouth opening. It may be congenital or result from scarring, burns, surgery, trauma, or connective-tissue disease and can impair eating, speech, oral hygiene, and dental treatment.

\medterm{Micro-Surgery} Surgery performed with magnification and fine instruments to repair very small structures such as nerves and vessels.

\medterm{Microsurgery} Microsurgery uses magnification and very small instruments to repair tiny blood vessels, nerves, or other structures.

\medterm{Microtechnology} Microtechnology is the design and manufacture of devices or structures with micrometer-scale features. Medical applications include microfluidic diagnostics, sensors, drug-delivery systems, and minimally invasive instruments.

\medterm{Microtia} A birth defect in which the outer ear is smaller than usual or
does not form fully. The ear canal may also be narrow or absent, which can reduce hearing.

\medterm{Microtome} An instrument that cuts extremely thin slices of tissue for examination under a microscope. The slices are mounted on slides, stained, and studied to identify normal structures or disease.

\medterm{Microtubule} A tiny hollow protein tube inside a cell. Microtubules help maintain cell shape, move materials within the cell, support cell division, and form part of cilia and other moving structures.

\medterm{Microtubule Bundle} A microtubule bundle is an organized group of parallel or otherwise coordinated microtubules within a cell. Bundles provide structural support and guide transport or movement, including in axons, cilia, and dividing cells.

\medterm{Microvascular} Relating to the body's smallest blood vessels, including tiny arteries, capillaries, and veins. These vessels deliver oxygen and nutrients, remove waste, and control blood flow through tissues.

\medterm{Microvascular Disease} Disease affecting the body's smallest blood vessels, which can reduce blood flow and oxygen delivery to organs or tissues. Diabetes, high blood pressure, smoking, inflammation, and aging can damage these vessels.

\medterm{Microvasculature} The network of the body's smallest blood vessels, including tiny arteries, capillaries, and veins. It delivers oxygen and nutrients, removes waste, controls tissue blood flow, and helps regulate fluid movement between blood and tissues.

\medterm{Microvillus} A microvillus is a microscopic, actin-supported projection of a cell membrane that increases surface area. Microvilli are abundant on intestinal epithelial cells, where they support nutrient absorption, and on renal tubular cells.

\medterm{Microvolt} A microvolt ($\mu$V) is one-millionth of a volt. It is used to express small bioelectric signals, such as those recorded by electroencephalography, electromyography, electrocardiography, or evoked-potential testing.

\medterm{Microwave Ablation} A procedure that uses microwave energy to heat and destroy selected tissue, usually under imaging guidance.

\medterm{Microwave Therapy} Treatment that uses microwave energy to heat or destroy selected tissue. It has specialized uses, such as treating some tumors or controlling bleeding, and the risks depend on the body area and device used.

\medterm{Micturition} The process of passing urine. The bladder squeezes while the muscles at its outlet relax, under coordinated control from the brain and nerves.

\medterm{Midaxillary Line} An imaginary vertical line on the side of the body passing through the apex of the
axilla (armpit).

\medterm{Midazolam} Midazolam is a short-acting benzodiazepine used for procedural sedation, anesthesia, and treatment of acute seizures in appropriate settings. It enhances GABA-A receptor activity and can cause profound sedation and respiratory depression, especially with opioids or other depressants.

\medterm{Midazolam Hydrochloride} Midazolam hydrochloride is the hydrochloride salt of the short-acting benzodiazepine midazolam. It is used for procedural sedation, anesthesia, and selected seizure emergencies and can cause dose-related sedation, hypotension, and respiratory depression.

\medterm{Mid-Brain} The midbrain, or mesencephalon, is the upper part of the brainstem involved in eye movements, motor pathways, and auditory and visual reflexes.

\medterm{Midbrain} The upper part of the brainstem, between the lower brainstem and the main brain. It helps control eye movements, movement, alertness, and automatic responses to sights and sounds.

\medterm{Midclavicular Line} An imaginary vertical line on the front of the body passing through the midpoint of the
clavicle (collarbone).

\medterm{Middle Ear} An air-filled space between the eardrum and inner ear that contains the malleus, incus, and stapes. These ossicles transmit
sound vibrations to the inner ear, and the auditory tube helps equalize middle-ear
pressure.

\medterm{Middle East Respiratory Syndrome} A serious respiratory infection caused by a coronavirus, usually linked to close contact with infected camels or people. It may cause fever, cough, pneumonia, breathing failure, or kidney and other organ problems.

\synonyms
MERS; Camel Flu

\medterm{Middle East Respiratory Syndrome} Alternate name; see \textit{MERS}.

\medterm{Middle Insomnia} Middle insomnia is repeated or prolonged awakening during the middle portion of the sleep period, often with difficulty returning to sleep. It may occur with stress, mood disorders, sleep apnea, pain, substances, medications, or other sleep disorders.

\medterm{Middle Lobe Syndrome} Chronic or recurrent atelectasis and infection of the right middle lobe due to bronchial
obstruction. The long, narrow right middle lobe bronchus is susceptible to compression
by lymph nodes, tumors, or mucous plugging. It may be caused by bronchiectasis,
malignancy, or granulomatous disease. Treatment depends on the underlying cause.

\medterm{Mid-Dorsal} Situated in the middle region of the back or on the dorsal side of a structure.

\medterm{Mid-Forearm} The middle part of the forearm, between the elbow and wrist. Measurements or injections described as mid-forearm use this region as a reference point.

\medterm{Mid-Gut} The embryonic middle portion of the primitive gut, which develops into much of the small intestine and proximal colon.

\medterm{Midgut Volvulus} A life-threatening twisting of the intestine that can cut off its blood supply, usually because the intestine formed abnormally before birth. Babies may suddenly develop green vomiting, a swollen or painful abdomen, blood in the stool, or extreme illness. It is an emergency requiring immediate imaging and surgery to untwist and protect the bowel.

\synonyms
Volvulus Neonatorum; Midgut Torsion

\medterm{Mid-Leg} The central part of the lower limb between the knee and ankle. It contains the shin bones, muscles, nerves, blood vessels, and soft tissues that support walking and ankle movement.

\medterm{Mid-Line} The imaginary or anatomical line dividing a structure or body into right and left portions.

\medterm{Midomafetamine} Midomafetamine is another name for 3,4-methylenedioxymethamphetamine, a psychoactive amphetamine derivative. It is being studied in controlled psychotherapy-assisted research but can cause hyperthermia, hyponatremia, cardiovascular toxicity, serotonin toxicity, and misuse.

\medterm{Midostaurin} Midostaurin is a multikinase inhibitor used with chemotherapy for newly diagnosed acute myeloid leukemia with a susceptible FLT3 mutation and for advanced systemic mastocytosis. It can cause nausea, cytopenias, infection, and other treatment-related toxicities.

\medterm{Midriff} Alternate informal term; see \textit{Abdomen}.

\medterm{Mid-Stream Urine Specimen} A urine sample collected after initial urine is voided, reducing contamination from the urethral opening and used commonly for culture.

\medterm{Mid-Systolic Click} A sharp heart sound heard during the middle or later part of the heartbeat, often associated with mitral-valve prolapse. It may be followed by a murmur if the valve allows blood to leak backward.

\medterm{Midtrimester} The middle part of pregnancy, usually covering about weeks 14 through 27. During this period the fetus grows rapidly, routine anatomy screening is often performed, and pregnancy symptoms and risks vary from person to person.

\medterm{Mid-Upper-Arm Circumference} The circumference of the upper arm measured at the midpoint between the acromion and
olecranon processes. It is a practical anthropometric measure of muscle and fat stores,
particularly useful when body weight or height cannot be measured reliably.

\medterm{Midwife} A trained and regulated professional who provides care during pregnancy, labor, birth, and the postnatal period.

\medterm{Midwifery} The professional practice of supporting and managing normal pregnancy, childbirth, postpartum recovery, and newborn care.

\medterm{Mifepristone} Mifepristone is a progesterone-receptor antagonist and, at higher doses, a glucocorticoid-receptor antagonist. It is used with misoprostol for medication abortion under applicable protocols and for selected cases of hyperglycemia associated with Cushing syndrome.

\medterm{Migraine} A recurring neurological disorder that causes attacks of moderate to severe, often throbbing or pulsing head pain. The pain may occur on one or both sides of the head and may be accompanied by nausea, vomiting, or increased sensitivity to light, sound, or smells. Some people have an aura, such as flashing lights, blind spots, tingling, numbness, or temporary speech problems, which usually develops gradually and resolves completely. Migraine attacks may last from several hours to a few days, and some occur with little or no headache.

\synonyms
Migraine Headache; Hemicrania

\medterm{Migraine And Stroke Symptoms} Migraine aura usually develops gradually and resolves, whereas stroke symptoms often begin suddenly and cause persistent focal deficits such as weakness, speech trouble, facial droop, or vision loss. New or atypical neurologic symptoms require emergency stroke assessment.

\medterm{Migraine With Aura} A primary headache disorder characterized by recurrent attacks of fully reversible
visual, sensory, or speech symptoms that develop gradually over at least five minutes
and last no more than sixty minutes, followed by headache with the characteristics of
migraine.

\medterm{Migraine Without Aura} A primary headache disorder characterized by recurrent attacks of moderate to severe,
pulsating, unilateral headache lasting four to seventy-two hours, associated with
nausea, vomiting, photophobia, and phonophobia.

\medterm{Migrainous Neuralgia} An older term for cluster headache, a disorder causing recurrent attacks of severe unilateral pain, usually around one eye, with ipsilateral autonomic symptoms such as tearing or nasal congestion.

\medterm{Migrating Motor Complex} Repeating waves of muscle activity that move through the stomach and small intestine between meals. They help clear leftover food, secretions, and some bacteria before the next meal.

\medterm{Migration} Movement of cells, organisms, substances, or a body structure from one place to another. Normal examples include immune-cell movement and fetal development; abnormal migration can contribute to infection, cancer spread, or blood clots.

\medterm{Migratory Thrombophlebitis} Recurrent or shifting venous thrombosis associated with an underlying malignancy; an alternate term for the thrombotic manifestation of Trousseau syndrome.

\medterm{Mikulicz Disease} An older term for persistent, usually symmetric enlargement and inflammation of the lacrimal and salivary glands, now commonly associated with IgG4-related disease.

\medterm{Mikulicz Syndrome} A historical term for symmetrical, persistent swelling of the tear and salivary glands, often associated with IgG4-related disease or other autoimmune conditions. The underlying cause should be evaluated rather than assumed.

\medterm{Mikulicz's Disease} An older term for persistent, usually symmetrical enlargement of the lacrimal and salivary glands, now often considered within IgG4-related disease or Sjögren syndrome.

\medterm{Milan System for Salivary Cytopathology} A six-level system for reporting results from a needle sample of a salivary-gland lump. Results range from an insufficient sample or noncancerous finding to a finding suspicious for or showing cancer. The system helps estimate cancer risk and guide further care.

\medterm{Mild} Mild describes relatively low intensity, severity, or extent of a symptom, sign, disease, or abnormality. The meaning depends on the condition and the clinical scale used and does not by itself indicate that no treatment or evaluation is needed.

\medterm{Mild Cognitive Impairment} A noticeable decline in memory or another thinking skill that is greater than expected for age but does not substantially prevent everyday activities. It may remain stable, improve, or progress to dementia depending on the cause.

\medterm{Mildew Remover Poisoning} Poisoning from a mildew-removal product, often containing bleach, acids, or other corrosive chemicals. It can burn the skin, eyes, mouth, and digestive tract or release dangerous fumes; leave the area and seek urgent advice.

\medterm{Mile} A mile is a unit of distance equal to 1,609.344 meters. It may be used in descriptions of walking or running distance, although scientific and clinical measurements generally use SI units.

\medterm{Milia} Milia are tiny, firm, white or yellow bumps caused by small cysts containing keratin, a protein in the skin. They are common in newborns and may appear on the face in adults; they are harmless and usually go away on their own.

\medterm{Miliaria} A sweat-retention disorder caused by blockage of eccrine ducts, producing small itchy or prickly papules.

\medterm{Miliaria Crystallina} A harmless heat rash made of tiny clear blisters when sweat becomes trapped just beneath the skin’s surface. It can occur in hot, humid conditions or during fever and usually settles when the skin is cooled and kept dry.

\medterm{Miliaria Rubra} An itchy prickly heat rash caused by sweat-duct blockage in the epidermis, producing small red papules or vesicles.

\medterm{Miliary} Describing many tiny spots or lesions that resemble millet seeds. The term is often used for tuberculosis or another disease spread widely through an organ or the bloodstream.

\medterm{Miliary TB} Tuberculosis that has spread through the bloodstream, producing many tiny areas of infection in the lungs and sometimes the liver, brain, bone marrow, or other organs. It can cause fever, weight loss, night sweats, cough, or organ-specific symptoms.

\medterm{Miliary Tuberculosis} Tuberculosis that spreads through the bloodstream and forms many tiny areas of infection. It may affect the lungs, liver, brain, bone marrow, or other organs and can cause prolonged fever, weight loss, night sweats, cough, or organ-specific symptoms.

\medterm{Milieu Therapy} Milieu therapy uses the structured social environment of a treatment setting as part of care, especially in some psychiatric and rehabilitation programs.

\medterm{Military Personnel} Military personnel are members of the armed forces, including active-duty, reserve, and related service members. Their health care may address occupational exposures, injuries, infectious risks, combat-related trauma, and deployment-related mental-health conditions.

\medterm{Milium} A small, firm, white or yellow bump filled with keratin, usually appearing on the face. It is harmless, does not spread by infection, and may disappear on its own or be removed safely by a clinician.

\medterm{Milk} A liquid made by the mammary glands to nourish an infant. Human breast milk provides water, energy, protein, fats, vitamins, and immune substances; its amount and composition change during feeding and as the infant grows.

\medterm{Milk Allergy} An immune reaction to proteins in cow's milk that may cause skin, gastrointestinal, or respiratory symptoms and, rarely, anaphylaxis.

\synonyms
Cow's Milk Protein Allergy; CMPA

\medterm{Milk Intolerance} Alternate name; see \textit{Lactose Intolerance}.

\medterm{Milk Intolerance} Alternate name; see \textit{Lactose Intolerance}.

\medterm{Milk Of Magnesia} A liquid suspension of magnesium hydroxide used to neutralize stomach acid and relieve constipation. It draws water into the bowel and can cause diarrhea, especially when taken too often or in kidney disease.

\medterm{Milk-Alkali Syndrome} A condition caused by taking too much calcium and absorbable alkali, leading to high blood calcium, excessive blood alkalinity, and kidney problems. Severe cases can cause confusion, vomiting, dehydration, or abnormal heart rhythms.

\medterm{Milker's Node} A localized nodular skin infection caused by parapoxvirus acquired from infected cattle, also called orf-like milker's nodule.

\medterm{Milk-Leg} An older term for postpartum swelling of a leg, historically associated with deep venous thrombosis or lymphatic obstruction.

\medterm{Milky Urine} Urine with a cloudy or white appearance caused by chyle, pus, crystals, fat, or other substances; evaluation identifies the cause.

\medterm{Miller's Lung} Hypersensitivity pneumonitis caused by inhaling material associated with the granary weevil (wheat weevil, \textit{Sitophilus granarius}), especially in people working with contaminated grain or flour.

\medterm{Millicurie} A millicurie (mCi) is one-thousandth of a curie, equal to approximately 37 megabecquerels. It is used to express the activity of radioactive sources and some radiopharmaceutical doses.

\medterm{Milliequivalent} A unit used to measure the amount and electrical combining power of charged substances, such as sodium, potassium, calcium, or chloride, in medicines and laboratory results.

\medterm{Milligram} A metric unit of mass equal to one-thousandth of a gram. Its abbreviation is mg, and it is commonly used for medicine doses and small quantities of substances.

\medterm{Milligram-Percent} An older concentration unit meaning milligrams of a substance in 100 milliliters of fluid. Its number is the same as milligrams per deciliter, written mg/dL, which is preferred in current medical reporting.

\medterm{Milligray} A unit of absorbed ionizing-radiation dose equal to one-thousandth of a gray, or 0.001 joule per kilogram; its abbreviation is mGy.

\medterm{Milliliter} A unit of volume equal to one-thousandth of a liter and one cubic centimeter. It is commonly used to measure liquid medicines, blood, urine, laboratory samples, and other small amounts of fluid.

\medterm{Millimeter} A unit of length equal to one-thousandth of a meter.

\medterm{Millimolar} A concentration unit equal to one thousandth of a mole per liter, written as mmol/L.

\medterm{Millimole} A unit measuring the amount of a chemical substance equal to one-thousandth of a mole. Laboratory tests use millimoles to report substances such as glucose, sodium, potassium, and other electrolytes.

\medterm{Million} A number equal to 1,000,000; laboratory counts may use it as a numerical multiplier.

\medterm{Milliosmole} A milliosmole (mOsm) is one-thousandth of an osmole, a unit reflecting the number of osmotically active particles in a solution. Osmolality and osmolarity are commonly reported in mOsm/kg or mOsm/L.

\medterm{Millipede Toxin} Irritating defensive chemicals released by some millipedes. Contact can cause burning, redness, blistering, or dark staining of skin and severe eye irritation; rinse exposed areas with water and seek care for eye symptoms.

\medterm{Millisecond} A millisecond (ms) is one-thousandth of a second. It is used to measure brief physiologic events, reaction times, nerve conduction intervals, and timing of electrical or imaging signals.

\medterm{Milliunit} A milliunit (mU) is one-thousandth of a unit of biologic or enzymatic activity. Because a unit is defined according to the substance or assay, the clinical meaning of a milliunit depends on the specific test or medication.

\medterm{Millivolt} A unit of electrical voltage equal to one-thousandth of a volt. Small voltages are measured in millivolts in tests of the heart, brain, nerves, muscles, and other electrical activity in the body.

\medterm{Milrinone} A medicine given by vein for short-term support in severe heart failure or after heart surgery. It helps the heart squeeze more strongly and widens blood vessels, but can lower blood pressure or trigger abnormal heart rhythms.

\medterm{Milroy Disease} An inherited primary lymphedema caused most often by pathogenic variants affecting
VEGFR3 signaling. It typically presents with chronic swelling of the lower limbs from
birth or early infancy and may be associated with papillomatosis, hydrocele, or
recurrent skin infection.

\medterm{Milroy's Disease} An inherited form of primary lymphedema, usually present from birth or early life and caused commonly by lymphatic-development abnormalities.

\medterm{Mimesis} Imitation or mimicry; in medicine, the term can refer to simulation of an organic disease by psychological symptoms or to one organic disease resembling another.

\medterm{Mimicry} Resemblance of one organism, structure, or clinical pattern to another; in medicine, it may describe a disorder that imitates another disease.

\medterm{Minamata Disease} A nervous-system disorder caused by poisoning with methylmercury, often after eating contaminated seafood. It can cause numbness, poor coordination, vision or hearing problems, muscle weakness, speech difficulty, and developmental harm in children.

\medterm{Mind} The collection of mental processes and capacities involved in consciousness, thought, perception, memory, emotion, judgment, and behavior.

\medterm{Mindful Eating} Paying attention to hunger and fullness, eating slowly, and noticing food's taste and texture. It may improve awareness of eating habits and help some people reduce distracted eating or overeating.

\medterm{Mineral} An inorganic element or compound required for body structure, fluid balance, enzyme activity, or other physiological functions.

\medterm{Mineral Oil} Mineral oil is a purified mixture of liquid hydrocarbons derived from petroleum. It is used in some topical products, cosmetics, and laxatives; aspiration of liquid mineral oil can cause lipoid pneumonia.

\medterm{Mineral Oil Overdose} Swallowing a large amount of mineral oil is dangerous mainly because it can enter the lungs during vomiting or reflux and cause lipoid pneumonia. Do not induce vomiting; obtain product-specific poison-control advice.

\medterm{Mineral Spirits Poisoning} Poisoning from mineral spirits, petroleum-based solvents used in paints and cleaning products. Swallowing or inhalation may cause dizziness, drowsiness, chemical pneumonitis, or breathing failure, especially after aspiration; do not induce vomiting.

\medterm{Mineralocorticoid} A steroid hormone, mainly aldosterone, that helps the kidneys retain sodium and water while removing potassium. Too much can raise blood pressure and lower potassium; too little can cause dehydration and weakness.

\medterm{Mineralocorticoid Receptor} A protein inside cells that responds mainly to aldosterone, a hormone made by the adrenal glands. It helps the kidneys retain salt and water and remove potassium, thereby affecting blood pressure and body-fluid balance.

\medterm{Mineralocorticoids} Steroid hormones, chiefly aldosterone, that promote sodium retention and potassium and hydrogen secretion in the kidney.

\medterm{Minerals} Naturally occurring elements needed by the body for bones, nerves, muscles, blood formation, fluid balance, and other functions. They include calcium, iron, magnesium, potassium, sodium, and zinc; both deficiency and excess can be harmful.

\medterm{Miner's Anaemia} Anemia associated historically with mining, commonly due to hookworm infection, chronic blood loss, nutritional deficiency, or occupational exposure.

\medterm{Miner's Asthma} Asthma caused or worsened by repeated inhalation of dust, fumes, or other substances during mining work. It can cause wheezing, cough, chest tightness, and breathlessness that improve after leaving the exposure.

\medterm{Miner's Nystagmus} An older term for visual disturbance and nystagmus reported in underground miners, associated with poor lighting and occupational conditions.

\medterm{Minimal Change Disease} A kidney disorder, especially common in children, in which the kidney filters look nearly normal on routine examination but leak large amounts of protein. It can cause swelling around the eyes, legs, or abdomen and is often treatable.

\medterm{Minimal Residual Disease} A small number of cancer cells that remain after treatment and cannot be seen on routine examination or ordinary scans. Highly sensitive blood, bone-marrow, tissue, or molecular tests may detect them and estimate relapse risk.

\medterm{Minimally Invasive Gastrointestinal Surgery} Gastrointestinal operations performed through small incisions or natural orifices using laparoscopic, robotic, or endoscopic techniques, often reducing recovery time while retaining procedure-specific risks.

\medterm{Minimally Invasive Hip Replacement} Hip-replacement surgery performed through a smaller or muscle-sparing approach to reduce tissue disruption. It still replaces damaged joint surfaces and has risks such as infection, blood clots, dislocation, fracture, nerve injury, or implant loosening.

\medterm{Minimally Invasive Mitral Valve Surgery} Surgical mitral valve repair or replacement performed through a right mini-thoracotomy or video-assisted thoracoscopic approach, avoiding a traditional median sternotomy. Results in reduced surgical trauma, less postoperative pain, and faster recovery.

\synonyms
MIMVS; Mini-Thoracotomy Mitral Repair; Port-Access Mitral Surgery

\medterm{Minimally Invasive Spine Surgery} Spine surgery performed through small incisions using specialized instruments, often reducing muscle disruption and recovery time. It still carries risks such as nerve injury, infection, bleeding, and persistent symptoms.

\medterm{Minimally Invasive Surgery} Surgery performed through small incisions or natural openings using specialized instruments, usually reducing tissue disruption and recovery time.

\medterm{Mini-Mental State Examination} A brief questionnaire that checks orientation, attention, memory, language, and simple thinking skills. The score can be affected by education, language, hearing, vision, culture, or delirium and cannot diagnose dementia by itself.

\medterm{Minimization} A psychological or communication process of downplaying the importance or severity of symptoms, events, or problems.

\medterm{Minimum Alveolar Concentration} A standard measure of how much inhaled anesthetic is needed to prevent movement in half of people during a specified surgical stimulus. It helps compare anesthetic strength and guide dosing during anesthesia.

\medterm{Minimum Bactericidal Concentration} The lowest concentration of an antimicrobial that kills a defined proportion of the
original bacterial inoculum, commonly demonstrated by subculture without growth. MBC is
distinct from MIC and is not routinely required for every infection.

\medterm{Minimum Effect Concentration} The lowest concentration of a drug in the blood or at its site of action that produces
a defined measurable effect.

\medterm{Minimum Inhibitory Concentration} The smallest amount of an antibiotic or other antimicrobial that stops a particular germ from visibly growing in a laboratory test. It helps guide treatment but must be interpreted with the infection and the patient's condition.

\medterm{Minimum Lethal Concentration} The lowest concentration of a substance associated with death in a defined test
population or experimental setting.

\medterm{Minimum Lethal Dose} The smallest dose of a substance reported to cause death under specified conditions; it varies by species, route, and individual.

\medterm{Minisatellite Repeat} A minisatellite repeat is a tandemly repeated DNA sequence with repeat units commonly ranging from about 10 to 100 base pairs. Variation in repeat number can serve as a genetic marker and, in some loci, influence gene regulation or disease risk.

\medterm{Mini-Stroke} Alternate name; see \textit{Transient Ischemic Attack}.

\medterm{Mini-Stroke} A lay term usually referring to a transient ischemic attack, in which focal neurologic symptoms resolve without persistent infarction.

\medterm{MINOCA} A heart attack occurring without a major blockage in the coronary arteries. It is an initial diagnosis; further tests are needed to find causes such as a small clot, artery spasm, artery-wall tear, or inflammation.

\medterm{Minocycline} An antibiotic used for some bacterial infections and for inflammatory acne. It can cause nausea, dizziness, skin or tooth discoloration, sun sensitivity, and rarely liver or immune problems; the course should be taken as prescribed.

\medterm{Minor} Smaller, lesser, or less significant than a related structure or condition; for example, the teres minor muscle is smaller than the teres major muscle.

\medterm{Minor Response} A minor response is a limited improvement after treatment that does not meet the predefined criteria for a partial or complete response. The exact threshold depends on the disease, study protocol, and outcome being measured.

\medterm{Minority Health} The health status and healthcare experiences of populations affected by social, economic, cultural, or structural disadvantage. Differences in exposure, income, environment, discrimination, and access to care can contribute to health disparities.

\medterm{Minoxidil} A vasodilator used topically for pattern hair loss and orally in selected cases of resistant hypertension.

\medterm{Minute Ventilation} The total amount of air breathed in one minute, calculated from breath size multiplied by breathing rate. It may rise with exercise or illness and fall with sedation, muscle weakness, airway disease, or poor breathing effort.

\medterm{Miosis} Unusually small pupils or normal pupil narrowing in response to light. It may be caused by medicines or toxins, eye disease, or problems affecting nerves in the brain or eye.

\medterm{Miosis} Abnormally small pupils. It may result from bright light, medicines such as opioids, eye disease, nerve disorders, or poisoning.

\medterm{Miosis Disorder} Abnormally small pupils or an abnormal pupillary constriction response, caused by medications, neurologic disease, ocular disease, or exposure to certain toxins.

\medterm{Miotic} A medicine or substance that makes the pupil smaller. Miotics may be used in glaucoma or during some eye procedures and can affect vision in dim light.

\synonyms
Miotic Agent; Pupillary Constrictor

\medterm{Miracidium} The ciliated larval stage of a trematode that hatches from an egg and infects a suitable snail host.

\medterm{Mirage Syndrome} A rare genetic disorder, usually related to SAMD9 variants, that can cause growth restriction, adrenal insufficiency, recurrent infections, and distinctive developmental findings.

\medterm{MIRD Schema} A method used to estimate how much radiation reaches organs and tissues after a radioactive medicine is given. It helps specialists plan treatment and assess safety for radiopharmaceutical procedures.

\medterm{Mirizzi Syndrome} A gallstone stuck near the gallbladder outlet presses on a nearby bile duct and blocks the flow of bile. It can cause jaundice, abdominal pain, fever, or infection and may require an endoscopic or surgical procedure.

\medterm{Mirror Therapy} A rehabilitation exercise in which a mirror makes movement of the stronger limb appear to come from the affected limb. It may help some people with phantom-limb pain or movement recovery after stroke.

\medterm{Mirror-Writing} Writing in which letters or words are reversed as if seen in a mirror; it may occur transiently in childhood or with neurologic disease.

\medterm{Mirtazapine} Mirtazapine is a noradrenergic and specific serotonergic antidepressant used for major depressive disorder. It commonly causes sedation and increased appetite or weight and can rarely cause neutropenia or other serious adverse effects.

\medterm{Misalignment} Misalignment is abnormal positioning or orientation of two or more structures relative to their intended relationship. It may affect bones, joints, teeth, eyes, implants, or medical devices and can impair function or healing.

\medterm{Miscarriage} Spontaneous loss of a pregnancy before fetal viability, commonly caused by chromosomal abnormalities but requiring clinical assessment.

\synonyms
Spontaneous Abortion; Early Pregnancy Loss

\medterm{Misdiagnosis} Misdiagnosis is assignment of an incorrect diagnosis or failure to recognize the correct condition. It can lead to inappropriate treatment, delayed care, avoidable harm, or unnecessary testing and may result from atypical presentation, incomplete information, or system factors.

\medterm{Miserere} An old term for vomiting that smells or looks like bowel contents because of severe blockage of the intestine. This is a medical emergency that may cause dehydration, infection, or bowel damage and needs urgent treatment.

\medterm{Mismatch Repair} A system that checks newly copied DNA and corrects small errors. When it fails, genetic changes can accumulate and increase the risk of cancers such as colorectal and womb cancer.

\medterm{Mismatch Repair Deficiency} Loss of DNA mismatch-repair function that produces microsatellite instability and a high
tumor mutational burden. Seen in Lynch syndrome and some sporadic colorectal,
endometrial, and other cancers, and associated with responsiveness to immune checkpoint
inhibitors.

\medterm{MISME} Abbreviation; see \textit{Neurofibromatosis Type 2}.

\medterm{Misophonia} Misophonia is an intense emotional or physiologic reaction to specific repetitive sounds, such as chewing or tapping. It is not simple annoyance; behavioral strategies, sound management, and therapy can reduce impairment, while assessment considers anxiety, obsessive-compulsive symptoms, and sensory sensitivity.

\medterm{Misoprostol} A prostaglandin E1 analogue used for ulcer prevention in selected patients, medical abortion, cervical preparation, and obstetric indications.

\medterm{Misoprostol} A synthetic prostaglandin analogue used for selected gynecologic and obstetric indications, including cervical ripening,
labor induction, and medical management of some pregnancies. It can cause uterine
contractions, bleeding, gastrointestinal effects, or uterine tachysystole; use depends on
the clinical setting and gestational age.

\synonyms
Prostaglandin E1 Analogue

\medterm{Missed Dose} A missed dose is a scheduled dose of a medicine that was not taken or administered. The appropriate response depends on the drug and timing; doubling a dose can be dangerous, so the product instructions or a clinician or pharmacist should guide the next dose.

\medterm{Missed Miscarriage} A nonviable intrauterine pregnancy that has not been expelled and may produce few or no symptoms of miscarriage.
Diagnosis requires appropriate ultrasound or other clinical evaluation before management is selected.

\synonyms
Missed Abortion; Silent Miscarriage; Early Pregnancy Failure

\medterm{Missense Mutation} A missense mutation is a DNA sequence change that substitutes one amino acid for another in a protein. Its effect may be harmful, beneficial, or neutral depending on the affected residue, protein function, and genetic context.

\medterm{Missense Variant} A single-nucleotide substitution that changes one amino acid in a protein. Its clinical
effect ranges from benign to pathogenic and must be assessed using molecular,
functional, population, and clinical evidence.

\medterm{Mistletoe Poisoning} Poisoning after eating mistletoe leaves or berries. It may cause nausea, vomiting, diarrhea, abdominal pain, dizziness, low blood pressure, slow or irregular heartbeat, seizures, or confusion; children and symptomatic people need urgent assessment.

\medterm{Mite} A tiny arachnid; some mites cause dermatitis or transmit disease, and house-dust mites can trigger allergy.

\medterm{Mites} Tiny arthropods related to ticks and spiders. Some live harmlessly on skin, while others cause allergy, scabies, or irritation; treatment depends on the species and condition.

\medterm{Mithramycin} An older antineoplastic antibiotic, also called plicamycin, whose use is limited by substantial toxicity.

\medterm{Mitobronitol} Mitobronitol is an alkylating antineoplastic compound formerly investigated or used in selected hematologic malignancies. Its cytotoxic effects can damage DNA and cause myelosuppression and other treatment-related toxicities.

\medterm{Mitochondria} Small structures inside cells that make much of the energy cells need. They also help control calcium, cell survival, and several metabolic processes; mitochondrial disorders can affect organs that need high energy, such as the brain, muscles, heart, and eyes.

\medterm{Mitochondrial Cristae} Folds of the inner membrane of mitochondria, the energy-producing parts of cells. The folds provide space for the machinery that turns nutrients into usable cellular energy.

\medterm{Mitochondrial Diabetes} Diabetes caused by a change in mitochondrial DNA, which is usually passed through the mother. It may occur with hearing loss, muscle problems, or other mitochondrial symptoms; insulin production is often reduced and genetic testing may confirm the cause.

\medterm{Mitochondrial Disease} A group of genetic disorders caused by impaired mitochondrial energy production. Symptoms vary and
may affect the brain, muscles, nerves, heart, eyes, or multiple organs; onset and severity can
differ widely even within one family.

\medterm{Mitochondrial DNA} Small circular DNA molecules located primarily in mitochondria that encode some components of oxidative phosphorylation and are inherited predominantly through the maternal line.

\medterm{Mitochondrial DNA Depletion Syndrome MDS MDDS} A group of inherited disorders in which some tissues contain too little mitochondrial DNA to make energy normally. Depending on the form, the brain, muscles, liver, or several organs may be affected, often beginning in infancy or childhood.

\medterm{Mitochondrial Fission} Mitochondrial fission is division of one mitochondrion into two or more mitochondria. It helps distribute mitochondria, remove damaged portions through mitophagy, and respond to cell-cycle or metabolic signals; excessive fission may contribute to disease.

\medterm{Mitochondrial Fusion} Mitochondrial fusion is joining of mitochondria into a connected network. It permits exchange of mitochondrial contents and helps maintain bioenergetic function and genome integrity; fusion and fission are normally balanced dynamically.

\medterm{Mitochondrial Inheritance} Inheritance of mitochondrial DNA, which is transmitted predominantly through the egg and therefore usually
from the mother. Examples include Leber hereditary optic neuropathy and MELAS. Heteroplasmy
means that mutant and normal mitochondrial DNA coexist in a cell.

\medterm{Mitochondrial Matrix} The mitochondrial matrix is the protein-rich compartment enclosed by the inner mitochondrial membrane. It contains mitochondrial DNA, ribosomes, and enzymes for the citric-acid cycle, fatty-acid oxidation, and other metabolic processes.

\medterm{Mitochondrial Membrane} The mitochondrial membrane system consists of an outer membrane and a highly folded inner membrane separated by an intermembrane space. It controls transport and houses the respiratory chain and ATP-producing machinery.

\medterm{Mitochondrial Myopathy} A muscle disorder caused by faulty mitochondria, the cell structures that make energy. It may cause weakness, exercise intolerance, drooping eyelids, hearing loss, seizures, or problems in other organs.

\medterm{Mitochondrial Ribosome} A mitochondrial ribosome is a ribosome within mitochondria that translates proteins encoded by mitochondrial DNA. Its structure differs from cytosolic ribosomes and defects in mitochondrial translation can cause energy-production disorders.

\medterm{Mitochondrial Swelling} Enlargement of mitochondria caused by disruption of ion and water balance, often indicating cellular injury, impaired oxidative phosphorylation, or permeability-transition changes.

\medterm{Mitochondrion} A mitochondrion is a membrane-bound organelle that generates much of a cell’s ATP through oxidative phosphorylation. It also participates in apoptosis, calcium handling, metabolism, and production or regulation of reactive oxygen species.

\medterm{Mitogen} A mitogen is a substance that stimulates a cell, especially a lymphocyte, to enter the cell cycle and proliferate. Mitogens include growth factors, lectins, and other signals that activate proliferative pathways.

\medterm{Mitolactol} Mitolactol is an alkylating antineoplastic compound formerly used or investigated in cancer treatment. It damages DNA and can cause cytotoxic effects such as myelosuppression and gastrointestinal toxicity.

\medterm{Mitomycin} Mitomycin is a chemotherapy medicine that damages DNA and is used for selected cancers; it can suppress bone marrow and injure the kidneys or bladder.

\medterm{Mitomycins} Mitomycins are antitumor antibiotics that become activated in cells and cross-link DNA, inhibiting replication and transcription. Mitomycin C is the best-known member and is used in selected cancers and ophthalmic or surgical applications.

\medterm{Mitophagy} The process by which a cell identifies and breaks down damaged or unnecessary mitochondria, the cell's energy-producing structures. This helps maintain healthy cells; abnormal mitophagy may contribute to some diseases.

\medterm{Mitosis} Cell division that produces two genetically similar daughter cells with the same chromosome number as the parent
cell.

\medterm{Mitosporic Fungi} Fungi that reproduce by asexual spores; the term is mainly used for fungi whose sexual stage is unknown or not observed.

\medterm{Mitotane} Mitotane is an adrenolytic drug used for selected adrenocortical carcinomas and to suppress excess adrenal steroid production. It can cause gastrointestinal and neurologic effects, adrenal insufficiency, and numerous drug interactions.

\medterm{Mitotic Anaphase} Mitotic anaphase is the stage of mitosis in which sister chromatids separate and move toward opposite poles of the cell. Accurate spindle attachment and chromosome segregation are required for formation of genetically similar daughter cells.

\medterm{Mitotic Figure} A cell seen under a microscope while it is dividing. The number of dividing cells in a tissue sample can help estimate how quickly a tumor or other tissue is growing.

\medterm{Mitotic Metaphase} Mitotic metaphase is the stage of mitosis in which duplicated chromosomes align at the cell equator with sister kinetochores attached to opposite spindle poles. A checkpoint helps prevent premature chromosome separation.

\medterm{Mitotic Prophase} Mitotic prophase is the early stage of mitosis in which chromatin condenses into visible chromosomes, the nucleolus disappears, and the mitotic spindle begins to form. In late prophase or prometaphase, the nuclear envelope breaks down.

\medterm{Mitotic Recombination} Mitotic recombination is exchange or rearrangement of DNA between homologous chromosomes during or around mitotic cell division. It can repair DNA or generate loss of heterozygosity and is relevant to tumor evolution and genetic disease.

\medterm{Mitotic Telophase} Mitotic telophase is the stage after chromosome separation when chromosomes reach opposite poles, decondense, and become enclosed in newly forming nuclear envelopes. Cytokinesis usually follows or overlaps with telophase.

\medterm{Mitoxantrone} A chemotherapy and immunosuppressive medicine that interferes with DNA and is used for selected cancers and some forms of multiple sclerosis.

\medterm{MitraClip} A transcatheter device that creates a double orifice in the mitral valve by clipping the
anterior and posterior leaflets together, reducing mitral regurgitation. It is indicated
for patients with significant (3+ or 4+) mitral regurgitation who are at high risk for
surgical mitral valve repair.

\medterm{Mitral} Relating to the heart's mitral valve, which lies between the left upper and lower chambers. The valve opens for blood flow forward and closes to prevent backward leakage during each heartbeat.

\medterm{Mitral Atresia} A serious heart defect present from birth in which the opening of the mitral valve is absent or completely closed. Blood cannot flow normally into the heart’s main pumping chamber, so newborns need urgent specialist care.

\medterm{Mitral Facies} A distinctive physical appearance characterized by a purplish or violaceous malar flush and patchy cyanosis across the cheeks, bridge of the nose, and lips. It is caused by chronic pulmonary hypertension, systemic vasoconstriction, and reduced cardiac output in severe mitral stenosis.

\medterm{Mitral Insufficiency} Failure of the mitral valve to close completely, allowing blood to leak back into the left atrium during systole.

\medterm{Mitral Regurgitation} Leakage through the mitral valve when the heart contracts, allowing blood to flow backward into the left upper chamber. It may cause breathlessness, tiredness, palpitations, or heart failure and can result from valve disease, infection, or heart damage.

\medterm{Mitral Stenosis} Narrowing of the mitral valve opening, making it harder for blood to move from the left upper chamber to the left lower chamber. It can cause breathlessness, tiredness, palpitations, or swelling and is often related to previous rheumatic heart disease.

\medterm{Mitral Stenosis Murmur} A low rumbling heart sound heard when blood passes through a narrowed mitral valve. It may be accompanied by a sharp opening sound and is best assessed with a stethoscope and heart imaging.

\medterm{Mitral Valve} The two-flap valve between the heart’s left upper and lower chambers. It opens to let blood pass forward and closes to prevent backward flow; it may become narrowed, leaky, infected, or weakened.

\medterm{Mitral Valve Insufficiency} Leakage of the mitral valve that allows blood to flow backward into the left atrium when the heart contracts. It may cause breathlessness, fatigue, palpitations, swelling, or heart failure.

\medterm{Mitral Valve Prolapse} A condition in which the flaps of the mitral valve bulge backward during a heartbeat. Many people have no symptoms, but some develop palpitations, chest discomfort, breathlessness, or leakage through the valve.

\synonyms
MVP; Barlow Syndrome; Click-Murmur Syndrome; Floppy Mitral Valve

\medterm{Mitral Valve Regurgitation} Backward leakage of blood through the mitral valve when the left ventricle contracts, which can cause breathlessness, fatigue, atrial enlargement, or heart failure.

\synonyms
MR; Mitral Incompetence; Mitral Insufficiency

\medterm{Mitral Valve Stenosis} Narrowing of the mitral valve that restricts blood flow from the left atrium to the left ventricle.

\synonyms
MS; Mitral Stenosis

\medterm{Mitral Valvotomy} A procedure that opens a narrowed mitral valve so blood can move more easily from the upper to the lower left heart chamber. It may use a balloon catheter or surgery and is suitable only for selected valve anatomy.

\medterm{Mitral Valvuloplasty} Mitral valvuloplasty is a procedure that widens a narrowed mitral valve, often with a balloon catheter or surgical repair.

\medterm{Mitral-Valve Repair} A surgical or transcatheter procedure to restore mitral valve competence without
replacement, typically for regurgitation from prolapse, chordal rupture, or ischemic
dysfunction. Repair preserves native valve tissue, maintains left ventricular function,
and avoids lifelong anticoagulation required by mechanical prostheses.

\medterm{Mitsuokella} Mitsuokella is a genus of anaerobic bacteria found mainly in the mouth and digestive tract; human infection is uncommon.

\medterm{Mittelschmerz} Recurrent one-sided lower-abdominal or pelvic pain occurring around ovulation. It is usually brief or lasts up to a day or two, but severe, persistent, or atypical pain requires assessment for other causes.

\medterm{Mixed Connective Tissue Disease} An autoimmune illness with features of lupus, scleroderma, and inflammatory muscle disease. It may cause color changes in fingers with cold, swollen hands, joint pain, muscle weakness, swallowing problems, or inflammation around organs.

\synonyms
MCTD; Sharp Syndrome

\medterm{Mixed Glioma} Mixed glioma is an older term for a glial tumor showing features of more than one glial lineage, historically including oligoastrocytoma. Modern classification relies mainly on molecular findings and generally assigns such tumors to a more specific entity.

\medterm{Mixed Hearing Loss} Hearing loss caused by both a problem carrying sound through the outer or middle ear and a problem in the inner ear or hearing nerve. The best treatment depends on the causes found.

\medterm{Mixed Oedema} Edema caused by more than one mechanism or affecting more than one tissue compartment; interpretation depends on clinical context.

\medterm{Mixed Urinary Incontinence} Urine leakage caused by both a sudden strong urge to urinate and pressure from coughing, sneezing, laughing, or exercise. Treatment may include bladder training, pelvic-floor exercises, medicines, or procedures.

\medterm{Mixed-Lineage Leukemia} An acute leukemia in which the malignant blasts show defining features of more than one hematopoietic lineage, commonly both myeloid and lymphoid; it is also called mixed-phenotype acute leukemia.

\medterm{Miyoshi Myopathy} A rare inherited muscle disease that usually begins in adolescence or early adulthood with progressive weakness of the calf muscles. People may have difficulty standing on tiptoe or climbing stairs; weakness later spreads to the thighs and hips. Blood tests often show high muscle-enzyme levels. Physical therapy helps maintain movement and function.

\synonyms
Miyoshi Distal Myopathy; Distal Dysferlinopathy

\medterm{MME} Membrane metalloendopeptidase, also called neprilysin, CD10, neutral endopeptidase, or CALLA; it is a cell-surface enzyme and marker expressed in several tissues and in some acute lymphoblastic leukemias.

\medterm{MéNièRe's Disease} An inner-ear disorder causing repeated attacks of spinning dizziness, changing hearing loss, ringing, and pressure or fullness in one ear. Symptoms may last minutes to hours and can affect balance and daily life.

\medterm{MéNièRe’S Disease} An inner-ear disorder causing repeated attacks of spinning dizziness, changing hearing loss, ringing, and pressure or fullness in one ear. Symptoms may last minutes to hours and can affect balance and daily life.

\synonyms
Meniere Disease; Endolymphatic Hydrops

\medterm{MéNéTrier’S Disease} A rare stomach disorder in which the stomach lining becomes unusually thick and loses protein into the digestive tract. It may cause upper-abdominal pain, nausea, vomiting, weight loss, swelling, or low blood protein.

\medterm{Mobile} Capable of moving or being moved; a mass or structure described as mobile is not fixed to adjacent tissue.

\medterm{Mobile Genetic Elements} Genetic elements that can move within and between genomes. These include plasmids,
transposons, and integrons. Mobile genetic elements can carry antibiotic resistance
genes and virulence factors. They facilitate horizontal gene transfer and contribute to
the rapid evolution of bacteria.

\medterm{Mobile Health} The use of phones, tablets, wearable devices, or similar portable technology to support health care. Uses include reminders, remote monitoring, education, communication, symptom tracking, and some diagnostic tools; privacy, accuracy, accessibility, and clinical oversight remain important.
\synonyms
mHealth; Digital Health

\medterm{Mobile Health Unit} A mobile health unit is a vehicle or transportable facility equipped to provide health services at changing locations. It may deliver screening, vaccination, primary care, diagnostic testing, emergency support, or outreach to underserved communities.

\medterm{Mobility} The ability to move or be moved, including functional movement of a person or range of motion of a joint.

\medterm{Mobilization} The process of restoring or encouraging movement, such as early patient activity, joint movement, or surgical release of tissue.

\medterm{Mobiluncus} A genus of curved anaerobic bacteria associated with bacterial vaginosis and occasionally isolated from other clinical specimens.

\medterm{Mobiluncus Curtisii} A curved bacterium that grows without oxygen and may be present with bacterial vaginosis. Its detection can support assessment, but symptoms, examination, and other test results are needed to interpret vaginal findings.

\medterm{Mobiluncus Mulieris} A curved bacterium that grows without oxygen and may be associated with bacterial vaginosis. Finding it can support a diagnosis, but symptoms, examination, and other laboratory results are also important.

\medterm{Mobius Syndrome} A rare condition present from birth that affects certain facial and eye-moving nerves. It can cause facial weakness or paralysis, trouble closing the eyes, limited outward eye movement, and feeding or speech difficulties.

\medterm{Modafinil} A prescription medicine that promotes wakefulness in narcolepsy, shift-work sleep disorder, and persistent sleepiness from obstructive sleep apnea when airway treatment is adequate. It can cause headache, anxiety, insomnia, serious rash, or dependence and is not a substitute for sleep or airway treatment.

\medterm{Modality} A particular method used for treatment, diagnosis, imaging, or sensory stimulation. Examples include surgery, medicines, ultrasound, MRI, physiotherapy, and forms of counseling; the best modality depends on the health problem and the person's needs.

\medterm{Mode} The most frequent value in a dataset or a particular manner or method of operation.

\medterm{Model} A simplified representation used to describe, explain, predict, or study a biological or clinical process.

\medterm{Moderate To Severe Disabilities} Moderate to severe disability describes substantial limitations in physical, cognitive, sensory, communication, or adaptive functioning that require ongoing support. Needs vary widely and may include medical care, rehabilitation, assistive technology, education, and assistance with daily activities.

\medterm{Moderation} Moderation is reduction or control of the intensity, amount, or severity of a process, symptom, behavior, or treatment effect. In clinical research, a moderator is a factor that changes the strength or direction of an association.

\medterm{Moderator} A factor that changes the strength or direction of a relationship or treatment effect; it may also mean a person who guides a discussion.

\medterm{Modes Of Inheritance} Patterns by which genetic traits pass through families, including autosomal dominant, autosomal recessive, X-linked, mitochondrial, and other patterns.

\medterm{Modification} A change made to a structure, process, treatment, or behavior.

\medterm{Modified Biophysical Profile} A pregnancy assessment that combines a fetal heart-rate test with measurement of the amniotic-fluid volume. A nonreactive heart-rate result or low fluid level may suggest that further assessment or delivery is needed, depending on gestational age and the clinical situation.

\synonyms
Modified BPP; NST with Amniotic Fluid Index

\medterm{Modulator} A substance or factor that alters the activity of a biological pathway, receptor, gene, or physiological process.

\medterm{Modus Operandi} The characteristic method by which a person performs an action; it is chiefly a forensic or legal term.

\medterm{Modus Vivendi} A practical arrangement or way of living that allows people or groups with differences to coexist.

\medterm{Moebius Syndrome} A congenital disorder characterized mainly by facial paralysis and inability to abduct the eyes, caused by developmental cranial-nerve abnormalities.

\medterm{Moellerella} A rare group of bacteria that can occasionally cause urinary, bloodstream, or other infection, mainly in people who are seriously ill or have weakened defenses. Laboratory identification is needed to determine whether a specimen represents true infection.

\medterm{Moellerella Wisconsensis} A rare type of bacterium that can occasionally cause urinary, bloodstream, wound, or other infections, especially in people with serious illness or weakened immunity. Laboratory testing identifies it and guides antibiotic choice.

\medterm{Moeller's Glossitis} An older term for a smooth, painful, inflamed tongue associated especially with nutritional deficiency, including iron, folate, or vitamin B12 deficiency.

\medterm{MOG Antibody-Associated Disorder} Alternate name; see \textit{MOG Antibody Disease}.

\medterm{MOGAD} Abbreviation; see \textit{MOG Antibody Disease}.

\medterm{Mogibacterium} A genus of anaerobic bacteria found in the mouth and intestine and occasionally associated with dental or other infection.

\medterm{Mohs Micrographic Surgery} A skin-cancer operation that removes the visible tumor a thin layer at a time and examines each layer under a microscope until no cancer remains. It preserves as much healthy tissue as possible; bleeding, infection, scarring, or recurrence can occur.

\medterm{Mohs Surgery} A micrographically controlled technique in which skin cancer is excised in thin layers
and examined margin-by-margin, achieving high cure rates for basal cell and squamous
cell carcinoma while sparing normal tissue.

\medterm{Moisture} Water or other liquid present on a surface or in tissue; excessive moisture can cause maceration and skin breakdown.

\medterm{Molar} Relating to a back tooth used for chewing, or to a chemical concentration measured in moles per liter.

\medterm{Molar Pregnancy} A type of gestational trophoblastic disease characterized by abnormal proliferation of
trophoblastic tissue. See Hydatidiform Mole.

\medterm{Molar Pregnancy} An abnormal pregnancy in which placental tissue grows excessively instead of developing normally. In a complete mole no embryo forms; in a partial mole abnormal fetal tissue may be present. It can cause bleeding, severe nausea, or unusually high pregnancy-hormone levels and requires specialist follow-up.

\synonyms
Hydatid Mole; Gestational Trophoblastic Disease

\medterm{Molasses} A viscous syrup produced during sugar refining; it is rich in sugar and may contain variable amounts of minerals.

\medterm{Mold} A multicellular fungus that grows as branching filaments; molds can cause allergy, food spoilage, or opportunistic infection.

\medterm{Mold Allergy} An allergic response to mold spores that may cause nasal, eye, skin, or lower-respiratory symptoms.

\synonyms
Fungal Hypersensitivity; Mold-Induced Rhinitis

\medterm{Mold Exposure} Mold exposure can irritate the eyes, skin, and airways and may worsen asthma or allergic rhinitis; invasive infection is uncommon and mainly affects severely immunocompromised people. Reducing moisture and remediating contaminated areas are more useful than unproven detoxification.

\medterm{Mole} A common usually harmless skin growth made of pigment-producing cells. It may be flat or raised and brown, black, pink, or skin-colored. A new or changing mole, or one unlike the others, should be checked for skin cancer.

\medterm{Molecular} Relating to molecules or processes occurring at the molecular level.

\medterm{Molecular Chaperone} A molecular chaperone is a protein that assists other proteins with folding, assembly, transport, or refolding without becoming part of the final structure. Chaperones can also help prevent or resolve protein aggregation.

\medterm{Molecular Diagnosis} Molecular diagnosis uses analysis of DNA, RNA, proteins, or other molecular markers to identify a disease, pathogen, inherited variant, tumor subtype, or treatment-relevant characteristic. Results require appropriate assay validation and clinical interpretation.

\medterm{Molecular Dynamics} A computer method that predicts how atoms and molecules move and interact over time. Researchers use it to study protein shape, drug binding, chemical reactions, and other biological processes; its results depend on the model and assumptions used.

\medterm{Molecular Imaging} Imaging that shows biological molecules, cell activity, or biochemical processes in the body, often using a targeted radioactive or fluorescent tracer. It can help characterize tumors, inflammation, infection, or organ function.

\medterm{Molecular Pathology} The study of disease-related changes in genes, gene activity, or proteins within cells and tissues. It can help identify a disease, classify a tumor, estimate likely outcome, detect inherited risk, or select a treatment.

\medterm{Molecular Pathology Of Colorectal Cancer} The assessment of molecular alterations that classify colorectal tumors and guide prognosis or therapy, including mismatch-repair deficiency or microsatellite instability, RAS mutations, BRAF mutation, HER2 amplification, and selected gene fusions.

\medterm{Molecular Probe} A molecular probe is a labeled or otherwise detectable molecule designed to bind a particular nucleic-acid sequence, protein, metabolite, or other target. Probes are used in imaging, assays, microscopy, and molecular diagnostics.

\medterm{Molecular Response} A molecular response is a measurable change in a molecular marker, such as a pathogen load, oncogenic transcript, mutation burden, or biomarker, after disease progression or treatment. Its definition and clinical significance depend on the disease and assay.

\medterm{Molecular Testing} Analysis of DNA, RNA, or other molecular features to identify variants, fusions, expression patterns, or infectious agents for diagnosis, prognosis, or treatment selection.

\medterm{Molecular Weight} The mass of one molecule, usually expressed in atomic mass units; for a substance used in chemistry, its numerical value corresponds to molar mass in grams per mole. It helps calculate concentrations and doses.

\medterm{Molecule} A group of chemically bonded atoms forming the smallest unit of a substance that retains its chemical properties.

\medterm{Moles} Moles are common collections of pigment-producing melanocytes and are usually benign. A lesion that changes in size, shape, color, symmetry, or symptoms, or differs markedly from others, should be examined for melanoma.

\medterm{Molimina} The symptoms and physical changes that occur before menstruation, including breast tenderness, bloating, or mood changes.

\medterm{Mollusc} A soft-bodied invertebrate; the term is also used in names of mollusc-borne or mollusc-associated parasitic diseases.

\medterm{Mollusca} The animal group containing snails, slugs, clams, and related species. Some mollusks can cause illness through venom, contaminated food, or parasites they carry; medical risk depends on the species and exposure.

\medterm{Molluscipoxvirus} A group of poxviruses that includes the virus causing molluscum contagiosum. Infection produces small, firm, usually painless skin bumps with a central dimple and spreads through close skin contact or shared items.

\medterm{Molluscum Contagiosum} A common skin infection caused by a poxvirus. It produces small, firm, painless, flesh-colored bumps with a central dip and spreads through close skin contact or shared objects. Most cases clear without treatment, although this may take months; scratching can spread the bumps or cause infection.

\medterm{Molten Metal Fumes} Fumes produced during welding, smelting, or casting that can cause metal-fume fever, airway irritation, or metal toxicity.

\medterm{Molteno Implant} A non-valved glaucoma drainage device (aqueous shunt) surgically implanted to reduce elevated intraocular pressure (IOP) in refractory and complicated glaucomas. It consists of a flexible silicone tube inserted into the anterior chamber (or pars plana) that shunts aqueous humor to one or two thin episcleral polypropylene plates sutured to the sclera beneath Tenon's capsule in the equatorial region, where fluid passively diffuses through a fibrous capsule into the orbital vascular and lymphatic circulation. Because the device is non-valved and offers no intrinsic resistance to flow, the tube is temporarily occluded at surgery (using an absorbable ligature or intraluminal ripcord suture) during the first 4--6 weeks until fibrous bleb encapsulation develops, thereby preventing early postoperative hypotony, anterior chamber collapse, and choroidal detachment. It is indicated when conventional filtration surgery (trabeculectomy) has failed or has a high probability of failure, including neovascular glaucoma, uveitic glaucoma, and post-penetrating keratoplasty glaucoma.

\synonyms
Molteno Glaucoma Drainage Device; Molteno Aqueous Shunt; Molteno Seton

\medterm{Molybdenum} An essential trace element required in small amounts for enzymes; excessive exposure can affect copper balance.

\medterm{Molybdenum-99} A radioactive form of molybdenum that changes into technetium-99m. It is used in generator systems to produce technetium-99m, a radioactive tracer used for many medical scans.

\medterm{Molybdosis} Illness caused by excessive molybdenum exposure, which may disturb copper metabolism and cause systemic effects.

\medterm{Mometasone} Mometasone is a potent synthetic glucocorticoid used topically or by inhalation for inflammatory skin or airway disease. It reduces local inflammation but may cause skin atrophy, oral candidiasis, adrenal effects, or other corticosteroid adverse effects.

\medterm{Mometasone Furoate} Mometasone furoate is the furoate ester of the corticosteroid mometasone. Depending on formulation, it is used for inflammatory skin disease, allergic rhinitis, or asthma; local and systemic corticosteroid precautions apply.

\medterm{Momordica} A genus of plants including bitter melon; preparations may affect blood glucose and can interact with medicines.

\medterm{Monascus} Monascus is a genus of fungi used in some fermented foods; some species produce pigments, while contamination can produce harmful mycotoxins.

\medterm{Mondor’S Disease} Superficial thrombophlebitis of the anterior chest wall veins, typically the lateral
thoracic, thoracoepigastric, or superior epigastric veins. Presents as a tender,
cord-like subcutaneous structure with skin retraction. Usually self-limited; associated
with breast surgery, trauma, or hypercoagulable states.

\medterm{Mongolian Spot} A Mongolian spot, now preferably called congenital dermal melanocytosis, is a flat blue-gray patch caused by melanocytes in the dermis, usually present at birth. It is benign and often fades during childhood but should be documented to distinguish it from bruising.

\medterm{Monilethrix} Monilethrix is an inherited hair-shaft disorder in which hair has regularly spaced narrow segments and breaks easily, producing a beaded appearance. It commonly causes short, brittle scalp hair and may involve follicles elsewhere.

\medterm{Monitored Anesthesia Care} Anesthesia care in which a trained professional gives medicines for comfort or pain relief and continuously watches breathing, heart rate, blood pressure, and oxygen. The level of sedation can be adjusted and emergency support is available.

\medterm{Monitoring Device} A monitoring device measures and displays or records a physiologic variable, such as heart rhythm, oxygen saturation, blood pressure, glucose, temperature, or respiratory rate. Accuracy, calibration, and clinical context determine how results should be interpreted.

\medterm{Monitoring Test} A test repeated over time to follow a disease, treatment response, medicine safety, or a possible complication. The result is interpreted with symptoms and earlier measurements.

\medterm{Monkeypox} Monkeypox, now generally called mpox, is a zoonotic infection caused by monkeypox virus, an orthopoxvirus. It can cause fever, lymphadenopathy, and a characteristic rash or lesions and spreads through close contact, including skin-to-skin or respiratory contact.

\medterm{Mono} Informal term; see \textit{Infectious mononucleosis}.

\medterm{Mono} Alternate name; see \textit{Infectious Mononucleosis}.

\medterm{Monoamine Oxidase Inhibitors} A class of antidepressants that slows the breakdown of several brain chemical messengers. They can help selected people with depression but have important interactions with medicines and some foods, requiring careful prescribing.

\medterm{Monobactams} A class of antibiotics used mainly for certain infections caused by oxygen-using Gram-negative bacteria. Aztreonam is the main example; laboratory testing and allergy history help guide its use.

\medterm{Monoblast} An early bone marrow cell that develops into a monocyte, a type of infection-fighting white blood cell.

\medterm{Monoblast Count} A monoblast count is the number or proportion of monoblasts, immature precursors of monocytes, in a blood or bone-marrow sample. Increased blasts may support evaluation for acute myeloid leukemia or another marrow disorder and require expert morphologic or immunophenotypic interpretation.

\medterm{Monochorionic Twin Pregnancy} A twin pregnancy in which identical twins share a single placenta. Because the babies share placental blood vessels, the pregnancy carries unique risks such as twin-to-twin transfusion syndrome and unequal growth, requiring frequent ultrasound monitoring throughout pregnancy. Twins may have separate inner fluid sacs (diamniotic) or rarely share a single sac (monoamniotic).
\synonyms Monochorionic twins; MCDA twins; MCMA twins.

\medterm{Monoclonal} Made by, or arising from, one original cell and therefore very similar to that cell. A monoclonal antibody is designed to recognize one specific target in the body or in a laboratory test.

\medterm{Monoclonal Antibodies} Antibodies produced from a single B-cell clone and directed against one antigenic target; they are used in diagnosis and treatment.

\medterm{Monoclonal Antibody} A laboratory-made immune protein designed to recognize one specific target. Monoclonal antibodies are used in tests and treatments for cancer, inflammatory disease, infection, and other conditions.

\medterm{Monoclonal Gammopathy Of Undetermined Significance} A condition in which abnormal plasma cells make a small amount of one unusual antibody. It usually causes no symptoms or organ damage but can rarely progress to multiple myeloma or another blood cancer, so monitoring is recommended.

\medterm{Monocrotophos} Monocrotophos is a highly toxic organophosphorus insecticide that inhibits acetylcholinesterase. Exposure can cause acute cholinergic poisoning with secretions, bronchospasm, vomiting, bradycardia, weakness, seizures, and respiratory failure.

\medterm{Monocular} Involving or affecting one eye, or occurring with vision from one eye. Monocular vision can reduce depth perception and may result from eye disease, injury, or loss of vision in the other eye.

\medterm{Monocyte} A type of white blood cell that travels in the blood and can enter tissues, where it becomes a macrophage or another immune cell. It helps remove germs and damaged material and coordinates inflammation and repair.

\medterm{Monocytes} Large white blood cells that circulate in the blood and can enter tissues, where they become macrophages or some dendritic cells. They help remove debris and coordinate inflammation and immune defense.

\medterm{Monocytopenia} An abnormally low number of monocytes, a type of white blood cell, in the blood. It may occur with bone-marrow disorders, medicines, infections, or severe illness and is interpreted with the full blood count.

\medterm{Monocytosis} Monocytosis is an increased absolute number of circulating monocytes. It may be reactive to infection or inflammation or associated with a hematologic neoplasm; persistent or marked monocytosis requires appropriate evaluation.

\medterm{Monogenic} Describing a condition or trait caused mainly by a change in one gene. The change may be inherited from a parent or arise for the first time in the affected person.

\medterm{Monogenic Obesity} Severe or early-onset obesity caused mainly by a disease-causing change in one gene. The affected gene may alter appetite or the body's control of energy use.

\medterm{Monograph} A monograph is a detailed, authoritative document devoted to one subject. In health care it may describe a drug’s identity, quality, pharmacology, indications, dosing, safety, or regulatory standards.

\medterm{Monoiodotyrosine} Monoiodotyrosine is an iodinated tyrosine residue formed during thyroid-hormone synthesis. Coupling of iodinated tyrosine residues within thyroglobulin contributes to production of thyroxine and triiodothyronine.

\medterm{Monokines} Monokines are cytokines produced chiefly by monocytes and macrophages, including mediators such as tumor necrosis factor and interleukin-1. They regulate inflammation, immune responses, fever, and communication among cells.

\medterm{Monomer} A monomer is a relatively small molecule that can chemically bond with other monomers to form a polymer. Biological examples include amino acids and nucleotides; some industrial monomers can irritate tissues or be toxic.

\medterm{Monomethylhydrazine} Monomethylhydrazine is a volatile, corrosive, highly toxic chemical used as a rocket propellant. Exposure can injure the eyes, skin, lungs, liver, and nervous system and may cause seizures or systemic poisoning.

\medterm{Mononegavirales} An order of viruses with nonsegmented, negative-sense RNA genomes, including rabies, measles, Ebola, and respiratory syncytial viruses.

\medterm{Mononeuropathy} Mononeuropathy is dysfunction or injury of a single peripheral nerve. Common causes include compression, trauma, ischemia, inflammation, and metabolic disease; symptoms follow the nerve’s distribution and may include pain, numbness, weakness, or autonomic changes.

\medterm{Mononuclear Leukocyte} A white blood cell with one main, unsegmented nucleus. Monocytes and lymphocytes are examples; they help remove damaged material, recognize germs, and coordinate immune responses.

\medterm{Mononucleosis} Common shortened term; see \textit{Infectious mononucleosis}.

\medterm{Mononucleosis Spot Test} A rapid blood test detecting heterophile antibodies associated with infectious mononucleosis; it may be negative early in illness or in young children, so other testing may be needed.

\medterm{Monoplegia} Complete paralysis of one arm or one leg. It can result from a stroke, brain or spinal-cord disease, or severe damage to a nerve. The cause is found through a neurological examination and, when needed, brain or spine imaging and nerve tests.

\medterm{Monopodium} A single foot or one-sided foot structure; the term is uncommon and may be used descriptively in developmental anatomy.

\medterm{Monorchidism} The presence of only one testis, either because the other failed to develop, descended, or was removed.

\medterm{Monosaccharide} The simplest form of sugar, which cannot be split into a smaller carbohydrate. Glucose and fructose are examples; cells use these sugars mainly for energy.

\medterm{Monosodium Glutamate} MSG, the sodium salt of glutamic acid, used as a flavor enhancer because it contributes umami taste; it occurs naturally in foods such as tomatoes and aged cheese.

\medterm{Monosodium Urate} Needle-shaped crystals formed from uric acid that can collect in joints and other tissues when the blood uric-acid level is high. They cause the inflammation of gout and may be identified in joint fluid with a microscope.

\medterm{Monosomy} A genetic condition in which a cell has one copy of a chromosome instead of the usual two. Effects depend on the chromosome and whether all cells or only some cells are affected.

\medterm{Monospot Test} A rapid blood test that detects heterophile antibodies associated with infectious mononucleosis. It may be
negative early in illness or in young children and does not directly detect Epstein-Barr
virus; specific antibody testing may be needed.

\synonyms
Heterophile Antibody Test; Paul-Bunnell Test

\medterm{Monotest} A rapid blood test that looks for heterophile antibodies produced during some cases of infectious mononucleosis. A negative result does not always exclude infection, especially early in illness or in young children.

\medterm{Monounsaturated Fat} A fat whose fatty acids contain one carbon-carbon double bond. It occurs in foods such as olive oil, avocados, nuts, and some seeds.

\medterm{Monozygotic Twin} A monozygotic twin develops when one fertilized egg divides into two embryos. The twins share nearly identical nuclear genomes but can differ through developmental, epigenetic, mitochondrial, and environmental factors.

\medterm{Mons Veneris} The rounded fatty prominence over the pubic bone in front of the external female genitalia.

\medterm{Montelukast} Montelukast is a leukotriene CysLT1-receptor antagonist used for asthma prevention and allergic rhinitis; it is not a rescue bronchodilator, and neuropsychiatric adverse effects require attention.

\medterm{Montelukast Sodium} Montelukast sodium is the sodium salt of montelukast, a leukotriene-receptor antagonist used for asthma prevention and allergic rhinitis. It is not a rescue bronchodilator and carries an important warning for possible neuropsychiatric adverse effects.

\medterm{Montevideo Units} A measurement of the strength of uterine contractions over 10 minutes during labor. It is calculated from pressure readings obtained by a catheter inside the uterus and helps clinicians assess whether contractions are strong enough to support labor progress.

\medterm{Montgomery Tubercles} Small oil-producing glands in the darker skin around the nipple. They often become more noticeable during pregnancy and breastfeeding and release oil that helps protect and lubricate the nipple; they are usually normal.

\medterm{Montreal Cognitive Assessment} A brief screening test of memory, attention, language, problem-solving, visual-spatial skills, and orientation. It can identify possible thinking problems but does not by itself diagnose dementia; education, language, sensory impairment, and acute illness affect the score.

\medterm{Mood} A lasting emotional state experienced by a person and noticed by others. Mood may be described as depressed, anxious, irritable, elevated, or stable and can change with illness, medicines, stress, or life events.

\medterm{Mood Disorder} Alternate name; see \textit{Mood Disorders}.

\medterm{Mood Disorders} A group of mental disorders characterized primarily by clinically significant depression, mania, hypomania, or combinations of these mood episodes. The group includes depressive disorders, bipolar disorders, and cyclothymic disorder.

\medterm{Mood Stabilizers} Medicines used mainly to treat or prevent episodes of mania and depression in bipolar disorder. Examples include lithium, valproate, carbamazepine, and lamotrigine; each has different benefits, interactions, monitoring needs, and risks during pregnancy.

\medterm{Mood Swing} A mood swing is a noticeable change in emotional state or affect over time. Occasional changes can be normal, while frequent, severe, or impairing swings may occur with mood disorders, substance use, medical illness, medications, or sleep disruption.

\medterm{Morale} Morale is the confidence, motivation, and sense of shared purpose within a person or group. It can affect teamwork and occupational well-being but is not itself a specific medical diagnosis.

\medterm{Moraxella} Moraxella is a genus of bacteria that can colonize the respiratory tract and cause ear, sinus, eye, or respiratory infections.

\medterm{Moraxella Atlantae} Moraxella atlantae is a rare Moraxella species; its recovery from a clinical sample should be interpreted with the site and evidence of infection.

\medterm{Moraxella Lincolnii} Moraxella lincolnii is a rare respiratory-tract bacterium; clinical importance depends on the specimen and evidence of infection.

\medterm{Moraxella Oblonga} Moraxella oblonga is a rare Moraxella species found in respiratory or oral specimens; significance depends on clinical evidence of infection.

\medterm{Moraxella Osloensis} An uncommon bacterium that can occasionally cause bloodstream, respiratory, wound, or other opportunistic infections, usually in people with serious illness or weakened immunity.

\medterm{Moraxellaceae} A family of bacteria that includes Moraxella and related organisms. Some normally live on the skin or in the airways, while others can cause ear, sinus, lung, wound, or bloodstream infections.

\medterm{Morbidity} Illness or disability in a person or group of people. It can also
mean how many people have a disease or how much illness or disability a disease causes.

\medterm{Morbilli} The medical or Latin name for measles, an acute contagious viral infection with fever and a characteristic rash.

\medterm{Morbilli Form} Resembling measles in appearance, especially a widespread rash made of flat and slightly raised spots. This description does not prove that a person has measles because medicines, infections, and other illnesses can cause similar rashes.

\medterm{Morbillivirus} A genus of paramyxoviruses that includes measles virus and related animal viruses causing fever, respiratory illness, rash, or neurologic disease.

\medterm{Morbus} A Latin word meaning disease. It appears in some older or formal disease names, but the word itself does not describe the cause, symptoms, or treatment of the condition.

\medterm{Morcellation} The division of tissue into smaller fragments, often during minimally invasive surgery; it may spread an unrecognized cancer and requires appropriate case selection.

\medterm{Mordant} A substance that helps fix a dye to tissue or a material, used in some histological staining methods.

\medterm{Morgagni Hernia} A congenital diaphragmatic hernia occurring through the anterior retrosternal foramen of
Morgagni, resulting from incomplete fusion of parts of the diaphragm. Abdominal contents
may move into the chest and cause respiratory or gastrointestinal symptoms.

\medterm{Morganella} Morganella is a genus of gram-negative bacteria that can cause urinary, wound, bloodstream, and other opportunistic infections.

\medterm{Morganella Morganii} A gram-negative bacterium that can cause urinary, wound, bloodstream, or postoperative infections, especially in hospitalized or medically vulnerable people. Antibiotic choice requires susceptibility testing because resistance is common.

\medterm{Morgellons Disease} Morgellons is a controversial condition involving distressing skin sensations and the belief that fibers or materials emerge from the skin. Symptoms are real to the person and require respectful evaluation for dermatologic, neurologic, psychiatric, toxic, or infectious causes without reinforcing an unproven explanation.

\medterm{Moribund} Critically ill and close to death, with little or no chance of recovery from the underlying condition. A moribund person may be unconscious, have very low blood pressure, cold or mottled skin, and irregular or failing breathing. Care focuses on urgent treatment when possible and comfort when death is inevitable.

\medterm{Morning Sickness} Nausea, with or without vomiting, during pregnancy. It often begins early and improves by the middle of pregnancy, but severe or persistent vomiting can cause dehydration, weight loss, and nutrient problems and may need medical treatment.

\medterm{Moro Reflex} A normal newborn startle response that causes the baby to spread out the arms, open the hands, and then draw the arms inward. It usually fades by about 3 to 6 months; an absent or unequal response needs assessment.

\medterm{Morphea} A localized sclerosing skin disorder that produces areas of thickened, hardened, or discolored skin and differs from systemic sclerosis by usually lacking internal-organ involvement.

\medterm{Morpheus} The Greek and Roman mythological god of dreams; the name is the source of \textit{morphine}, an opioid named for its association with sleep and dreams.

\medterm{Morphinans} Morphinans are a group of chemical compounds with a fused-ring structure related to morphine. The class includes opioid analgesics and other compounds with varying receptor activity, potency, and clinical effects.

\medterm{Morphine} Morphine is a potent opioid analgesic used for selected moderate-to-severe pain and palliative-care indications. It activates opioid receptors and can cause respiratory depression, sedation, constipation, nausea, tolerance, dependence, and overdose.

\medterm{Morphine Overdose} Taking too much morphine can cause extreme sleepiness, pinpoint pupils, slow or stopped breathing, low blood pressure, coma, and death. Call emergency services immediately; naloxone may reverse opioid breathing suppression when available.

\medterm{Morphine Sulfate} Morphine sulfate is a salt formulation of the opioid analgesic morphine. It is available in immediate- and extended-release preparations for selected pain indications and carries risks of respiratory depression, misuse, dependence, and fatal overdose.

\medterm{Morphogenesis} The biological process by which cells and tissues change shape, move, and organize to form organs and body structures during development or tissue repair.

\medterm{Morphologic Response} A morphologic response is a treatment-related change in the visible or microscopic structure of a lesion, tumor, tissue, or cell population. It may include shrinkage, necrosis, differentiation, or altered cellular features and is interpreted using disease-specific criteria.

\medterm{Morphology} The study or description of the form, shape, and structural features of organisms, organs, tissues, cells, or other biological structures.

\medterm{Morquio Syndrome} A rare inherited disorder in which complex sugars build up and interfere mainly with bone and cartilage growth. It can cause short stature, an unusually shaped chest and spine, loose joints, cloudy corneas, and pressure on the spinal cord in the neck. Intelligence is usually unaffected. Treatment may include enzyme replacement, surgery, and careful monitoring of breathing and the spine.

\synonyms
Mucopolysaccharidosis Type Iv; MPS IV

\medterm{Mortality} Death, or the number of deaths in a particular group during a stated period. A mortality measure may describe all causes of death or deaths from a specific disease and population.

\medterm{Mortality Rate} The number of deaths in a specified population and time period, often reported for a particular disease or cause. It helps compare the impact of illness between groups when the populations are defined appropriately.

\medterm{Morton Neuroma} Thickening and irritation of a nerve in the ball of the foot, usually between the third and fourth toes. It can cause burning pain, numbness, or the feeling of a pebble under the foot, often worsened by tight shoes or walking.

\synonyms
Interdigital Neuroma; Morton's Metatarsalgia; Plantar Interdigital Neuroma

\medterm{Morton's Metatarsalgia} Pain in the forefoot beneath the metatarsal heads, caused by pressure or irritation and distinct from Morton neuroma.

\medterm{Mortons Neuroma} Alternate spelling; see \textit{Morton's Neuroma}.

\medterm{Morton’S Neuroma} Alternate spelling; see \textit{Morton neuroma}.

\medterm{Morton's Neuroma} Thickening and irritation of tissue around a nerve, usually between the third and fourth toes. It causes burning forefoot pain or a feeling of a pebble in the shoe and may improve with wider footwear.

\medterm{Morton's Toe} A foot shape in which the second toe is longer than the big toe, often because the first metatarsal is relatively short. It is usually a normal variation but may contribute to pressure or forefoot pain.

\medterm{Morula} A morula is an early embryo consisting of a compact ball of blastomeres formed after several cleavage divisions. It develops into a blastocyst before implantation in the uterine lining.

\medterm{Morvan's Disease} A rare autoimmune neurologic syndrome with severe insomnia, hallucinations, confusion, pain, and peripheral nerve hyperexcitability.

\medterm{Mosaic} Describing a person or tissue composed of two or more genetically different cell populations derived from one zygote, often because of a post-zygotic mutation or chromosomal event; mosaicism can alter the severity and distribution of a genetic disorder.

\medterm{Mosaicism} The presence of groups of cells in one person that have different genetic material, even though all the cells came from the same fertilized egg. It may or may not cause disease, and its effects depend on which cells are affected.

\medterm{Mosquito} A small blood-feeding insect whose bites may transmit malaria, dengue, West Nile virus, and other infections.

\medterm{Mosquito Clamp} A small hemostatic clamp used to compress small blood vessels.

\medterm{Mosquito-Borne} Spread by mosquitoes, usually describing an infection or parasite carried from one host to another by a mosquito bite. Examples include malaria, dengue, West Nile virus, yellow fever, and some forms of encephalitis.

\medterm{Motilin} A hormone made mainly in the small intestine that helps trigger waves of movement through the stomach and small intestine, especially between meals. Its activity contributes to the pattern that clears leftover food and secretions from the upper digestive tract.

\medterm{Motility} The ability of an organ to move its contents or itself. Intestinal motility moves food and waste forward, while abnormal motility can cause pain, constipation, diarrhea, vomiting, or swallowing problems.

\medterm{Motion Analysis} A detailed study of how a person moves during walking, reaching, or other activities. Clinicians use measurements and video to find movement problems, understand their causes, plan treatment, and track improvement.

\medterm{Motion Artifact} Image distortion or blurring caused by movement of the patient, organs, or equipment during image acquisition.

\medterm{Motion Correction} Processing or acquisition methods that reduce image blurring and errors caused by patient, breathing, or organ motion. Examples include image registration, tracking, gating, and navigator methods.

\medterm{Motion Perception} Motion perception is the ability of the visual, vestibular, and proprioceptive systems to detect and interpret movement of the body or surrounding objects. Disturbance may cause vertigo, oscillopsia, imbalance, or visual-motion sensitivity.

\medterm{Motion Sickness} Nausea, dizziness, sweating, or vomiting caused by sensory mismatch during travel or other motion. It can occur in cars,
boats, aircraft, amusement rides, or virtual-reality environments and usually improves when motion stops.
Severe or persistent symptoms require assessment for other causes.

\medterm{Motion Sickness} A disturbance of spatial orientation and equilibrium; see \textit{Kinetosis}.

\medterm{Motivation} Motivation is the set of processes that initiate, direct, and sustain goal-directed behavior. Reduced or excessive motivation may occur with mood disorders, neurologic disease, substance use, medication effects, or other conditions.

\medterm{Motivational Interviewing} A collaborative counseling approach that helps a person explore mixed feelings and strengthen their own reasons for changing a health behavior. The clinician listens without judgment, asks open questions, reflects the person’s views, and supports informed choice.

\medterm{Motor} Relating to movement or to the nerves and brain pathways that control movement. A motor problem may cause weakness, stiffness, poor coordination, tremor, or loss of normal control.

\medterm{Motor Aphasia} Difficulty speaking or writing caused by damage to language areas of the brain, often after a stroke. A person may know what they want to say but struggle to find words, form sentences, or produce speech.

\medterm{Motor Cortex} Areas at the front of the brain that plan and control voluntary movement. Each side mainly controls the opposite side of the body. Injury can cause weakness, loss of fine movement, or difficulty planning purposeful actions.

\medterm{Motor Delay} Motor delay is slower-than-expected acquisition of gross- or fine-motor skills in a child. It may reflect neuromuscular, developmental, genetic, sensory, environmental, or systemic causes and warrants age-appropriate developmental assessment.

\medterm{Motor End-Plate} The specialized area where a nerve signal reaches a skeletal muscle fiber. Chemical signals released there start muscle contraction, and damage can cause weakness, fatigue, or abnormal muscle function.

\medterm{Motor Fluctuations Causes In Parkinson’S Disease} Changes between better and worse movement in Parkinson disease, often as the effect of a medicine wears off before the next dose. Periods of stiffness, slowness, or involuntary movement may occur and can require treatment adjustment.

\medterm{Motor Learning} A relatively lasting change in the ability to perform a motor skill resulting from
practice, experience, feedback, and adaptation. Rehabilitation uses task-specific
repetition and meaningful practice to promote motor learning.

\medterm{Motor Neuron} A nerve cell that carries signals from the brain or spinal cord to skeletal muscle, allowing movement. Damage can cause weakness, muscle wasting, stiffness, cramps, or loss of voluntary control.

\synonyms
Motor Nerve Cell; Motoneuron

\medterm{Motor Neuron Disease} Alternate spelling; see \textit{Motor Neurone Disease}.

\medterm{Motor Neurone Disease} A group of progressive disorders affecting upper and lower motor neurons, including amyotrophic lateral sclerosis. They cause worsening weakness, muscle wasting, stiffness, cramps, and difficulty with movement; speech, swallowing, and breathing may become affected as the disease advances. Cognitive or behavioral changes occur in some forms.

\medterm{Motor Neurones} Nerve cells that carry signals from the brain or spinal cord to skeletal muscles, producing movement. Injury or disease can cause weakness, wasting, stiffness, cramps, poor coordination, or loss of voluntary control.

\medterm{Motor Stereotypy} A repeated, predictable movement such as hand flapping, rocking, or body waving. It may occur during normal development or with autism, intellectual disability, neurologic disease, stress, or other conditions.

\medterm{Motor Tic} A motor tic is a sudden, brief, repetitive movement such as blinking, facial grimacing, or shoulder shrugging. Tics are often temporarily suppressible and may occur in transient tic disorder, persistent tic disorder, or Tourette syndrome.

\medterm{Motor Unit} One movement-controlling nerve cell and all the muscle fibers it activates. Small units allow precise movements, while large units produce more force. Damage to motor units can cause weakness, wasting, cramps, or abnormal results on nerve and muscle tests.

\medterm{Motor Vehicle Accident} A collision or other road-traffic event causing injury or property damage. Medical evaluation may be needed for visible or delayed injuries, including head, neck, chest, abdominal, or musculoskeletal trauma.

\medterm{Mottled Enamel} Patchy, discolored, or structurally altered tooth enamel, classically caused by excessive fluoride during tooth development.

\medterm{Mottled Teeth} Teeth with irregular white, brown, or pitted enamel, often due to dental fluorosis or developmental enamel defects.

\medterm{Moulage} A model, cast, or realistic reproduction of a body part or lesion used for teaching, simulation, or forensic documentation.

\medterm{Mould} British spelling of mold, a multicellular fungus that grows as branching filaments.

\medterm{Mountain Sickness} Illness caused by rapid ascent to high altitude, ranging from acute mountain sickness to life-threatening pulmonary or cerebral edema.

\medterm{Mouth} The oral cavity that serves as the entry point of the digestive tract, used for eating,
breathing, and speaking.

\medterm{Mouth Cancer} Cancer arising in the lips, gums, tongue, cheeks, floor of the mouth, palate, or other oral tissues.

\synonyms
Oral Cancer; Oral Cavity Carcinoma; Oral Squamous Cell Carcinoma

\medterm{Mouth Diseases} Disorders affecting the oral cavity, including infections, ulcers, inflammatory disease, dental disease, mucosal lesions, and tumors.

\medterm{Mouth Neoplasm} A mouth neoplasm is an abnormal growth in the oral cavity that may be benign or malignant; persistent sores, lumps, or patches need assessment.

\medterm{Mouth Sores} Painful ulcers, blisters, raw areas, or inflamed patches inside or around the mouth. Trauma, infection, immune disease, medicines, nutritional deficiency, and cancer can cause them; sores lasting more than two weeks need assessment.

\medterm{Mouth Ulcer} Open sores affecting the oral mucosa, commonly caused by aphthous ulcers, trauma,
infection, irritation, or systemic inflammatory disease. Most resolve within two weeks;
persistent, recurrent, unusually large, or indurated ulcers require evaluation.

\synonyms
Oral Ulcer; Aphthous Ulcer; Canker Sore

\medterm{Mouth-To-Mouth Resuscitation} Rescue breathing in which air is delivered from a rescuer to a person who is not breathing normally, usually as part of CPR.

\medterm{Mouth-Wash} A liquid used to rinse the mouth, freshen breath, or deliver an antiseptic, fluoride, or other active ingredient.

\medterm{Mouthwash} A liquid rinsed around the mouth and then expelled to freshen breath, deliver an active ingredient, or reduce selected oral microorganisms.

\medterm{Mouthwash Overdose} Swallowing too much mouthwash may cause alcohol poisoning, vomiting, drowsiness, low blood sugar, breathing problems, or coma; products containing fluoride or other ingredients can cause additional toxicity. Contact poison control promptly.

\medterm{Movement} Movement is a change in position of the body or a body part produced by coordinated activity of the nervous system, muscles, and joints. Abnormal movement may be excessive, reduced, involuntary, poorly coordinated, or restricted.

\medterm{Movement Disorder} A neurologic condition causing excessive, reduced, or otherwise abnormal movement.
Hyperkinetic disorders include tremor, dystonia, chorea, myoclonus, tics, and
dyskinesia; hypokinetic disorders include bradykinesia and parkinsonism.

\medterm{Movement Disorders} Alternate name; see \textit{Movement Disorders}.

\medterm{Movement Incoordination} Impairment in the ability to execute smooth, coordinated voluntary movements (ataxia), resulting from lesions or dysfunction within the cerebellum, spinocerebellar tracts, vestibular apparatus, or proprioceptive pathways. Manifests as dysmetria, intention tremor, and gait instability.

\synonyms
Ataxic Incoordination; Motor Incoordination; Dysdiadochokinesia

\medterm{Moxalactam} Moxalactam is a beta-lactam antibiotic of the oxacephem group with activity against many gram-negative bacteria and some gram-positive organisms. It is rarely used today and may cause allergic reactions, diarrhea, or bleeding related to impaired coagulation.

\medterm{Moxibustion} A traditional treatment that warms selected skin areas with burning mugwort or another heated material. Evidence of benefit varies by condition, and burns, smoke exposure, allergic reactions, or infection can occur.

\medterm{Moxifloxacin} An antibiotic used for selected bacterial infections. It interferes with bacterial reproduction but can cause serious side effects, including tendon injury, nerve symptoms, blood-sugar changes, or dangerous heart rhythms, so use is selective.

\medterm{Moyamoya Disease} A rare condition in which major arteries supplying the brain gradually narrow or close. Small fragile blood vessels form to compensate, but the condition can cause transient attacks, strokes, headaches, seizures, or thinking problems.

\medterm{Moynihan's Hump} An anatomical variation in which the right hepatic artery forms an abnormally tortuous, caterpillar-like loop that extends close to the gallbladder neck and cystic duct within Calot's triangle. During gallbladder removal (cholecystectomy), this loop can be mistaken for the smaller cystic artery; accidental clipping or injury can cause severe arterial bleeding or damage blood flow to the right lobe of the liver.

\synonyms
Moynihan Caterpillar Turn; Caterpillar Turn of the Right Hepatic Artery; Hepatic Artery Tortuosity

\medterm{MPH} Master of Public Health, a graduate degree focused on population health, disease prevention, health policy, epidemiology, and related public-health practice.

\medterm{MPO} Myeloperoxidase, an enzyme in neutrophil and monocyte granules that uses hydrogen peroxide and halide ions to generate antimicrobial oxidants; MPO antibodies are also used in evaluating some autoimmune vasculitides.

\medterm{Mpox} A viral infection that can cause fever, swollen lymph nodes, tiredness, and a rash or painful lesions on the skin, mouth, or genital area. It spreads mainly through close contact with lesions, body fluids, or contaminated materials.

\medterm{Mpox} A viral disease caused by monkeypox virus, an orthopoxvirus related to smallpox. It usually causes a rash or lesions, fever, swollen lymph nodes, and body aches; spread occurs mainly through close contact, and severity varies with the virus type and the person’s health.

\synonyms
Simian Pox; Monkeypox Virus Infection

\medterm{MPTP} A toxic chemical that can permanently damage movement-controlling nerve cells and produce symptoms resembling Parkinson disease. It is mainly important in research and is not a routine medical treatment.

\medterm{MR Angiography} An MRI scan that creates pictures of blood vessels without using X-rays. It can help assess narrowed or blocked arteries, aneurysms, dissections, clots, and abnormal blood-vessel connections.

\medterm{MR Enterography} A specialized MRI scan of the small intestine after drinking contrast liquid. It can show inflammation, narrowing, fistulas, abscesses, bleeding, or other complications of Crohn disease and does not use X-rays.

\medterm{MR Neurography} Magnetic resonance imaging optimized to show peripheral nerves and abnormalities such as swelling, compression, or injury.

\medterm{MRA} Magnetic resonance angiography, a noninvasive MRI-based technique that produces images of blood vessels and can help evaluate aneurysms, stenoses, malformations, and atherosclerosis.

\medterm{MRCP} Magnetic resonance cholangiopancreatography, an MRI technique that noninvasively images the bile ducts, gallbladder, and pancreatic duct without requiring endoscopic catheterization.

\medterm{MRD} Minimal residual disease: a small number of malignant cells remaining after treatment, detectable by sensitive methods such as flow cytometry or PCR even when routine tests show remission.

\medterm{MRI} Magnetic resonance imaging uses a strong magnetic field, radiofrequency energy, and computer processing to produce detailed images of body tissues without ionizing radiation. Contrast material may be used when clinically appropriate.

\medterm{MRI And Low Back Pain} MRI may be used for low back pain when serious disease is suspected, neurologic deficits are present, or symptoms persist despite appropriate care; many uncomplicated cases do not need early MRI.

\medterm{MRI Brain} An MRI scan used to examine the brain for tumors, stroke, bleeding, infection, inflammation, multiple sclerosis, or degenerative disease. Different settings highlight different tissues and abnormalities; contrast may be used when needed.

\medterm{MRI Contrast} A contrast agent administered during magnetic resonance imaging to improve visibility of selected tissues or abnormalities.

\medterm{MRI Joint} Evaluation of musculoskeletal structures: cartilage, ligaments, tendons, menisci,
labrum. Knee: ACL, meniscus tears. Shoulder: rotator cuff, labrum. Hip: labral tears,
FAI.

\medterm{MRI Sequences} Different image settings used during an MRI scan to highlight anatomy, fluid, fat, bleeding, restricted water movement, or inflammation. Radiologists combine several settings because each shows different types of abnormality.

\medterm{MRI-Targeted Fusion Prostate Biopsy} A prostate biopsy guided by combining a previous MRI with real-time ultrasound. The images help the clinician direct needles toward areas that look suspicious for cancer, sometimes along with routine samples from other areas. The procedure can cause blood in the urine or semen, infection, pain, or temporary difficulty passing urine.

\synonyms
Fusion Biopsy; Mri-Ultrasound Fusion Biopsy

\medterm{Mrna Precursor} An initial RNA transcript that is processed by capping, splicing, and polyadenylation to form mature messenger RNA.

\medterm{MRSA} A strain of Staphylococcus aureus resistant to several common antibiotics. It may live harmlessly on the skin or cause skin, wound, lung, bone, joint, or bloodstream infection; treatment is guided by culture results.

\medterm{MRSA Decolonization} Measures intended to reduce MRSA carriage on the skin or in the nose. Protocols may include topical
agents and antiseptic bathing, but the indication, regimen, and duration depend on the
person's risk, infection history, setting, and local guidance.

\medterm{MRSA Infection} Methicillin-resistant Staphylococcus aureus is a strain of staphylococcus resistant to several beta-lactam antibiotics. It can colonize without illness or cause skin, wound, lung, bone, joint, or bloodstream infection; culture-guided treatment and hygiene reduce spread.

\synonyms
Methicillin-Resistant Staphylococcus Aureus Infection

\medterm{MS Alternative Therapies} Alternate name; see \textit{Alternative Therapy for Multiple Sclerosis}.

\medterm{MSG Symptom Complex} A group of symptoms such as headache, flushing, or palpitations attributed by some people to monosodium glutamate; evidence for a consistent syndrome is limited and symptoms may have other causes.

\medterm{MST} Multiple subpial transection, a neurosurgical procedure in which shallow cuts are made in the cerebral cortex to interrupt seizure pathways when the epileptogenic region cannot safely be removed.

\medterm{MTOR} A protein inside cells that helps control growth, division, energy use, and response to nutrients. Medicines that block mTOR, such as rapamycin-related drugs, are used in some cancers, transplant treatments, and other selected conditions.

\medterm{Mu-Chain Disease} A rare B-cell disorder in which abnormal immunoglobulin heavy chains are produced, sometimes causing anemia, enlarged lymph nodes, or other features of lymphoproliferative disease.

\medterm{Mucin} Mucin is a high-molecular-weight glycoprotein that forms mucus and helps lubricate and protect epithelial surfaces. Abnormal mucin production or accumulation occurs in some respiratory, gastrointestinal, and neoplastic diseases.

\medterm{Mucinoses} A group of disorders in which excessive mucin, a gel-like substance in connective tissue, builds up in the skin or other tissues. Causes may be primary, metabolic, autoimmune, or associated with another disease.

\medterm{Mucinous} Mucinous means containing, producing, or resembling mucin. In pathology, the term describes lesions or tumors with extracellular or intracellular mucin and must be interpreted with the organ, architecture, and other histologic features.

\medterm{Mucinous Adenocarcinoma} Mucinous adenocarcinoma is a gland-forming cancer that produces abundant mucus and can arise in organs such as the colon, pancreas, lung, or breast.

\medterm{Mucinous Borderline Ovarian Tumor} An ovarian mucinous epithelial tumor with cellular proliferation and atypia but without destructive stromal invasion; it may be large and unilateral.

\medterm{Mucinous Cyst Of The Vulva} A benign mucus-filled cyst of the vulva, arising from obstruction or metaplastic mucinous epithelium in a vulvar gland or duct; persistent or atypical masses require examination and possible biopsy.

\medterm{Mucinous Cystadenoma} A mucinous cystadenoma is a usually benign cystic epithelial tumor that produces mucin, commonly arising in the ovary or pancreas. Large lesions may cause pain or pressure, and pathology is needed to distinguish benign from borderline or malignant mucinous tumors.

\medterm{Mucinous Cystadenoma of Pancreas} Alternate name; see \textit{Pancreatic Cysts}.

\medterm{Mucinous Ovarian Carcinoma} A malignant ovarian epithelial tumor composed of invasive mucin-producing cells; metastatic gastrointestinal or pancreatobiliary carcinoma must be excluded.

\medterm{Muckle-Wells Syndrome} An inherited cryopyrin-associated autoinflammatory syndrome causing recurrent fever, rash, joint or muscle pain, and progressive sensorineural hearing loss. Anti-inflammatory treatment can reduce inflammation and prevent complications.

\medterm{Mucocele} A soft, mucus-filled swelling caused by leakage or blockage of a small salivary gland, most often on the inner lower lip after minor injury. It may come and go, but a persistent or recurrent swelling should be examined and may need removal.

\medterm{Mucocutaneous Leishmaniasis} A chronic, mutilating metastatic form of leishmaniasis caused by New World species of the Leishmania subgenus Viannia (principally Leishmania braziliensis), which disseminates from a primary cutaneous ulcer to mucosal tissues through hematogenous or lymphatic routes. It causes progressive, destructive granulomatous ulceration and cartilaginous necrosis of the nasal septum, palate, pharynx, and larynx, clinically presenting with nasal obstruction, epistaxis, septal perforation, and severe facial disfigurement.

\synonyms
Espundia; American Mucocutaneous Leishmaniasis; Nasopharyngeal Leishmaniasis

\medterm{Mucocutaneous Lymph Node Syndrome} Alternate name; see \textit{Kawasaki Disease}.

\medterm{Mucoepidermoid Carcinoma} Mucoepidermoid carcinoma is a cancer containing mucus-producing and squamous-like cells, most often arising in salivary glands or the lungs.

\medterm{Mucoid} Resembling mucus or having a thick, slippery, mucus-like appearance. The word describes a substance or finding and does not by itself identify its cause.

\medterm{Mucolipidoses} A group of inherited lysosomal storage disorders in which cells cannot properly process certain lipids and carbohydrate-containing substances. Material accumulates in tissues and may cause developmental delay, skeletal changes, vision problems, and organ disease.

\medterm{Mucolytic Therapy} Treatment that makes thick mucus easier to move and cough up. Different medicines work by breaking mucus proteins, removing extra DNA, or drawing water into airway secretions; the choice depends on the disease and the person's ability to clear mucus.

\medterm{Mucolytics} Medicines that make thick mucus thinner and easier to clear from the airways. They may help some people cough up mucus, but treatment depends on the cause of congestion and the person's overall condition.

\medterm{Mucopolysaccharides} Long chains of sugar molecules that form part of connective tissue and help support joints, skin, cartilage, and other structures. Inherited problems breaking them down can cause storage diseases.

\medterm{Mucopolysaccharidoses} A group of inherited lysosomal storage diseases in which glycosaminoglycans accumulate because one of several degradation enzymes is deficient. The disorders can affect the skeleton, joints, airways, heart, eyes, abdominal organs, and nervous system, with severity varying by type.

\medterm{Mucopolysaccharidosis} The singular name for any of a group of inherited lysosomal storage disorders in which glycosaminoglycans accumulate because of deficient enzyme activity. Types differ in their genetic cause, age of onset, and effects on the skeleton, organs, airways, eyes, heart, and nervous system.

\medterm{Mucopolysaccharidosis Type I} An inherited lysosomal storage disorder caused by deficient alpha-L-iduronidase, usually because of disease-causing variants in the \textit{IDUA} gene. It ranges from severe Hurler syndrome to milder forms and may affect the skeleton, airways, heart, eyes, abdominal organs, and nervous system.

\synonyms
Hurler Syndrome; Hurler-Scheie Syndrome; Scheie Syndrome

\medterm{Mucopolysaccharidosis Type I} Mucopolysaccharidosis type I is an inherited alpha-L-iduronidase deficiency ranging from severe Hurler syndrome to milder forms, with skeletal, airway, cardiac, eye, and neurologic disease.

\medterm{Mucopolysaccharidosis Type II} Mucopolysaccharidosis type II, or Hunter syndrome, is an X-linked iduronate-2-sulfatase deficiency causing accumulation of glycosaminoglycans and multisystem disease.

\medterm{Mucopolysaccharidosis Type III} Mucopolysaccharidosis type III, or Sanfilippo syndrome, is an inherited disorder of heparan-sulfate breakdown causing progressive neurodevelopmental and behavioral decline.

\medterm{Mucopolysaccharidosis Type IV} Mucopolysaccharidosis type IV, or Morquio syndrome, is an inherited disorder of keratan-sulfate breakdown causing skeletal and connective-tissue abnormalities with variable neurologic involvement.

\medterm{Mucopolysaccharidosis VII} Mucopolysaccharidosis VII, or Sly syndrome, is a lysosomal storage disorder caused by deficient beta-glucuronidase activity. Accumulation of glycosaminoglycans can cause skeletal, connective-tissue, visceral, developmental, and sometimes neurologic abnormalities.

\medterm{Mucopus} A mixture of mucus and pus produced by inflammation or infection. It may appear in the nose, sinuses, airways, bowel, cervix, or other mucus-lined areas and should be interpreted with symptoms.

\medterm{Mucor} Mucor is a genus of molds; related Mucorales fungi can cause mucormycosis, a serious invasive infection mainly in people with severe immune or metabolic problems.

\medterm{Mucor Ramosissimus} A type of mold found in the environment that rarely causes infection in people. When infection occurs, it is usually in someone with severely weakened immune defenses and requires identification of the fungus and specialist treatment.

\medterm{Mucorales} Mucorales is an order of molds that includes fungi capable of causing mucormycosis, a rapidly invasive infection of the sinuses, lungs, skin, or other tissues.

\medterm{Mucormycosis} A rare but rapidly destructive infection caused by certain molds. It most often affects people with severely weakened immunity or uncontrolled diabetes and can damage the sinuses, eyes, lungs, skin, brain, or several organs.

\medterm{Mucosa} Mucosa is the moist epithelial lining of body passages such as the mouth, nose, airways, gastrointestinal tract, and reproductive or urinary tracts.

\medterm{Mucosa-Associated Lymphoid Tissue Lymphoma} Alternate name; see \textit{MALT Lymphoma}.

\medterm{Mucosal} Relating to a mucous membrane, the moist lining of areas such as the mouth, nose, airways, digestive tract, and genital or urinary tract. Mucosa helps protect tissues, produce mucus, absorb substances in some organs, and support local immune defense.

\medterm{Mucosal Immunity} Mucosal immunity is immune defense at surfaces lining the respiratory, gastrointestinal, genitourinary, and other tracts. It includes epithelial barriers, mucus, secretory IgA, antimicrobial molecules, innate cells, and locally regulated adaptive responses.

\medterm{Mucosal Infection} An infection involving a moist epithelial lining such as the mouth, eye, nose, airway, gastrointestinal tract, or genital tract.

\medterm{Mucosal Irritation} Inflammation or injury of moist linings in the eyes, nose, throat, mouth, stomach, or intestines. It may cause burning, pain, redness, swelling, or sores after chemicals, medicines, infection, or friction.

\medterm{Mucosal Melanoma} A melanoma arising from melanocytes in mucosal surfaces such as the sinonasal tract, oral cavity, anorectal region, or genital tract; it is uncommon and often diagnosed at an advanced stage.

\medterm{Mucosal Warts} Warts caused by human papillomavirus on moist lining tissues such as the mouth, genitals, or anus. They may be painless or cause itching, bleeding, or discomfort; some HPV types increase cancer risk and need medical assessment.

\medterm{Mucositis} Painful inflammation and sores in the moist lining of the mouth or digestive tract, often after chemotherapy or radiation. It can make eating and swallowing difficult and may increase the risk of infection.

\medterm{Mucous} Relating to mucus, the viscous fluid produced by mucous membranes and glands that lubricates surfaces and helps trap particles and microorganisms.

\medterm{Mucous Cysts} Small fluid-filled lumps near the end joint of a finger, often associated with osteoarthritis. They may cause tenderness, nail changes, or leakage and should not be punctured without medical advice.
\synonyms
Digital Mucous Cyst; Myxoid Cyst

\medterm{Mucous Membrane} A moist tissue lining body cavities that open to the outside, including the nose, mouth, and lungs.

\synonyms
Mucosa

\medterm{Mucous Membrane Pemphigoid} A chronic autoimmune subepidermal blistering disease affecting mucous membranes, especially the eyes and mouth; scarring can impair vision or function.

\medterm{Mucus} Mucus is a viscoelastic secretion containing water, mucin glycoproteins, salts, antimicrobial factors, and cellular material. It lubricates and protects epithelial surfaces and traps particles and microorganisms for clearance.

\medterm{Mucus Plugging} Accumulation of thick, tenacious mucus within airways causing obstruction. It is common
in asthma, CF, bronchiectasis, and chronic bronchitis. It impairs gas exchange, promotes
infection, and causes atelectasis. Treatment includes airway clearance techniques,
mucolytics, and bronchodilators.

\medterm{Muehrcke's Nails} Paired transverse white bands in the nail beds, usually associated with hypoalbuminemia or other systemic illness. The bands remain fixed as the nail grows because they arise from the nail bed rather than the nail plate.

\medterm{Muenke Syndrome} A craniosynostosis syndrome caused by a pathogenic FGFR3 variant. Premature fusion of
one or more skull sutures may produce an abnormal head shape and can be accompanied by
developmental, hearing, limb, or neurological findings.

\medterm{MUGA Scan} A heart scan that uses a small radioactive tracer and records the heart's pumping with each heartbeat. It measures how much blood the main chambers pump and can assess wall movement.

\synonyms
Radionuclide Angiography; Radionuclide Ventriculography

\medterm{Mulberry Molars} Permanent back teeth with many small, rounded, irregular cusps. They are classically linked with congenital syphilis but can also result from other developmental problems affecting tooth formation before birth.

\medterm{Mulibrey Nanism} Mulibrey nanism is a rare inherited multisystem disorder, usually caused by biallelic variants in \textit{TRIM37}. It may cause severe growth restriction, characteristic facial features, pericardial or hepatic disease, insulin resistance, and other organ involvement.

\medterm{Muller Cells} Supporting glial cells of the retina that help nourish retinal nerve cells and maintain the retina's structure and chemical environment.

\medterm{Multicenter Study} A multicenter study is a clinical or research investigation conducted at two or more sites under a coordinated protocol. It can improve enrollment and generalizability but requires harmonized methods, quality control, and site-specific oversight.

\medterm{Multicentric} Having more than one center or focus; in oncology, describing separate primary tumor foci, often also called multifocal or polycentric.

\medterm{Multicystic Dysplastic Kidney} A birth defect in which one kidney develops as a cluster of cysts and has little or no working tissue. The other kidney often provides most kidney function but may also need monitoring for related urinary problems.

\medterm{Multidetector Computed Tomography} A CT scan that uses many detectors to collect multiple thin image slices at the same time. It produces detailed pictures quickly and is useful for injuries, blood vessels, organs, and some virtual examinations, but uses X-rays.

\medterm{Multidisciplinary Disaster Victim Identification} The coordinated process of identifying victims of mass-fatality incidents (natural
disasters, transport crashes, terrorist attacks) using multiple disciplines:
fingerprinting, dental records, DNA profiling, forensic anthropology/radiology, and
personal effects, with ante-mortem data collection and comparison.

\medterm{Multidisciplinary Team} Healthcare professionals from different fields who work together on one person's care. Depending on the problem, the team may include doctors, nurses, therapists, pharmacists, psychologists, dietitians, and social workers.

\medterm{Multidrug-Resistant TB} Tuberculosis caused by bacteria resistant to at least the two most important first-line medicines, isoniazid and rifampin. It needs specialized medicines, susceptibility testing, and close monitoring because treatment is harder and may last many months.

\medterm{Multienzyme Complexes} Groups of enzymes organized together so that several steps in a chemical pathway occur efficiently. This arrangement can speed reactions, limit loss of intermediates, and coordinate metabolism inside cells.

\medterm{Multifactorial} Caused or influenced by several factors rather than one single cause. These may include genes, age, lifestyle, infections, medicines, environment, and other health conditions, with each factor adding to the overall risk.

\medterm{Multifocal} Arising from or involving two or more distinct foci or sites, such as separate tumor foci or vascular lesions.

\medterm{Multifocal Atrial Tachycardia} An atrial tachyarrhythmia with at least three distinct P-wave morphologies, irregular
rhythm, and a rate above 100 beats per minute. It is strongly associated with severe
pulmonary disease, hypoxia, and metabolic stress.

\medterm{Multigravida} A person who has been pregnant two or more times, including pregnancies that did not continue.

\medterm{Multi-Infarct Dementia} A decline in memory and thinking caused by several strokes that damage different parts of the brain. Problems may include slowed thinking, poor planning, mood changes, walking difficulty, or memory loss; the term is now often included under vascular cognitive impairment.

\synonyms
Vascular Dementia; Vascular Cognitive Impairment

\medterm{Multilocular} Having several separate spaces or compartments. The word may describe a cyst, tumor, or other lesion seen on imaging; its significance depends on the organ involved, the appearance of the walls, and the person's symptoms.

\medterm{Multi-Minicore Disease} A rare inherited muscle disorder in which muscle fibers contain multiple small areas of damage. It can cause low muscle tone, weakness, spinal curvature, and breathing problems.
\synonyms
MmD; Minicore Myopathy; Multicore Disease

\medterm{Multimodal Analgesia} Use of analgesic techniques with different mechanisms to improve pain control while
reducing reliance on any single drug, especially opioids. Components may include
acetaminophen, nonsteroidal anti-inflammatory drugs, regional anesthesia, local
anesthetics, and selected adjuvants.

\medterm{Multimodal Imaging} Use of two or more imaging methods, such as CT, MRI, ultrasound, or nuclear imaging, to examine the same disease or body region. Combining methods can provide complementary structural, functional, or molecular information.

\medterm{Multimorbidity} The coexistence of two or more chronic health conditions in the same person, without
requiring one to be considered the primary disease. It complicates treatment, increases
polypharmacy risk, and requires coordinated patient-centered care.

\medterm{Multinodular Goiter} An enlarged thyroid containing multiple nodules. Thyroid function may be normal, underactive, or overactive; evaluation commonly includes thyroid-function tests and ultrasound, with further assessment of suspicious nodules or compressive symptoms.

\medterm{Multinucleated Hepatocyte} A multinucleated hepatocyte is a liver cell containing two or more nuclei, usually from nuclear division without complete cytokinesis. It may be a normal or reactive finding, but its significance depends on the number, morphology, and associated liver abnormalities.

\medterm{Multipara} A person who has had two or more pregnancies that reached the stage of potential fetal viability, regardless of whether the pregnancies resulted in live births.

\medterm{Multiparametric MRI Of The Prostate} A detailed prostate MRI that combines several types of images to show prostate structure, water movement, and blood flow. It helps identify suspicious areas and guide biopsy or treatment, but an abnormal scan does not by itself prove cancer.

\medterm{Multiparous} Having given birth two or more times to a fetus or fetuses beyond the stage of viability.

\medterm{Multiple Bacterial Abscesses} Several pockets of pus caused by bacterial infection in one or more organs. They may develop when bacteria spread through the blood or extend from nearby tissue and can cause fever, pain, or organ problems; imaging helps locate them.

\medterm{Multiple Biliary Hamartomas} Multiple small, harmless birth-related changes in the liver's bile ducts that look like tiny cysts on scans. They usually cause no symptoms and are often found by chance, but imaging helps distinguish them from cancer or other liver disease.

\medterm{Multiple Endocrine Neoplasia} An inherited condition causing tumors or overactivity in two or more hormone-making glands during a person's lifetime.

\medterm{Multiple Endocrine Neoplasia I} MEN type 1 is an inherited tumor syndrome involving parathyroid, pituitary, and pancreatic or duodenal neuroendocrine tumors, often caused by MEN1 variants.

\medterm{Multiple Endocrine Neoplasia II} MEN type 2 is an inherited syndrome with medullary thyroid cancer, pheochromocytoma, and sometimes parathyroid disease, commonly involving RET variants.

\medterm{Multiple Endocrine Neoplasia Type 1} An inherited condition that causes tumors in several hormone-making organs, especially the parathyroid glands, pituitary gland, and pancreas or nearby intestine. It is usually caused by a change in the MEN1 gene and requires lifelong monitoring.

\medterm{Multiple Endocrine Neoplasia Type 2} An inherited condition that greatly increases the risk of medullary thyroid cancer and may cause adrenal or parathyroid tumors. Type 2B can also cause nerve-related mouth bumps and distinctive body features; lifelong specialist monitoring is needed.

\medterm{Multiple Endocrine Neoplasia Type 2A} An inherited condition that greatly increases the risk of medullary thyroid cancer and can also cause adrenal tumors and overactive parathyroid glands. Regular genetic and specialist monitoring helps detect these problems early.

\synonyms
MEN 2A; MEN2A; Sipple Syndrome

\medterm{Multiple Endocrine Neoplasia Type 2B} An inherited RET-associated syndrome characterized by medullary thyroid
carcinoma, pheochromocytoma, mucosal neuromas, and a marfanoid habitus. The thyroid-cancer risk is high and requires
early specialist genetic and surgical assessment.

\medterm{Multiple Endocrine Neoplasia Type 4} A rare, autosomal dominant inherited tumor syndrome caused by germline mutations in the \textit{CDKN1B} gene, which encodes the p27$^{Kip1}$ tumor suppressor protein. It presents with features overlapping with MEN1, most commonly characterized by primary hyperparathyroidism (parathyroid adenomas causing high blood calcium) and anterior pituitary adenomas (such as prolactinomas or growth hormone-secreting tumors). Patients may also develop pancreatic neuroendocrine tumors, adrenal adenomas, and reproductive organ tumors. It is typically diagnosed in patients with MEN1-like symptoms who test negative for mutations in the \textit{MEN1} gene.

\synonyms
MEN 4; MEN4; Multiple Endocrine Neoplasia Type IV; CDKN1B-Related MEN

\medterm{Multiple Exostoses} Multiple exostoses, or hereditary multiple osteochondromas, is an inherited disorder causing multiple cartilage-capped bony projections near growth plates. Lesions can deform bones, restrict motion, compress nerves or vessels, and rarely undergo malignant transformation.

\medterm{Multiple Fibrofolliculomas} Numerous harmless small skin-colored growths arising from hair follicles, usually on the face and trunk. They can be a sign of Birt-Hogg-Dubé syndrome, which also increases the risk of kidney tumors and collapsed lung.

\medterm{Multiple Inherited Schwannomas} Alternate name; see \textit{Neurofibromatosis Type 2}.

\medterm{Multiple Mononeuropathy} Damage to several separate nerves in different parts of the body, often developing unevenly. It can cause patchy pain, numbness, weakness, or muscle wasting and may result from blood-vessel inflammation, diabetes, infection, or another systemic disease.

\medterm{Multiple Myeloma} A cancer of plasma cells that usually develops in the bone marrow and often produces a monoclonal immunoglobulin. It can
cause bone lesions, anemia, kidney dysfunction, high calcium, infection, or other organ
damage.

\synonyms
Plasma Cell Myeloma; Kahler Disease; Myelomatosis

\medterm{Multiple Myeloma: Deadly To Chronic} Multiple myeloma is a plasma-cell cancer that produces a monoclonal protein and can cause bone lesions, anemia, kidney injury, high calcium, infections, and fatigue. It is often treatable as a chronic illness, though relapse and complications remain possible.

\medterm{Multiple Organ Dysfunction Syndrome} Abnormal function of two or more organs during severe illness, such as sepsis, shock, trauma, or pancreatitis. Organ function may worsen over time and requires close monitoring and treatment of the underlying cause.

\medterm{Multiple Organ Failure} A life-threatening condition in which two or more organs stop working properly, usually during severe infection, shock, injury, or widespread inflammation. Intensive care supports the organs while doctors treat the underlying cause.

\medterm{Multiple Pregnancy} A pregnancy involving two or more fetuses. Twins may develop from two separate eggs or from one fertilized egg that divides; triplets and higher-order pregnancies are less common. Multiple pregnancies have higher risks of premature birth, growth problems, and pregnancy complications and require closer monitoring.

\medterm{Multiple Sclerosis} A chronic immune-mediated disease of the brain, spinal cord, and optic nerves. The immune system mistakenly attacks myelin, the protective covering around nerve fibers, and may also damage the nerve fibers themselves. This creates areas of inflammation and scarring, called lesions, that slow or block nerve signals. Symptoms vary and may include vision changes, numbness, weakness, stiffness, poor coordination, balance problems, fatigue, thinking difficulties, or bladder and bowel problems. Some people have attacks followed by partial or complete recovery, while others develop steadily worsening symptoms; mixed patterns also occur.

\synonyms
MS; Disseminated Sclerosis

\medterm{Multiple Sclerosis MS And Pregnancy} Multiple sclerosis usually does not prevent pregnancy, but symptoms, disease-modifying medicines, relapse risk, delivery planning, and postpartum care require coordination. Some treatments must be stopped or changed, and relapse risk can shift after delivery.

\medterm{Multiple Sclerosis MS Early Signs And Types} Multiple sclerosis is an immune-mediated disease of the central nervous system with episodes or progression of visual, sensory, motor, balance, bladder, cognitive, or fatigue symptoms. Relapsing, progressive, and other patterns guide testing and treatment.

\medterm{Multiple Sleep Latency Test} A daytime sleep study measuring how quickly a person falls asleep during several scheduled nap opportunities. It helps evaluate excessive daytime sleepiness and narcolepsy after nighttime sleep has been assessed.

\medterm{Multiple System Atrophy} A progressive brain disorder affecting movement and automatic body functions. It may cause stiffness and slowness, poor balance, fainting from low blood pressure, bladder problems, speech changes, or difficulty swallowing.

\synonyms
Shy-Drager Syndrome; MSA

\medterm{Multiple System Atrophy Cerebellar Type} An adult-onset neurodegenerative synucleinopathy (MSA-C) characterized primarily by progressive cerebellar ataxia (gait instability, limb dysmetria, dysarthria, oculomotor abnormalities) accompanied by severe autonomic failure and variable pyramidal signs.

\synonyms
MSA-C; Cerebellar Multiple System Atrophy; Olivopontocerebellar Atrophy Variant

\medterm{Multiple System Atrophy Parkinsonian Type} A progressive neurodegenerative disorder (MSA-P) characterized predominantly by parkinsonian features (rigidity, akinesia, postural instability) that respond poorly to levodopa, combined with early severe autonomic failure (orthostatic hypotension, urinary incontinence) and variable pyramidal dysfunction.

\synonyms
MSA-P; Parkinsonian Multiple System Atrophy; Striatonigral Degeneration Variant

\medterm{Multiple System Failure} Failure of two or more organ systems, also called multiple-organ failure or multiple-organ dysfunction syndrome; it may occur during severe illness such as sepsis, shock, or major trauma.

\medterm{Multiple Vitamin Overdose} Taking too many multivitamins can cause poisoning, especially from iron, vitamin A, or vitamin D. Symptoms may include vomiting, abdominal pain, confusion, liver injury, kidney injury, or dangerous calcium levels; urgent advice is needed.

\medterm{Multiplex PCR Panel} A laboratory test that looks for genetic material from several germs in one sample. Different panels test respiratory, intestinal, blood, or other specimens; a detected germ may represent infection, harmless carriage, or material left after an infection.

\synonyms
Multiplex Polymerase Chain Reaction; Syndromic Molecular Panel

\medterm{Multiplicity} Multiplicity is the state of having more than one occurrence, copy, lesion, tumor, or pregnancy component. Its interpretation depends on context, such as multiple gestation, multifocal disease, or copy number in genetics.

\medterm{Multipolar Neuron} A nerve cell with one long output branch and several receiving branches. This common type of neuron includes many cells that control movement or pass signals within the brain and spinal cord.

\medterm{Multisystem Inflammatory Syndrome In Children} A serious delayed inflammatory illness occurring in some children after infection with SARS-CoV-2, involving fever and two or more organ systems and potentially causing shock or cardiac dysfunction.

\medterm{Multivesicular Body} A multivesicular body is an endosomal compartment containing many small intraluminal vesicles. It can fuse with the plasma membrane to release exosomes or with a lysosome to degrade its cargo.

\medterm{Mumps} A contagious viral infection caused by a paramyxovirus and spread through respiratory droplets, saliva, and direct contact. It typically begins with fever, headache, and fatigue, followed by tender swelling of one or both parotid salivary glands (parotitis) below and in front of the ears, making chewing and swallowing painful. Complications are more common in adolescents and adults and include orchitis (painful testicular inflammation that can cause atrophy and reduced fertility), viral meningitis, pancreatitis, and permanent sensorineural hearing loss. Treatment is supportive with rest, fluids, and pain relief. It is highly preventable with the standard two-dose Measles-Mumps-Rubella (MMR) vaccine.
\synonyms{Epidemic Parotitis; Mumps Parotitis; Paramyxoviral Parotitis}

\medterm{Mumps Virus} The virus that causes mumps, a contagious infection spread through saliva and respiratory droplets. It commonly causes fever and painful swelling near the jaw, and can sometimes inflame the brain, pancreas, or reproductive organs; vaccination reduces risk.

\medterm{Munchausen Syndrome} Alternate name; see \textit{Factitious Disorder}.

\medterm{Mupirocin} Mupirocin is a topical antibiotic that inhibits bacterial isoleucyl-tRNA synthetase. It is used for selected skin infections and nasal eradication of certain staphylococci; resistance can develop with inappropriate or prolonged use.

\medterm{Mural} Relating to a wall, especially the wall of a body cavity, organ, or blood vessel.

\medterm{Mural Thrombosis} Formation of a blood clot along the wall of a heart chamber or blood vessel rather than freely within the lumen.

\medterm{Mural Thrombus} A mural thrombus is a blood clot that forms on the wall of a heart chamber or large
blood vessel, typically in areas of damaged endothelium, stasis, or turbulent flow.

\medterm{Muramidase} Muramidase, or lysozyme, is an enzyme that cleaves the peptidoglycan backbone of bacterial cell walls. It is present in secretions and innate immune cells and contributes to defense against susceptible bacteria.

\medterm{Murine Typhus} An acute endemic rickettsial illness caused by Rickettsia typhi and transmitted to humans primarily through the feces of infected rodent fleas (Xenopsylla cheopis). It typically presents with high prolonged fever, severe frontal headache, myalgias, arthralgias, and a maculopapular rash that typically spares the face, palms, and soles.

\synonyms
Endemic Typhus; Flea-Borne Typhus

\medterm{Murmur} An extra sound heard during a heartbeat, usually caused by fast or turbulent blood flow through the
heart or nearby blood vessels. Some murmurs are harmless, while others may be caused by heart valve or
structural problems.

\medterm{Murmur Classification} Heart murmurs are classified by timing, grade, shape, location, radiation, and pitch.
Timing: systolic (between S1 and S2), diastolic (between S2 and S1), or continuous.
Grade: I (barely audible) to VI (heard without stethoscope). Shape: crescendo,
decrescendo, or plateau.

\medterm{Murphy Sign} A sudden catch or stop in breathing when a doctor presses under the right ribcage as the patient takes a deep breath in, indicating an inflamed gallbladder.

\medterm{Murphy's Sign} Inspiratory arrest caused by pain when the right upper quadrant is palpated, traditionally suggesting acute cholecystitis.

\medterm{Musca} A fly genus including the housefly; some species can spread germs by carrying them on their bodies.

\medterm{Muscae Volitantes} Moving spots, threads, or cobweb-like shapes seen in the visual field, usually caused by shadows from vitreous opacities.

\medterm{Muscarine} Muscarine is a mushroom-derived agonist of muscarinic acetylcholine receptors. Ingestion of muscarine-containing mushrooms can cause a cholinergic syndrome with salivation, sweating, tearing, abdominal cramps, vomiting, diarrhea, bronchospasm, and bradycardia.

\medterm{Muscarinic Antagonist} A muscarinic antagonist blocks acetylcholine at muscarinic receptors in smooth muscle, glands, the heart, and the central nervous system. Examples include atropine, ipratropium, and oxybutynin; effects can include dry mouth, blurred vision, urinary retention, tachycardia, and confusion.

\medterm{Muscidae} Muscidae is a family of flies that includes the housefly. Some species can mechanically carry pathogens to food or wounds, and their larvae may occasionally cause myiasis.

\medterm{Muscimol} Muscimol is a psychoactive compound found in some Amanita mushrooms and an agonist at GABA-A receptors. Ingestion can cause altered perception, sedation, agitation, incoordination, or seizures and may require toxicologic evaluation.

\medterm{Muscle} Body tissue that shortens to create force and movement. Skeletal muscle moves the bones and supports posture, heart muscle pumps blood, and smooth muscle moves substances through organs such as the bowel and blood vessels.

\medterm{Muscle Aches} Pain or soreness in one or more muscles. Common causes include overuse, injury, infection, medicines, and conditions affecting the whole body.

\medterm{Muscle Atrophy} Loss or wasting of muscle tissue caused by disuse, nerve injury, inflammation, malnutrition, or systemic disease.
It may lead to weakness and reduced function; treatment focuses on the cause, nutrition, and
appropriately supervised rehabilitation.

\medterm{Muscle Biopsy} Removal of a small sample of skeletal muscle for microscopic, histochemical, immunologic, or molecular examination in suspected myopathy, inflammation, infection, or infiltrative disease.

\medterm{Muscle Cell} A cell specialized to contract; skeletal, cardiac, and smooth muscle cells perform different movement and organ functions.

\medterm{Muscle Contraction} The tightening of a muscle, allowing it to produce force or movement. Muscle fibers use calcium and energy from ATP to make their protein filaments slide against each other.

\medterm{Muscle Contracture} Persistent shortening or tightening of a muscle or nearby tissues that limits joint movement and does not fully relax with stretching. It may follow scarring, prolonged spasticity, immobility, injury, or, more rarely, a muscle-energy disorder.

\synonyms
Skeletal Muscle Contracture; Metabolic Contracture

\medterm{Muscle Cramp} A muscle cramp is a sudden, involuntary, painful contraction of a muscle. It may follow exertion, dehydration, electrolyte disturbance, pregnancy, medication use, nerve disease, or vascular disease; recurrent cramps need assessment of the cause.

\synonyms
Charley Horse; Involuntary Muscle Spasm

\medterm{Muscle Development} Muscle development is the formation, growth, differentiation, and maturation of muscle tissue. It depends on genetic programs, neural input, hormones, nutrition, mechanical loading, and repair processes.

\medterm{Muscle Disorder} A disease affecting muscle structure or function, causing weakness, pain, cramps, stiffness, wasting, or abnormal movement. Causes include inherited conditions, inflammation, nerve disease, medicines, infection, and metabolic problems.

\medterm{Muscle Fatigue} A reduced ability of a muscle to produce force after prolonged or repeated activity. It may result from exertion, poor sleep, illness, medicines, nerve problems, or muscle disease, and improves when the cause is treated.

\medterm{Muscle Fiber} A muscle fiber is an individual muscle cell containing contractile myofibrils. Skeletal muscle fibers vary in size and metabolic properties, while cardiac and smooth muscle cells have distinct structures and functions.

\medterm{Muscle Fibers} Individual muscle cells containing contractile proteins. Skeletal muscle fibers vary in size and properties, including endurance-oriented slow fibers and powerful, faster-fatiguing fibers.

\synonyms
Muscle Cells; Myocytes

\medterm{Muscle Function Loss} Reduced ability of a muscle or muscle group to generate movement or force, caused by muscle disease, nerve injury, neuromuscular-junction disorders, pain, or disuse.

\medterm{Muscle Hypertrophy} Muscle hypertrophy is enlargement of individual muscle fibers, usually in response to resistance, hormonal stimulation, or compensatory workload. It can be physiologic or associated with disease, depending on the muscle and cause.

\medterm{Muscle Neoplasm} A muscle neoplasm is an abnormal growth in skeletal or smooth muscle or nearby soft tissue and may be benign or malignant.

\medterm{Muscle Pain} Pain, soreness, or aching in one or more muscles caused by exertion, injury, infection, inflammation, medicines, or an underlying medical condition. Persistent, severe, or weakness-associated pain needs assessment.

\synonyms
Myalgia; Muscle Ache; Myodynia

\medterm{Muscle Protein} A protein found in muscle that provides structure, produces contraction, stores or transports substances, or regulates muscle activity. Important examples include actin, myosin, troponin, dystrophin, and myoglobin.

\medterm{Muscle Relaxant} A muscle relaxant is a medicine that reduces skeletal-muscle spasm or contraction. Centrally acting agents reduce muscle spasm through the nervous system, whereas neuromuscular blockers interrupt transmission at the neuromuscular junction and require respiratory support.

\medterm{Muscle Relaxants} Medicines that reduce skeletal-muscle spasm or neuromuscular transmission; they include centrally acting drugs and peripheral neuromuscular blockers.

\medterm{Muscle Rigidity} Muscle rigidity is abnormal resistance to passive movement caused by sustained muscle contraction or altered tone. It occurs in conditions such as Parkinson disease, dystonia, meningitis, neuroleptic malignant syndrome, and tetanus.

\medterm{Muscle Spasm} A sudden, involuntary tightening of a muscle that may cause pain, stiffness, or an abnormal position. Spasms can follow overuse, dehydration, mineral changes, nerve problems, medicines, pregnancy, or injury.

\medterm{Muscle Spasms} Muscle spasms are sudden involuntary contractions caused by exertion, dehydration, electrolyte disturbance, nerve disease, medicines, pregnancy, or musculoskeletal problems. Recurrent, widespread, painful, or weakness-associated spasms need evaluation for an underlying cause.

\medterm{Muscle Spindle} A small sensor inside skeletal muscle that detects how much the muscle is stretched and how quickly it changes length. It helps maintain posture, muscle tone, balance, and coordinated movement.

\medterm{Muscle Strain Treatment} Care for a stretched or torn muscle, usually involving temporary activity reduction, gradual movement, rehabilitation, and pain relief chosen safely for the person. Severe swelling, weakness, deformity, inability to use the limb, or persistent pain needs assessment.

\medterm{Muscle Strength} Muscle strength is the maximal force a muscle or muscle group can generate under specified conditions. It is assessed by functional or resisted testing and may be reduced by muscle, nerve, neuromuscular-junction, central, or systemic disorders.

\medterm{Muscle Tissue} Specialized tissue that contracts to produce movement, maintain posture, move substances through organs, and generate heat. Skeletal, heart, and smooth muscle have different structures and control systems for their jobs.

\medterm{Muscle Tonus} The slight, continuous activity present in resting muscles, which gives them resistance to being stretched. Too much or too little tone can occur with disorders of the brain, spinal cord, nerves, or muscles.

\medterm{Muscle Twitching} Brief involuntary contractions of a small group of muscle fibers, often called fasciculations; persistent or widespread twitching may require evaluation.

\medterm{Muscle Weakness} Reduced ability to generate force with one or more muscles. Causes include nerve or muscle disease, pain, inactivity,
medicines, metabolic disorders, and aging.

\medterm{Muscle-Contraction Headache} A pressure-like headache caused by prolonged tightening or sensitivity of muscles in the scalp, neck, or jaw. Stress, poor posture, fatigue, or jaw clenching may contribute; frequent, severe, or changing headaches need medical assessment.
\synonyms
Tension-Type Headache; Stress Headache

\medterm{Muscle-Invasive Bladder Cancer} Bladder cancer that has grown into the bladder muscle layer. It has a greater risk of spreading than cancer limited to the bladder lining and is classified by examination of bladder tissue.

\medterm{Muscular} Relating to muscle or made of muscle tissue. It may describe strength, body build, pain, weakness, or a structure containing muscle.

\medterm{Muscular Dystrophy} A group of more than 30 inherited muscle diseases that cause progressive weakness and breakdown of skeletal muscles. Different gene changes affect proteins that strengthen and protect muscle fibers, so the age of onset, muscles affected, and severity vary by type. Weakness may cause difficulty walking or doing everyday activities. Some forms also affect the heart, breathing muscles, swallowing, bones, brain, or other organs. Muscular dystrophy is not contagious and is not caused by ordinary injury or physical activity.

\synonyms
MD; Inherited Myopathy

\medterm{Muscular Endurance} The ability of a muscle or muscle group to keep working during repeated or prolonged effort. It depends on muscle strength, energy supply, blood flow, nerve control, conditioning, and overall health.

\medterm{Muscularis} A layer of smooth muscle in the wall of a hollow organ, especially the digestive tract. Its coordinated contractions move contents through the organ and can also change the diameter or pressure of the passage.
\synonyms
Muscular Layer; Tunica Muscularis

\medterm{Musculocutaneous Nerve} A nerve from the neck and upper chest that controls several muscles used to bend the elbow and turn the forearm. It also carries sensation from part of the outer forearm; injury can cause weakness or numbness there.

\medterm{Musculoskeletal} Relating to muscles, bones, joints, tendons, ligaments, and related tissues.

\medterm{Musculoskeletal Deformity} A musculoskeletal deformity is an abnormal shape, alignment, or structure of bone, joint, muscle, or related tissue. It may be congenital, developmental, traumatic, degenerative, neuromuscular, or caused by disease.

\medterm{Musculoskeletal Disease} A disorder affecting bones, joints, muscles, tendons, ligaments, cartilage, or other supporting tissues, causing pain or limited movement.

\medterm{Musculoskeletal Injury} Trauma to muscles, bones, joints, tendons, ligaments, cartilage, or related soft tissues, ranging from strains and sprains to fractures and dislocations.

\medterm{Musculoskeletal Pain} Pain arising from muscles, bones, joints, ligaments, tendons, or related soft tissues. Common causes
include injury, overuse, arthritis, and inflammation; persistent pain, weakness, swelling, or
neurologic symptoms should be assessed.

\medterm{Musculoskeletal System} The musculoskeletal system consists of bones, joints, cartilage, ligaments, tendons, muscles, and associated connective tissues. It provides support, protects organs, enables movement, and contributes to mineral storage and blood-cell production.

\medterm{Musculoskeletal Tumor} An abnormal growth arising in bone, muscle, cartilage, connective tissue, or related structures. It may be benign or malignant and can cause a lump, pain, swelling, reduced function, or fracture.

\medterm{Musculus} The Latin word for muscle, used in formal anatomical names. For example, ``musculus deltoideus'' means the deltoid muscle. Abbreviations may be used in anatomy texts to identify one or several muscles.

\medterm{MUSE} A medicated urethral system that delivers a pellet into the urethra to treat erectile dysfunction.
\synonyms
Medicated Urethral System for Erection

\medterm{Mushroom} The visible fruiting body of a fungus. Edible mushrooms can provide nutrients, but wild mushrooms may be poisonous or confused with safe species; uncertain mushrooms should not be eaten.

\medterm{Mushroom Poisoning} Illness caused by eating toxic fungi, with effects ranging from gastrointestinal upset to liver failure, neurologic toxicity, or death.

\medterm{Music Therapy} The planned use of music by a trained therapist to support health and wellbeing. It may help with pain, stress, communication, movement, memory, emotional expression, or coping during illness and rehabilitation.

\medterm{MuSK-Positive Myasthenia Gravis} A form of myasthenia gravis in which antibodies interfere with nerve-to-muscle communication. It often causes weakness of the face, mouth, throat, neck, or breathing muscles and requires specialist care.

\synonyms
MuSK Myasthenia Gravis; MuSK-Ab-Positive MG

\medterm{Mustard Compounds} Mustard compounds are chemical agents containing reactive sulfur or nitrogen mustard groups that can alkylate DNA and other molecules. Some are chemotherapy drugs; others are highly toxic chemical-warfare agents.

\medterm{Mustard Gas} Mustard gas, or sulfur mustard, is a vesicant chemical-warfare agent that damages skin, eyes, and respiratory tissues and can cause blistering, airway injury, and delayed complications. It is not a true gas at ordinary temperatures and has no specific widely available antidote.

\medterm{Mutagen} A chemical, physical agent, or biological factor capable of causing a genetic mutation. Mutagenicity does not by itself prove that an agent causes cancer.

\medterm{Mutant} A mutant is a cell, organism, or genetic sequence carrying a mutation relative to a reference or parental form. The change may have no effect, alter a trait, or cause disease, depending on its location and functional consequence.

\medterm{Mutant Proteins} Proteins whose amino-acid sequence differs because of a genetic mutation. The change may have no effect or may alter folding, stability, location, or function and contribute to inherited disease, cancer, or drug resistance.

\medterm{Mutate} To mutate means to undergo a change in genetic sequence. Mutations arise spontaneously or from replication errors, DNA damage, or mutagen exposure and may be inherited, acquired in somatic cells, harmful, neutral, or beneficial.

\medterm{Mutation} A heritable or acquired alteration in the DNA sequence. Variants may be
single-nucleotide substitutions, insertions, deletions, copy-number changes, or larger
rearrangements and may be benign, pathogenic, or of uncertain significance.

\medterm{Mutation Abnormality} A mutation abnormality is an atypical change in DNA sequence or chromosome structure identified relative to a reference. Its clinical significance must be classified using evidence, such as pathogenic, likely pathogenic, uncertain significance, likely benign, or benign.

\medterm{Mutation Rate} The frequency at which new DNA changes arise in a cell, organism, or generation. The rate differs between genes and species and can be influenced by DNA-copying errors, repair systems, radiation, chemicals, or other factors.

\medterm{Mute} Describing a person who does not speak, whether because of a physical, neurologic, developmental, or psychological reason or by choice; current clinical language should describe the specific communication limitation or preference.

\medterm{Mutism} Absence or refusal of speech despite adequate ability to speak, occurring in conditions such as catatonia, selective mutism, or neurologic disease.

\medterm{Myalgia} Muscle pain or aching. It may affect one area after injury or exercise, or occur throughout the body with infection, inflammation, medicine effects, thyroid disease, or other conditions. Severe pain with weakness or dark urine needs prompt medical assessment because muscle breakdown may be occurring.

\synonyms
Myodynia; Muscle Pain

\medterm{Myalgic Encephalomyelitis} A complex illness characterized by substantial reduction in activity, post-exertional
malaise, and unrefreshing sleep, often with cognitive or autonomic symptoms.

\medterm{Myalgic Encephalomyelitis} A complex illness involving substantial activity-limiting fatigue, post-exertional worsening, unrefreshing sleep, and often cognitive or orthostatic symptoms.

\synonyms
ME/CFS; SEID; Systemic Exertional Intolerance Disease

\medterm{Myalgic Encephalomyelitis/Chronic Fatigue Syndrome} A complex illness characterized by substantial reduction in
activity, post-exertional malaise, unrefreshing sleep, and cognitive or autonomic symptoms. Management is individualized
and focuses on pacing, symptom relief, and treatment of coexisting conditions.

\medterm{Myasthenia} Muscle weakness that often becomes more pronounced with activity and lessens after rest. The term may refer generally to this pattern of weakness or to specific disorders such as myasthenia gravis.

\medterm{Myasthenia Gravis} A chronic autoimmune disorder of the neuromuscular junction, where nerves communicate with skeletal muscles. Antibodies most often target acetylcholine receptors and sometimes proteins such as muscle-specific kinase (MuSK) or LRP4, disrupting signals that cause muscles to contract. Weakness typically worsens with repeated use and improves with rest. It may affect the eyelids and eye movements, facial expression, chewing, speech, swallowing, neck, arms, legs, or breathing muscles. Sensation is usually preserved; thymus abnormalities, including thymoma, may also occur.

\synonyms
MG

\medterm{Myasthenic Crisis} A life-threatening worsening of myasthenia gravis that weakens the muscles needed for breathing or swallowing. Infection, surgery, stress, pregnancy, or certain medicines may trigger it. Severe breathlessness, choking, inability to swallow, or rapidly worsening weakness requires emergency care, breathing support, and specialist treatment.

\synonyms
Myasthenia Gravis Crisis; Acute Myasthenic Exacerbation

\medterm{Myasthenic Syndrome} A disorder of neuromuscular transmission causing fluctuating weakness; it may refer to myasthenia gravis or a congenital or Lambert-Eaton syndrome depending on context.

\medterm{Mycelium} The network of fine branching threads that makes up the growing part of a fungus. It spreads through soil, food, or tissue and can produce the structures that form fungal spores.

\medterm{Mycetoma} A chronic localized infection of skin and subcutaneous tissue caused by fungi or filamentous bacteria, producing swelling, sinuses, and grains.

\medterm{Mycobacteriaceae} A family of slow-growing or rapidly growing, acid-fast bacteria that includes Mycobacterium species causing tuberculosis, leprosy, and nontuberculous infections.

\medterm{Mycobacterial Culture} A laboratory test that tries to grow Mycobacterium bacteria from a sample such as sputum, tissue, or fluid. It can identify tuberculosis or other mycobacteria and may show which medicines will work, but results can take weeks.

\medterm{Mycobacterial Infection} Infection caused by bacteria in the genus Mycobacterium. It includes tuberculosis, leprosy, and nontuberculous mycobacterial disease; symptoms and treatment depend on the species and affected organs.

\medterm{Mycobacterium} A genus of acid-fast bacteria that includes the causes of tuberculosis and leprosy and several nontuberculous infections.

\medterm{Mycobacterium Abscessus} A rapidly growing nontuberculous mycobacterium that can cause pulmonary, skin, and
soft-tissue infections and is often difficult to treat because of multidrug resistance.

\medterm{Mycobacterium Avium} A member of the Mycobacterium avium complex that can cause lung disease, lymph-node infection, or disseminated infection, especially in people with weakened immunity.

\medterm{Mycobacterium Avium Complex} A group of nontuberculous mycobacteria that can cause chronic pulmonary disease
or disseminated infection, particularly in people with impaired immunity. Treatment depends on the species, site,
severity, susceptibility, and patient factors.

\medterm{Mycobacterium Bovis} A tuberculosis-complex bacterium that mainly infects cattle but can cause tuberculosis in humans through unpasteurized dairy or animal exposure.

\medterm{Mycobacterium Chelonae} A rapidly growing nontuberculous mycobacterium that can cause skin, soft-tissue, bone, eye, or disseminated infection.

\medterm{Mycobacterium Fortuitum} A rapidly growing nontuberculous mycobacterium that can cause wound, catheter, bone, lung, or other opportunistic infections.

\medterm{Mycobacterium Gordonae} A usually nonpathogenic water-associated nontuberculous mycobacterium that can occasionally cause disease in people with weakened immunity.

\medterm{Mycobacterium Haemophilum} A nontuberculous mycobacterium that can cause skin, joint, and bone infections, especially in people with impaired immunity.

\medterm{Mycobacterium Infection} Infection caused by a Mycobacterium species, including tuberculosis, leprosy, and nontuberculous mycobacterial disease.

\medterm{Mycobacterium Kansasii} A nontuberculous mycobacterium that commonly causes chronic lung disease resembling tuberculosis and can occasionally cause disseminated infection.

\medterm{Mycobacterium Leprae} The bacterium that causes leprosy, a long-lasting infection affecting the skin and nerves. It usually spreads only after prolonged close contact, develops slowly, and can cause numbness, weakness, or skin lesions without treatment.

\medterm{Mycobacterium Marinum} A nontuberculous mycobacterium associated with fish tanks and aquatic exposure that can cause
chronic skin and soft-tissue infection.

\medterm{Mycobacterium Scrofulaceum} A nontuberculous mycobacterium associated especially with cervical lymph-node infection in children and occasional other disease.

\medterm{Mycobacterium Szulgai} A rare nontuberculous bacterium that can cause long-lasting lung disease or other infections, especially in people with underlying lung problems or weakened immunity. Diagnosis requires specialist laboratory testing and clinical assessment.

\medterm{Mycobacterium Tuberculosis} The bacterium that causes tuberculosis. It spreads through the air and most often affects the lungs, but can involve other organs. Latent infection has no symptoms and does not spread; active disease may cause cough, fever, night sweats, or weight loss.

\medterm{Mycobacterium Ulcerans} A mycobacterium that causes Buruli ulcer, a chronic destructive skin infection associated with the toxin mycolactone.

\medterm{Mycology} The branch of microbiology that studies fungi, including yeasts, moulds, and mushrooms.
Medical mycology focuses on fungal pathogens that cause disease in humans. Fungal
infections (mycoses) are classified as superficial, cutaneous, subcutaneous, or
systemic.

\medterm{Mycophenolate Mofetil} Mycophenolate mofetil is a prodrug of mycophenolic acid that inhibits inosine monophosphate dehydrogenase and lymphocyte proliferation. It is used as an immunosuppressant, especially after organ transplantation, and can cause infection, cytopenias, gastrointestinal effects, and serious fetal harm.

\medterm{Mycophenolic Acid} The active immune-suppressing substance formed from mycophenolate medicines. It blocks an enzyme needed by multiplying lymphocytes and helps prevent organ-transplant rejection, but can cause infections, low blood counts, diarrhea, and fetal harm.

\medterm{Mycoplasma} A genus of very small bacteria lacking a cell wall, including organisms that cause respiratory and genitourinary infections.

\medterm{Mycoplasma Genitalium} A sexually transmitted bacterium that can cause urethritis, cervicitis, pelvic
inflammatory disease, and, less commonly, proctitis. Many infections are asymptomatic.
Diagnosis is usually by a nucleic acid amplification test, and macrolide resistance is
common, so treatment should follow current susceptibility-informed guidance.

\medterm{Mycoplasma Hominis} A wall-less bacterium associated with the genital tract that can cause pelvic, postpartum, urinary, wound, or joint infection.

\medterm{Mycoplasma Pneumonia} Atypical pneumonia most common in ages 5--35. Dry cough, low-grade fever, headache.
Patchy infiltrates on CXR disproportionate to symptoms (``walking pneumonia'').

\medterm{Mycoplasma Pneumoniae} A bacterium without a cell wall that can cause bronchitis, sore throat, or pneumonia, especially in school-age children and young adults. Illness often develops gradually with cough, fever, headache, or tiredness; complications outside the lungs are uncommon but possible.

\medterm{Mycoplasma Pneumoniae-Associated Mucositis} Inflammatory lesions of the oral, ocular, or genital mucosa occurring as an
extrapulmonary manifestation of Mycoplasma pneumoniae infection, often presenting with
target-like skin lesions.

\medterm{Mycoplasmal Pneumonia} Pneumonia caused by Mycoplasma species, most often Mycoplasma pneumoniae, usually producing a gradual cough and fever.

\medterm{Mycoplasmataceae} A family of very small bacteria that lack a cell wall and includes Mycoplasma and related genera.

\medterm{Mycosis} An infection caused by a fungus, such as yeast or mold. Fungal infections may affect the skin, nails, mouth, lungs, or other organs, depending on the type.

\medterm{Mycosis Fungoides} The most common cutaneous T-cell lymphoma, presenting with progressive patches, plaques,
and tumors on the skin with an indolent course. Staging and therapy depend on skin
involvement and extracutaneous spread.

\synonyms
MF; Cutaneous T-Cell Lymphoma

\medterm{Mycosis Fungoides Lymphoma} Mycosis fungoides is a type of cutaneous T-cell lymphoma that usually begins with persistent patches or plaques on the skin.

\medterm{Mycotic Aneurysm} An infected arterial wall or aneurysm, usually caused by bacteria rather than fungi. It may develop from bloodstream infection or an infected embolus and can rupture, causing life-threatening hemorrhage.

\medterm{Mycotoxin} A poisonous substance made by some fungi that can contaminate food, crops, or indoor environments. Exposure may damage the liver, kidneys, nervous system, immune system, or other organs depending on the toxin and dose.

\medterm{Mydriasis} Dilation of the pupil, produced mainly by sympathetic activity through the dilator pupillae muscle or by
reduced parasympathetic activity. Causes include darkness, medicines, toxins, neurologic
disease, and ocular injury.

\medterm{Mydriatic Agents} A medicine that widens the pupil. Some, such as tropicamide or atropine, also temporarily prevent focusing; others, such as phenylephrine, mainly widen the pupil. These medicines are used to examine the inside of the eye and sometimes to ease pain and prevent scarring during eye inflammation. They can trigger a dangerous rise in eye pressure in people with narrow drainage angles.

\medterm{Myelin} Myelin is an insulating sheath around nerve fibers that speeds electrical conduction; it is produced by oligodendrocytes in the CNS and Schwann cells in the PNS.

\medterm{Myelin Sheath} The insulating covering around many nerve fibers that allows electrical signals to travel quickly. It is made by different support cells in the brain, spinal cord, and peripheral nerves and can be damaged in disorders such as multiple sclerosis.

\medterm{Myelination} Myelination is formation of a myelin sheath around an axon by oligodendrocytes in the central nervous system or Schwann cells in the peripheral nervous system. It increases the speed and efficiency of nerve conduction.

\medterm{Myelitis} Myelitis is inflammation of the spinal cord or, less commonly, other neural tissue described by the term. Transverse myelitis can cause weakness, sensory changes, and bladder or bowel dysfunction and requires urgent evaluation for infectious, immune, vascular, or compressive causes.

\medterm{Myeloablative Agonist} This term likely means a myeloablative agent: a treatment that severely or completely destroys blood-forming cells in the bone marrow. It may be used before a stem-cell transplant and causes profound, temporary low blood counts.

\medterm{Myeloablative Chemotherapy} Myeloablative chemotherapy destroys or severely suppresses bone-marrow blood-cell production and is often used before a stem-cell transplant.

\medterm{Myeloblast} An immature precursor of granulocytes in the bone marrow; excess myeloblasts are characteristic of acute myeloid leukemia.

\medterm{Myelocyte} A developing granulocyte precursor in bone marrow that is more mature than a myeloblast and normally absent or rare in peripheral blood.

\medterm{Myelocytic Leukaemia} A blood cancer arising from cells that normally develop into several blood-cell types. It includes acute and chronic myeloid forms, which can cause anemia, infections, bleeding, fatigue, or an enlarged spleen.

\medterm{Myelodysplastic Neoplasms} Alternate name; see \textit{Myelodysplastic Syndromes}.

\medterm{Myelodysplastic Or Myeloproliferative Neoplasms} Myelodysplastic and myeloproliferative neoplasms are clonal bone-marrow disorders causing abnormal blood-cell production. They may produce anemia, infection, bleeding, high counts, enlarged spleen, thrombosis, or progression to acute leukemia; classification and mutation testing guide care.

\medterm{Myelodysplastic Syndrome} A group of bone-marrow cancers in which blood cells develop abnormally and are produced poorly. It can cause anemia, infections, or bleeding and may progress to acute leukemia; diagnosis uses blood and bone-marrow testing.

\medterm{Myelodysplastic Syndromes} Alternate name; see \textit{Myelodysplastic Syndromes}.

\medterm{Myeloencephalitis} Myeloencephalitis is inflammation involving the spinal cord and brain, caused by infection, immune disease, toxins, or other neurologic processes.

\medterm{Myelofibrosis} A bone-marrow cancer in which scar tissue develops in the marrow and disrupts blood-cell production. It may cause anemia, enlarged spleen, fatigue, fever, night sweats, weight loss, or abnormal blood counts.

\synonyms
Idiopathic Myelofibrosis; Primary Myelofibrosis

\medterm{Myelogenous} Relating to blood-forming cells that develop from the bone marrow's myeloid cell line.

\synonyms
Myeloid

\medterm{Myelogram} An imaging study of the spinal canal performed after contrast material is introduced around the spinal cord.

\medterm{Myelography} Contrast injected into subarachnoid space for spinal imaging. Used when MRI
contraindicated. CT myelography more common now.

\medterm{Myeloid} Relating to bone marrow or to the non-lymphoid blood-cell lineages derived from hematopoietic stem cells.

\medterm{Myeloid Cell} A blood-forming cell from the myeloid lineage, which produces red cells, platelets, and several white blood-cell types.

\medterm{Myeloid Leukemia} Myeloid leukemia is a blood cancer arising from abnormal myeloid precursor cells in the bone marrow and blood.

\medterm{Myeloid Sarcoma} Myeloid sarcoma is a mass of immature myeloid cells occurring outside the bone marrow, often associated with acute myeloid leukemia.

\medterm{Myelolipoma} A usually benign tumor composed of mature adipose tissue and blood-forming tissue, most often arising in the adrenal gland and often found incidentally.

\medterm{Myeloma} A cancer of plasma cells, usually called multiple myeloma when it affects many bone-marrow sites. It can cause bone pain, anemia, kidney damage, high blood calcium, infections, or abnormal proteins in blood and urine.

\medterm{Myeloma Cast Nephropathy} Kidney damage caused by abnormal proteins from multiple myeloma blocking and injuring the kidney's filtering tubes. It can cause sudden kidney failure and requires urgent treatment of the myeloma and support of kidney function.

\medterm{Myelomalacia} Softening of the spinal cord, usually due to ischemia, trauma, hemorrhage, or other cord injury.

\medterm{Myelomeningocele} A birth defect in which the spinal cord and its protective coverings bulge through an opening in the spine. It can cause weakness or paralysis, loss of sensation, bladder or bowel problems, and fluid buildup around the brain.

\medterm{Myelomonocyte} A bone-marrow precursor that can develop toward either the monocyte or granulocyte lineages. An immature bone marrow cell that can develop into either monocytes or certain infection-fighting white blood cells.

\medterm{Myelopathy} A general term for any neurological deficit resulting from spinal cord pathology,
encompassing a wide range of causes and clinical presentations.

\medterm{Myeloperoxidase} An enzyme in neutrophil and monocyte granules that uses hydrogen peroxide and chloride to generate antimicrobial oxidants; antibodies against it occur in some autoimmune diseases.

\medterm{Myelopoiesis} The production and development of myeloid blood cells, including granulocytes, monocytes, and some precursor cells, in the bone marrow.

\medterm{Myeloproliferative Disorders} Clonal bone-marrow disorders causing overproduction of one or more blood-cell lineages, including polycythemia vera, essential thrombocythemia, and myelofibrosis.

\medterm{Myeloproliferative Neoplasm} A myeloproliferative neoplasm is a blood cancer in which the bone marrow produces too many mature or immature blood cells.

\medterm{Myelosuppression} Reduced bone marrow production of red cells, white cells, and platelets, causing anemia,
neutropenia, and thrombocytopenia; a common dose-limiting toxicity of chemotherapy and
radiation that increases infection and bleeding risk.

\medterm{Myelosuppressive Therapy} Myelosuppressive therapy reduces bone-marrow production of blood cells, sometimes intentionally during cancer treatment but with risk of anemia, infection, and bleeding.

\medterm{Myenteric Plexus} A network of enteric nerves between the circular and longitudinal muscle layers of the gastrointestinal tract that coordinates intestinal motility.

\medterm{Myhre Syndrome} A rare inherited connective-tissue disorder caused by variants in SMAD4 and characterized by short stature, stiff joints, distinctive facial features, hearing loss, and cardiovascular abnormalities.

\medterm{Myiasis} Infestation of living tissue by larvae (maggots) of various fly species, which feed on
host tissue and can cause pain, secondary infection, and tissue destruction. It may be
cutaneous, involving the skin and subcutaneous tissue, or affect other body sites
including the nose, ears, and wounds.

\medterm{Myoblast} An early muscle cell that develops and joins with others to form skeletal muscle fibers.

\medterm{Myoblastoma} A rare tumor made of cells resembling developing muscle cells. The term has been used for more than one tumor type, so diagnosis requires modern microscopic and molecular classification.
\synonyms
Granular Cell Tumor; Abrikossoff Tumor

\medterm{Myocardial} Relating to the myocardium, the muscular wall of the heart.

\medterm{Myocardial Biopsy} Removal of a small sample of heart muscle, usually by catheter, for microscopic or molecular examination of suspected myocarditis, rejection, infiltrative disease, or other cardiomyopathy.

\medterm{Myocardial Bridging} Myocardial bridging occurs when a segment of a coronary artery tunnels through heart muscle and may be compressed during contraction.

\medterm{Myocardial Contusion} A bruise of the heart muscle caused by blunt chest trauma, which may produce chest pain, arrhythmias, or impaired heart function.

\medterm{Myocardial Infarction} Ischemic necrosis (death) of heart muscle caused by an abrupt reduction or cessation of coronary blood supply, most commonly triggered by the rupture of an atherosclerotic plaque with overlying blood clot formation (thrombus). Symptoms include crushing retrosternal chest pain radiating to the jaw, neck, or left arm, shortness of breath, diaphoresis (sweating), and nausea. Diagnosis is confirmed by cardiac biomarkers (specifically elevated cardiac troponin levels) and electrocardiogram (ECG) changes, categorizing the event as either an ST-elevation myocardial infarction (STEMI, complete coronary occlusion) or non-ST-elevation myocardial infarction (NSTEMI, partial occlusion). Emergent treatment focuses on rapid coronary reperfusion via primary percutaneous coronary intervention (PCI/angioplasty) or clot-dissolving fibrinolytic therapy, along with antiplatelets, anticoagulants, beta-blockers, and statins.

\synonyms
Heart Attack; MI; Acute Myocardial Infarction; STEMI; NSTEMI

\medterm{Myocardial Ischemia} Reduced blood flow and oxygen supply to heart muscle, usually from narrowed coronary arteries, potentially causing chest pain or heart damage.

\synonyms
Cardiac Ischemia; Myocardial Hypoxia

\medterm{Myocardial Perfusion Imaging} A heart scan that compares blood flow to the heart muscle at rest and during exercise or medicine-induced stress. It can show areas with reduced blood supply or old heart-muscle damage.

\medterm{Myocardial Reperfusion} Restoration of blood flow to heart muscle after ischemia, which can limit injury but may also cause reperfusion-related tissue damage.

\medterm{Myocardial Rupture} A tear in the heart wall or other heart structure, usually a severe complication of myocardial infarction.

\medterm{Myocardial Viability} The ability of heart muscle weakened by poor blood flow to recover if blood supply is restored. Imaging tests can show whether the tissue is still alive and may help guide decisions about restoring blood flow.

\medterm{Myocarditis} Inflammation of the heart muscle, often after an infection or because of an immune reaction. It can cause chest pain, breathlessness, abnormal heart rhythms, reduced pumping, or sudden severe illness; severity ranges from mild inflammation to heart failure or shock.

\medterm{Myocarditis Diagnosis} Doctors assess suspected heart-muscle inflammation using symptoms, examination, an electrocardiogram, blood tests, heart imaging, and sometimes MRI or biopsy. No single test proves every case, so results are interpreted together and other causes are excluded.

\medterm{Myocarditis Treatment} Treatment depends on the cause and severity. Care may include rest from strenuous exercise, medicines for heart failure or abnormal rhythms, treatment of infection or autoimmune inflammation when appropriate, and hospital monitoring for severe cases.

\medterm{Myocardium} The thick muscular middle layer of the heart wall that contracts to pump blood throughout the body.

\medterm{Myoclonic Epilepsy} An epilepsy syndrome in which seizures include brief, shock-like involuntary muscle jerks; the specific syndrome determines associated seizure types, age of onset, and prognosis.

\medterm{Myoclonic Seizure} A generalized seizure characterized by brief, shock-like, involuntary single or multiple muscle contractions without loss of consciousness. Commonly affects the upper extremities, resulting in sudden dropping or throwing of objects, and is a hallmark feature of juvenile myoclonic epilepsy.

\medterm{Myoclonus} A sudden, brief, involuntary muscle jerk or group of jerks.

\medterm{Myocyte} A differentiated muscle cell specialized for contraction. Skeletal myocytes are long, multinucleated fibers; cardiac myocytes are branching cells joined by intercalated discs; and smooth-muscle myocytes are spindle-shaped cells without organized sarcomeres. Myocyte injury or death contributes to myopathy, rhabdomyolysis, and myocardial infarction.

\synonyms
Muscle Cell

\medterm{Myodegeneration} Degeneration or breakdown of muscle fibers, caused by inherited, inflammatory, ischemic, toxic, or other disorders.

\medterm{Myoepithelial Cell} A contractile cell around certain glands that helps push secretions from breast or salivary tissue.

\medterm{Myoepithelioma} A myoepithelioma is a tumor composed of myoepithelial cells and can be benign or malignant, commonly arising in salivary or other glands.

\medterm{Myoepithelioma Of Soft Tissue} A rare growth in soft tissue made of myoepithelial cells, which have features of both gland cells and muscle-like support cells. It may be harmless or cancerous; imaging and tissue examination determine treatment.

\medterm{Myofascial Pain Syndrome} A regional pain disorder involving sensitive trigger points in muscle and fascia, with pain that may be referred to other areas and often worsens with pressure or movement.

\medterm{Myofibril} A contractile thread-like structure inside a muscle fiber, composed of repeating sarcomeres containing actin and myosin.

\medterm{Myofibrillar Myopathy} Myofibrillar myopathy is a group of inherited muscle disorders with abnormal protein aggregates and progressive skeletal-muscle weakness, sometimes involving heart or breathing muscles.

\medterm{Myofibrils} Contractile structures inside muscle fibers made of repeating sarcomeres. Long contractile structures inside muscle fibers, made of repeating sarcomeres containing actin and myosin. Their shortening produces muscle contraction.

\medterm{Myofibroblastoma} A benign mesenchymal tumor composed of myofibroblastic cells in collagenous or myxoid stroma, most often arising in the mammary or inguinal region but potentially occurring at other soft-tissue sites.

\medterm{Myofibroblastoma Of The Breast} A rare benign stromal tumor composed of spindle-shaped myofibroblastic cells, usually occurring in older adults. It is typically a well-circumscribed painless mass and is diagnosed by histology with immunohistochemistry.

\medterm{Myofibroblasts} Cells that combine features of fibroblasts and smooth muscle cells and help contract wounds and produce scar tissue.

\medterm{Myofibromatosis} Myofibromatosis is a benign proliferation of myofibroblastic tumors, often in infants or children and involving skin, muscle, bone, or internal organs.

\medterm{Myofilament} A contractile protein filament within a muscle fiber; actin and myosin myofilaments interact through cross-bridges to produce muscle contraction.

\medterm{Myofilaments} Thin actin and thick myosin protein strands inside muscle cells. They slide past one another using energy from ATP, shortening the muscle and producing contraction.
\synonyms
Contractile Filaments; Actin and Myosin Filaments

\medterm{Myogenin} A muscle-specific transcription factor required for skeletal-muscle differentiation and used as an immunohistochemical marker in diagnosing some soft-tissue tumors.

\medterm{Myoglobin} An oxygen-binding protein stored in skeletal and heart muscle. When muscle is injured, myoglobin can enter the blood and urine; large amounts may contribute to kidney damage and are interpreted with other tests.

\medterm{Myoglobin Blood Test} A blood test measuring myoglobin released from injured skeletal or cardiac muscle; it rises early after muscle injury but is nonspecific and is not preferred over troponin for myocardial infarction.

\medterm{Myoglobin Urine Test} A urine test detecting myoglobin released during muscle breakdown, used in evaluating rhabdomyolysis; substantial myoglobinuria can contribute to acute kidney injury.

\medterm{Myoglobinuria} A condition in which muscle protein enters urine, often after severe muscle injury or breakdown.

\medterm{Myokymia} Involuntary, fine, continuous rippling contractions of muscle fibers, often seen around the eye but also associated with nerve hyperexcitability or neurologic disease.

\medterm{Myoma} A noncancerous tumor made of muscle tissue, often referring to a fibroid growing in the uterus.

\medterm{Myomalacia} Softening and breakdown of muscle tissue, sometimes affecting heart muscle after serious injury or poor blood supply.

\medterm{Myomectomy} Surgical removal of uterine fibroids while leaving the uterus in place. It may improve bleeding, pain, or pressure symptoms and can preserve the possibility of pregnancy, although risks and recurrence remain.

\medterm{Myometrium} The thick muscle layer forming most of the uterus. It contracts during menstruation and labor, responds to hormones, and may develop conditions such as fibroids, adenomyosis, or abnormal thickening.

\medterm{Myopathic Facies} A characteristic expressionless, mask-like facial appearance marked by drooping eyelids (ptosis), prominent facial muscle wasting, an inability to bury the eyelashes, and a transverse, tented upper lip (tapir mouth). It occurs due to weakness of the facial musculature in conditions such as facioscapulohumeral dystrophy and myotonic dystrophy.

\medterm{Myopathy} A disease that directly affects muscles and causes weakness, fatigue, cramps, or pain. Causes include inherited conditions, inflammation, medicines, hormone disorders, infections, toxins, and problems with muscle energy production.

\medterm{Myopericytoma} A myopericytoma is a usually benign perivascular tumor composed of cells around small blood vessels, often occurring in skin or soft tissue.

\medterm{Myopia} A refractive error in which light focuses in front of the retina, usually because the eyeball is too long or the eye focuses too strongly. Distant objects appear blurred while near objects are usually clearer. It can be corrected with glasses, contact lenses, or selected refractive procedures.

\synonyms
Short-Sightedness; Nearsightedness

\medterm{Myosin} A protein that uses energy to pull on actin, another muscle protein, and produce movement or force. It is important for contraction of skeletal, heart, and smooth muscle and also helps movement inside many other cells.

\medterm{Myosin Atpase} The activity by which the muscle protein myosin breaks down ATP, the cell’s energy source. This releases energy that lets myosin pull on actin and produce muscle contraction.

\medterm{Myositis} Inflammation of muscle tissue caused by autoimmune disease, infection, medicines, or injury. It may produce weakness, pain, swelling, or fatigue; evaluation often includes examination, blood tests, imaging, and sometimes muscle testing.

\medterm{Myositis Ossificans} A noncancerous condition in which abnormal bone forms inside a muscle or soft tissue, most often following a severe blunt injury, deep bruise, or surgery (commonly in the thigh or upper arm). It causes a firm, painful lump and joint stiffness, and usually matures and improves with rest and gentle physical therapy; aggressive massage or premature re-injury should be avoided.
\synonyms Traumatic myositis ossificans; Heterotopic ossification of muscle; Myositis ossifican.

\medterm{Myotonia} Delayed relaxation of a muscle after voluntary contraction or stimulation, often caused by an inherited or acquired muscle-channel disorder.

\medterm{Myotonia Congenita} An inherited skeletal-muscle channelopathy causing delayed relaxation after voluntary
contraction. Stiffness often improves with repeated movement and may be accompanied by
muscle hypertrophy; onset and severity vary between genetic forms.

\medterm{Myotonia Fluctuans} A rare inherited muscle-channel disorder causing stiffness that varies from day to day or after activity. It is usually mild and does not cause lasting weakness; genetic testing may identify a change in the \textit{SCN4A} gene.

\synonyms
Potassium-Aggravated Myotonia Fluctuans

\medterm{Myotonia Permanens} A rare inherited muscle-channel disorder causing continuous, severe muscle stiffness, including at rest. It may begin in infancy and affect the face, throat, or breathing muscles; it is usually related to a change in the \textit{SCN4A} gene and requires specialist care.

\synonyms
Severe Sodium Channel Myotonia; Continuous Myotonia

\medterm{Myotonic Disorder} A neuromuscular disorder characterized by delayed relaxation after muscle contraction, often with weakness, stiffness, or progressive muscle wasting.

\medterm{Myotonic Dystrophy} An inherited disorder causing delayed muscle relaxation, progressive weakness, and wasting. It can also affect the heart, breathing, swallowing, eyes, hormones, or thinking; symptoms and severity vary between types and people.

\medterm{Myotonic Dystrophy Type 2} An inherited multisystem disorder caused by CCTG repeat expansion in CNBP, producing myotonia, proximal or axial weakness, muscle pain, cataracts, and variable endocrine or cardiac features.

\medterm{Myringitis} Inflammation of the eardrum that may cause ear pain, blisters, drainage, or temporary hearing changes.

\medterm{Myringotomy} A small opening made in the eardrum to drain fluid or relieve pressure in the middle ear. A ventilation tube may be inserted to keep the opening open temporarily. Possible problems include infection, scarring, persistent drainage, or a hole that does not close.

\medterm{Myringotomy With Tympanostomy Tube} Incision of the tympanic membrane with placement of a ventilation tube to equalize
middle-ear pressure and permit drainage. It is used selectively for persistent
middle-ear effusion, recurrent acute otitis media, or associated hearing and
developmental problems.

\medterm{Myristica Oil Poisoning} Poisoning from concentrated nutmeg oil, which can cause nausea, vomiting, agitation, fast heartbeat, dizziness, hallucinations, confusion, seizures, or coma. Suspected ingestion needs prompt poison-control or emergency assessment.

\medterm{Myxedema} Severe hypothyroidism causing accumulation of mucin in tissues, with swelling, slowed body functions, and other symptoms.

\medterm{Myxedema Coma} A life-threatening medical emergency representing the decompensated, end-stage state of severe, long-standing, untreated or undertreated hypothyroidism. Typically precipitated by an acute physiological stressor (such as sepsis, cold exposure, myocardial infarction, stroke, or sedatives), it is clinically defined by a classic triad of profound hypothermia (often core temperature $<35^\circ\text{C}$), altered mental status (ranging from lethargy and confusion to stupor and coma), and multisystem organ slowing. Characteristic complications include severe sinus bradycardia, hypotension, hypoventilation leading to hypercapnic respiratory acidosis, dilutional hyponatremia, and diffuse non-pitting periorbital and peripheral myxedema.
\synonyms{Myxoedema Coma; Decompensated Hypothyroidism; Myxedematous Coma}

\medterm{Myxofibroma} A myxofibroma is a benign tumor containing fibrous and gelatinous myxoid tissue, often arising in the jaw or other connective tissue.

\medterm{Myxoid Cyst} A myxoid cyst is a gelatinous fluid-filled lesion near a joint or tendon, commonly on a finger near the distal joint and nail.

\medterm{Myxoma} A usually noncancerous tumor, most often in the upper left chamber of the heart. It may obstruct blood flow or release clots, causing breathlessness, fainting, fever, or stroke-like symptoms; heart ultrasound is commonly used for diagnosis.

\medterm{Myxopapillary Ependymoma} Myxopapillary ependymoma is a usually slow-growing spinal tumor arising from ependymal cells, commonly in the conus or filum terminale.

\medterm{Myxosarcoma} A rare cancer of soft tissue containing abundant jelly-like material. It can grow into nearby tissue or spread to distant organs; diagnosis requires examination of a tissue sample.

\medterm{MZL} Abbreviation; see \textit{Marginal Zone Lymphoma}.

\medterm{Neonatal Mechanical Ventilation} Invasive or non-invasive mechanical positive-pressure ventilation specifically tailored for neonates and infants with respiratory failure, severe respiratory distress syndrome (RDS), bronchopulmonary dysplasia, or apnea of prematurity. Emphasizes lung-protective strategies such as high-frequency oscillatory ventilation and synchronized intermittent mandatory ventilation.

\synonyms
Pediatric Mechanical Ventilation; Neonatal Ventilatory Support; Infant Mechanical Ventilation

\medterm{Night Sweats in Men} Alternate name; see \textit{Night Sweats}.

\medterm{Non-ST-Elevation Myocardial Infarction} A heart attack in which blood tests show heart-muscle injury but the electrocardiogram does not show persistent ST-segment elevation. It is usually caused by reduced or blocked coronary blood flow and may require urgent assessment and treatment.

\medterm{Open Mitral Valve Surgery} Conventional cardiac surgical repair or mechanical/bioprosthetic replacement of the mitral valve performed via a median sternotomy with cardiopulmonary bypass. Indicated for severe mitral regurgitation or mitral stenosis when percutaneous or minimally invasive approaches are not feasible.

\synonyms
Open Mitral Repair; Open Mitral Valve Replacement; Sternotomy Mitral Surgery

\medterm{Oral Cavity} The entry to the digestive tract, containing the teeth, tongue, and openings of salivary glands. It supports chewing, taste, swallowing, speech, and initial digestion.

\synonyms
Oral Cavity; Buccal Cavity

\medterm{Pediatric Midline Catheters} A peripheral vascular access device (typically 3 to 8 cm in length) inserted into a peripheral vein in the upper extremity of infants and children, with the catheter tip terminating distal to the axillary line. Utilized for intermediate-duration intravenous fluid therapy and medication administration lasting 1 to 4 weeks.

\synonyms
Pediatric Midline Venous Access; Infant Midline IV; Neonatal Midline Catheter

\medterm{Pediatric Myocarditis} Acute or chronic inflammation of the myocardium in infants and children, most commonly triggered by viral infections (enteroviruses, adenovirus, parvovirus B19, SARS-CoV-2) or post-infectious autoimmune responses. Manifests as tachypnea, poor feeding, tachycardia, cardiomegaly, arrhythmias, or acute congestive heart failure.

\synonyms
Childhood Myocarditis; Pediatric Myocardial Inflammation; Acute Pediatric Carditis

\medterm{Psychological Health} Alternate name; see \textit{Mental Health}.

\medterm{Rubeola} Alternate name; see \textit{Measles}.

\medterm{Spontaneous Miscarriage} Pregnancy loss occurring without an intentional medical or surgical procedure.

\medterm{ST-Elevation Myocardial Infarction} A heart attack associated with characteristic ST-segment elevation on an electrocardiogram, usually indicating sudden complete blockage of a coronary artery. It causes acute injury to heart muscle and requires rapid emergency evaluation and restoration of blood flow.

\medterm{Stretch Reflex} The automatic tightening of a muscle when it is suddenly stretched, helping maintain muscle tone and posture. Doctors test this reflex by tapping tendons such as the knee, ankle, biceps, or triceps to assess nerve pathways in the spinal cord and brain.

\medterm{Threatened Miscarriage} Vaginal bleeding occurring in early pregnancy prior to 20 weeks of gestation, in the presence of a closed cervical os and documented fetal cardiac activity. Management includes pelvic rest, ultrasound evaluation of fetal viability, maternal Rh immunoglobulin administration if indicated, and monitoring for progression.

\synonyms
Threatened Abortion; Early Pregnancy Bleeding; Impending Miscarriage

\medterm{Toxoplasma Meningoencephalitis} Alternate name; see \textit{Toxoplasmosis}.

\medterm{Tuberculous Meningitis} Severe granulomatous inflammation of the leptomeninges and basal cisterns caused by \textit{Mycobacterium tuberculosis}. Typically presents with a subacute prodrome of fever, headache, altered mental status, and cranial nerve palsies, complicated by obstructive hydrocephalus and cerebral vasculitis.

\synonyms
TB Meningitis; Tuberculous Meningoencephalitis; CNS Tuberculosis

\medterm{Type 2 Diabetes Medications} Alternate name; see \textit{Oral Antidiabetic Agents}.
